% Supplementary material for the BOBILib paper (separate document).
% Build from bobilib/paper:  make appendix   (latexmk -pdf -cd appendix/appendix.tex)
% Needs: generated/facts.tex and appendix/generated/*.tex (make data), ../refs.bib.
\newcommand{\authoremail}{EMAIL}
% Title of the main paper; keep it identical to the title in main.tex.
\newcommand{\papertitle}{Settling 648 unsolved BOBILib instances: published optima, new
  optimality proofs and infeasibility certificates}

\documentclass[11pt,letterpaper]{article}
\usepackage[letterpaper,margin=1in]{geometry}
\usepackage[T1]{fontenc}
\usepackage{mathptmx}
\renewcommand{\ttdefault}{lmtt}
\usepackage{amsmath,amsthm}
\usepackage{array,booktabs,longtable,multicol}
\usepackage{siunitx}
\usepackage[super]{cite}
\usepackage{titlesec}
\usepackage[labelsep=period,font=small,skip=4pt]{caption}
\usepackage[hidelinks,bookmarksnumbered]{hyperref}

% SPIE-like headings: "A. NOTATION ..." centred, bold, upper case; bold
% subsection headings; run-in bold paragraph headings (LaTeX default).
\titleformat{\section}{\normalfont\bfseries\centering}{\thesection.}{0.5em}{\MakeUppercase}
\titlespacing*{\section}{0pt}{16pt plus 4pt minus 2pt}{8pt plus 2pt}
\titleformat{\subsection}{\normalfont\bfseries}{\thesubsection}{0.6em}{}
\titlespacing*{\subsection}{0pt}{12pt plus 3pt minus 2pt}{5pt plus 1pt}
\renewcommand{\thesection}{\Alph{section}}
\numberwithin{equation}{section}
% Tables of the supplement are S1, S2, ... so that they cannot be confused
% with the tables of the paper.
\renewcommand{\thetable}{S\arabic{table}}

\newtheorem{theorem}{Theorem}[section]
\newtheorem{lemma}[theorem]{Lemma}
\newtheorem{proposition}[theorem]{Proposition}
\newtheorem{corollary}[theorem]{Corollary}
\theoremstyle{definition}
\newtheorem{definition}[theorem]{Definition}
\theoremstyle{remark}
\newtheorem{remark}[theorem]{Remark}

\setcounter{tocdepth}{1}
\makeatletter
% Tighter table of contents (one line per section).
\renewcommand*\l@section[2]{\addpenalty{\@secpenalty}\addvspace{2pt plus 1pt}%
  \@dottedtocline{1}{0pt}{1.6em}{\bfseries #1}{\bfseries #2}}
\makeatother
\setlength{\parskip}{0pt}
\setlength{\emergencystretch}{2em}

% Generated by paper/scripts/gen_facts.py from bobilib/results/; do not edit by hand.
% Each macro is computed from the results files and checked against the paper spec.
% Integers use {,} as thousands separator; percentages are bare numbers (add \%).

\newcommand{\NumInstances}{2{,}654}% BOBILib instances
\newcommand{\NumBobiOptimal}{1{,}027}% BOBILib status optimal
\newcommand{\NumBobiInfeasible}{33}% BOBILib status infeasible
\newcommand{\NumBobiFeasible}{1{,}175}% BOBILib status feasible
\newcommand{\NumBobiOpen}{419}% BOBILib status open
\newcommand{\NumUnsolved}{1{,}594}% feasible or open in BOBILib
\newcommand{\NumLitOptimal}{527}% already optimal in published work
\newcommand{\NumNewOptimal}{81}% proven optimal here
\newcommand{\NumNoFeasible}{40}% no bilevel feasible point (here)
\newcommand{\NumFirstPoint}{230}% first verified point here, BOBILib has none
\newcommand{\NumUnchanged}{716}% status unchanged
\newcommand{\NumSettled}{648}% settled = literature + new optimal + no feasible point
\newcommand{\PctSettled}{41}% settled share of the unsolved, percent
\newcommand{\NumOptimalNow}{608}% optimal now among the unsolved
\newcommand{\NumPacking}{1{,}254}% packing interdiction among the unsolved
\newcommand{\PctPacking}{79}% packing share of the unsolved, percent
\newcommand{\NumKipUnsolved}{392}% KIP not optimal in BOBILib
\newcommand{\NumMkipUnsolved}{862}% MKIP not optimal in BOBILib
\newcommand{\NumCliqueUnsolved}{60}% clique not optimal in BOBILib
\newcommand{\NumNetworkUnsolved}{2}% network not optimal in BOBILib
\newcommand{\NumGeneralizedUnsolved}{12}% generalized not optimal
\newcommand{\NumGeneralPure}{131}% general pure not optimal (incl. infeasible)
\newcommand{\NumGeneralMixed}{168}% general mixed not optimal (incl. infeasible)
\newcommand{\NumGeneralUnsolved}{266}% general feasible or open (not infeasible)
\newcommand{\NumKipInstances}{599}% BOBILib knapsack instances
\newcommand{\NumMkipInstances}{954}% BOBILib multidim-knapsack instances
\newcommand{\NumCliqueInstances}{220}% BOBILib clique instances
\newcommand{\NumNetworkInstances}{72}% BOBILib network instances
\newcommand{\NumGeneralizedInstances}{90}% BOBILib generalized instances
\newcommand{\NumAssignmentInstances}{24}% BOBILib assignment instances
\newcommand{\NumGeneralInstances}{695}% BOBILib general (MIPLIB-derived and other) instances
\newcommand{\NumLitKip}{376}% KIP optima of Weninger and Fukasawa
\newcommand{\NumLitLozano}{96}% optima of Lozano, Bergman and Cire (all families)
\newcommand{\NumLitOr}{87}% OR optima of Lozano et al.
\newcommand{\NumLitGeneralized}{4}% generalized optima of Lozano et al.
\newcommand{\NumLitMiplibLozano}{5}% MIPLIB 3 optima of Lozano et al.
\newcommand{\NumLitImkp}{52}% imkp optima of Fischetti et al. 2019
\newcommand{\NumLitMkip}{139}% MKIP literature optima (OR + imkp)
\newcommand{\NumLitMiplib}{8}% MIPLIB-derived literature optima
\newcommand{\NumLitFischettiSeventeen}{1}% air03-0.100000
\newcommand{\NumLitVasquez}{1}% p0282-0.500000
\newcommand{\NumLitEarlierRelease}{1}% neos-950242_0_100
\newcommand{\NumKipFiles}{583}% KIP files at published optima
\newcommand{\NumKipBobiAbove}{335}% KIP with BOBILib value above the published optimum
\newcommand{\KipBobiExcessMedian}{8}% median excess, percent
\newcommand{\KipBobiExcessMax}{672}% largest excess, percent
\newcommand{\ImkpBobiExcessMax}{23}% largest imkp excess over the optimum, percent
\newcommand{\NumNewKip}{16}% KIP optima proven here
\newcommand{\NumNewClique}{57}% clique optima proven here
\newcommand{\NumNewNetwork}{2}% network optima proven here
\newcommand{\NumNewHpr}{5}% MIPLIB-bound optima
\newcommand{\NumNewOr}{1}% OR closure
\newcommand{\NumRayExact}{38}% rays checked in exact arithmetic
\newcommand{\NumRayFloat}{2}% rays checked in floating point only
\newcommand{\NumOrFirstAnywhere}{88}% OR first points anywhere
\newcommand{\NumOrPointNewToBobilib}{52}% OR points where the literature had one, BOBILib none
\newcommand{\NumOrFirstPoints}{140}% OR points where BOBILib has none
\newcommand{\NumGeneralPoints}{62}% general first points without coupling
\newcommand{\NumCoupledPoints}{28}% general first points with coupling
\newcommand{\NumGeneralFirstPoints}{90}% general first points
\newcommand{\NumMkipUnchanged}{582}% MKIP unchanged
\newcommand{\NumCliqueOpen}{3}% clique still open
\newcommand{\NumGeneralizedUnchanged}{8}% generalized unchanged
\newcommand{\NumGeneralUnchanged}{123}% general unchanged
\newcommand{\NumGeneralNew}{135}% general instances settled or given a first point here
\newcommand{\KipOpenItemsMin}{200}% fewest items among the 16
\newcommand{\KipOpenItemsMax}{500}% most items among the 16
\newcommand{\NumKipSubsetSum}{8}% class 6 (subset sum)
\newcommand{\NumKipEvenOdd}{8}% class 7 (even-odd)
\newcommand{\NumKipOptEqualsBobi}{16}% optimum = BOBILib value
\newcommand{\KipHardestName}{interKP-500-100-6-3}% slowest of the 16 (CPU)
\newcommand{\KipHardestOptimum}{6{,}976}% its optimum
\newcommand{\KipHardestNodes}{10{,}079}% its search nodes
\newcommand{\KipHardestCpu}{10.0}% its CPU seconds (one thread)
\newcommand{\KipHardestIndepNodes}{10{,}116}% nodes of scripts/ss_check.py on it
\newcommand{\KipMostNodesName}{interKP-300-100-6-3}% instance with the most nodes
\newcommand{\KipMostNodes}{19{,}405}% most nodes among the 16
\newcommand{\KipTotalCpu}{32.5}% CPU seconds for all 16
\newcommand{\NumKipEasy}{12}% the 16 other than class 6, budget index 3
\newcommand{\KipEasyMaxCpu}{0.26}% slowest of those 12, CPU s
\newcommand{\KipEasyMaxNodes}{651}% most nodes among those 12
\newcommand{\KipMedianNodes}{334.5}% median nodes of the 16
\newcommand{\KipMedianCpu}{0.155}% median CPU s of the 16
\newcommand{\NumKipPrevSolved}{12}% instances the previous search version proved
\newcommand{\KipPrevMaxSeconds}{427}% slowest proof of the previous version, s
\newcommand{\NumOr}{810}% OR instances
\newcommand{\NumOrPoints}{690}% verified OR points
\newcommand{\NumOrBetter}{469}% points strictly better than every published value
\newcommand{\NumOrNoPublished}{139}% OR instances without a published point
\newcommand{\NumOrBetterLb}{191}% better lower bounds (all variant -100)
\newcommand{\NumOrBetterLbNotOne}{0}% better lower bounds outside -100
\newcommand{\NumOrOneKnapsack}{270}% one follower knapsack (-100)
\newcommand{\NumOrOneKnapsackPoints}{270}% points there
\newcommand{\NumOrSeveral}{540}% several follower knapsacks (-0, -50)
\newcommand{\NumOrSeveralPoints}{420}% points there
\newcommand{\NumOrSeveralHigh}{360}% several knapsacks, budget 50% or 75%
\newcommand{\NumOrSeveralHighPoints}{360}% all certified
\newcommand{\NumOrSeveralQuarter}{180}% several knapsacks, budget 25%
\newcommand{\NumOrSeveralQuarterPoints}{60}% certified at 25%
\newcommand{\NumOrSeveralQuarterNoPoint}{120}% without a point at 25%
\newcommand{\OrGapOneQuarterLit}{24.0}% median published gap, variant -100, budget 0.25
\newcommand{\OrGapOneQuarterNow}{4.1}% median gap now, variant -100, budget 0.25
\newcommand{\OrPairedOneQuarter}{90}% instances in that median
\newcommand{\OrGapOneHalfLit}{34.0}% median published gap, variant -100, budget 0.50
\newcommand{\OrGapOneHalfNow}{7.1}% median gap now, variant -100, budget 0.50
\newcommand{\OrPairedOneHalf}{90}% instances in that median
\newcommand{\OrGapOneThreeQuarterLit}{0.14}% median published gap, variant -100, budget 0.75
\newcommand{\OrGapOneThreeQuarterNow}{0.14}% median gap now, variant -100, budget 0.75
\newcommand{\OrPairedOneThreeQuarter}{90}% instances in that median
\newcommand{\OrGapZeroQuarterLit}{33.3}% median published gap, variant -0, budget 0.25
\newcommand{\OrGapZeroQuarterNow}{31.4}% median gap now, variant -0, budget 0.25
\newcommand{\OrPairedZeroQuarter}{55}% instances in that median
\newcommand{\OrNoPointZeroQuarter}{31}% no published point, -0, 0.25
\newcommand{\OrFirstZeroQuarter}{0}% first points among those
\newcommand{\OrGapZeroHalfLit}{36.6}% median published gap, variant -0, budget 0.50
\newcommand{\OrGapZeroHalfNow}{29.1}% median gap now, variant -0, budget 0.50
\newcommand{\OrPairedZeroHalf}{59}% instances in that median
\newcommand{\OrNoPointZeroHalf}{31}% no published point, -0, 0.50
\newcommand{\OrFirstZeroHalf}{31}% first points among those
\newcommand{\OrGapZeroThreeQuarterLit}{5.1}% median published gap, variant -0, budget 0.75
\newcommand{\OrGapZeroThreeQuarterNow}{0.04}% median gap now, variant -0, budget 0.75
\newcommand{\OrPairedZeroThreeQuarter}{63}% instances in that median
\newcommand{\OrNoPointZeroThreeQuarter}{27}% no published point, -0, 0.75
\newcommand{\OrFirstZeroThreeQuarter}{27}% first points among those
\newcommand{\OrGapFiftyQuarterLit}{29.6}% median published gap, variant -50, budget 0.25
\newcommand{\OrGapFiftyQuarterNow}{28.2}% median gap now, variant -50, budget 0.25
\newcommand{\OrPairedFiftyQuarter}{70}% instances in that median
\newcommand{\OrNoPointFiftyQuarter}{20}% no published point, -50, 0.25
\newcommand{\OrFirstFiftyQuarter}{0}% first points among those
\newcommand{\OrGapFiftyHalfLit}{41.3}% median published gap, variant -50, budget 0.50
\newcommand{\OrGapFiftyHalfNow}{34.8}% median gap now, variant -50, budget 0.50
\newcommand{\OrPairedFiftyHalf}{70}% instances in that median
\newcommand{\OrNoPointFiftyHalf}{20}% no published point, -50, 0.50
\newcommand{\OrFirstFiftyHalf}{20}% first points among those
\newcommand{\OrGapFiftyThreeQuarterLit}{11.8}% median published gap, variant -50, budget 0.75
\newcommand{\OrGapFiftyThreeQuarterNow}{1.4}% median gap now, variant -50, budget 0.75
\newcommand{\OrPairedFiftyThreeQuarter}{80}% instances in that median
\newcommand{\OrNoPointFiftyThreeQuarter}{10}% no published point, -50, 0.75
\newcommand{\OrFirstFiftyThreeQuarter}{10}% first points among those
\newcommand{\OrClosureName}{OR10x100-0.75\_4.dat-0}% the OR instance closed here
\newcommand{\OrClosureValue}{8{,}813}% its optimum
\newcommand{\OrClosureLitLb}{8{,}814}% published (inconsistent) lower bound
\newcommand{\NumOrInconsistent}{4}% OR instances whose published lower bound exceeds the published point
\newcommand{\NumOrLitOptimalPoint}{69}% published OR optima reproduced by a file here
\newcommand{\NumCliqueSet}{60}% instances of the clique set
\newcommand{\NumCliqueSevenThree}{27}% density 0.7 with value 3
\newcommand{\NumCliqueNineFour}{30}% density 0.9 with value 4
\newcommand{\CliqueBobiMin}{5}% smallest BOBILib value
\newcommand{\CliqueBobiMax}{13}% largest BOBILib value
\newcommand{\CliqueMantelLhsForty}{410}% m-k at n=40, d=0.7
\newcommand{\CliqueMantelRhsForty}{400}% floor(n^2/4) at n=40
\newcommand{\CliqueBrouwerLhsForty}{527}% m-k at n=40, d=0.9
\newcommand{\CliqueBrouwerRhsForty}{522}% t_3(n)-floor(n/3)+2 at n=40
\newcommand{\CliqueTuranForty}{533}% t_3(n) at n=40
\newcommand{\CliqueKFortySeven}{136}% budget k at n=40, d=0.7
\newcommand{\CliqueMFortySeven}{546}% edges m at n=40, d=0.7
\newcommand{\CliqueKFortyNine}{175}% budget k at n=40, d=0.9
\newcommand{\CliqueMFortyNine}{702}% edges m at n=40, d=0.9
\newcommand{\CliqueMantelLhsFifty}{643}% m-k at n=50, d=0.7
\newcommand{\CliqueMantelRhsFifty}{625}% floor(n^2/4) at n=50
\newcommand{\CliqueBrouwerLhsFifty}{827}% m-k at n=50, d=0.9
\newcommand{\CliqueBrouwerRhsFifty}{819}% t_3(n)-floor(n/3)+2 at n=50
\newcommand{\CliqueTuranFifty}{833}% t_3(n) at n=50
\newcommand{\CliqueKFiftySeven}{214}% budget k at n=50, d=0.7
\newcommand{\CliqueMFiftySeven}{857}% edges m at n=50, d=0.7
\newcommand{\CliqueKFiftyNine}{275}% budget k at n=50, d=0.9
\newcommand{\CliqueMFiftyNine}{1{,}102}% edges m at n=50, d=0.9
\newcommand{\CliqueMantelLhsSixty}{930}% m-k at n=60, d=0.7
\newcommand{\CliqueMantelRhsSixty}{900}% floor(n^2/4) at n=60
\newcommand{\CliqueBrouwerLhsSixty}{1{,}195}% m-k at n=60, d=0.9
\newcommand{\CliqueBrouwerRhsSixty}{1{,}182}% t_3(n)-floor(n/3)+2 at n=60
\newcommand{\CliqueTuranSixty}{1{,}200}% t_3(n) at n=60
\newcommand{\CliqueKSixtySeven}{309}% budget k at n=60, d=0.7
\newcommand{\CliqueMSixtySeven}{1{,}239}% edges m at n=60, d=0.7
\newcommand{\CliqueKSixtyNine}{398}% budget k at n=60, d=0.9
\newcommand{\CliqueMSixtyNine}{1{,}593}% edges m at n=60, d=0.9
\newcommand{\NumCliquePartCases}{170}% part-size cases decided by exact solvers
\newcommand{\CliqueOpenNames}{normalClique60-70-2, normalClique60-70-3, normalClique60-70-5}% the three open instances
\newcommand{\CliqueOpenBestThreeA}{311}% best 3-partition found (inside edges)
\newcommand{\CliqueOpenBestThreeB}{310}% best 3-partition found (inside edges)
\newcommand{\CliqueOpenBestThreeC}{310}% best 3-partition found (inside edges)
\newcommand{\CliqueOpenInsideFour}{202, 200, 202}% inside edges of the open instances' 4-partitions
\newcommand{\CliqueOpenLpBound}{296}% best LP bound on the K4-transversal number (floor)
\newcommand{\CliqueOpenK}{309}% budget of the open instances
\newcommand{\NetworkOptimum}{26}% optimum of both open trees
\newcommand{\NetworkAttacks}{2{,}138{,}410}% attacks enumerated per open tree
\newcommand{\NetworkOptAttacksA}{4}% optimal attacks, _006
\newcommand{\NetworkOptAttacksB}{28}% optimal attacks, _010
\newcommand{\NetworkNodes}{49}% nodes of the open trees
\newcommand{\NetworkPhi}{5}% attack budget
\newcommand{\NetworkLambda}{5}% protection budget
\newcommand{\NetworkTreesReproduced}{34}% BOBILib-optimal trees reproduced
\newcommand{\NetworkTotalAttacks}{9{,}340{,}760}% attacks over all 36 trees
\newcommand{\NetworkSecsEnumA}{80}% enumeration seconds, _006
\newcommand{\NetworkSecsEnumB}{70}% enumeration seconds, _010
\newcommand{\NetworkRederived}{54}% BOBILib optimal values re-derived
\newcommand{\NetworkBobiOptimal}{70}% BOBILib optimal network instances
\newcommand{\NumGeneralCandidates}{98}% open instances without coupling rows run through bobi general
\newcommand{\NumGeneralFiles}{64}% verified files in results/general
\newcommand{\NumGeneralFilesOptimal}{2}% of them optimal
\newcommand{\NumGeneralFilesFeasible}{62}% of them feasible only
\newcommand{\NumCoupledCandidates}{54}% open instances with coupling rows run through bobi coupled
\newcommand{\NumCoupledAttempted}{52}% attempted
\newcommand{\NumCoupledFiles}{29}% verified files in results/general-coupled
\newcommand{\NumCoupledOptimal}{1}% of them optimal
\newcommand{\NumCoupledTolerance}{9}% coupled files whose follower value uses BOBILib's tolerance
\newcommand{\NumCoupledByHpr}{21}% coupled points at the HPR's leader values
\newcommand{\NumCoupledByPenalty}{5}% coupled points from the penalised HPR
\newcommand{\NumCoupledByFollow}{3}% coupled points from best-response dynamics
\newcommand{\NumHprOptimal}{5}% points meeting the MIPLIB bound
\newcommand{\NumHprOpen}{3}% of them open before
\newcommand{\PigeonOptimum}{\ensuremath{-}9{,}000}% pigeon-10 optimum
\newcommand{\TimtabOptimum}{764{,}771.99999978}% timtab1_0_100 optimum (full precision)
\newcommand{\TimtabOptimumRounded}{764{,}772}% timtab1_0_100 optimum rounded
\newcommand{\NoswotOptimum}{\ensuremath{-}41}% noswot optimum
\newcommand{\NumRayFiles}{41}% direction files (40 open + 1 BOBILib-infeasible)
\newcommand{\NumRayChecks}{43}% exact checks run (ray_verify.jsonl)
\newcommand{\NumVerifyOrOne}{270}% bobi verify passes, OR one knapsack
\newcommand{\NumVerifyOrSeveral}{420}% bobi verify passes, OR several
\newcommand{\NumVerifyNetwork}{36}% bobi verify passes, network
\newcommand{\NumVerifyClique}{60}% bobi verify passes, clique
\newcommand{\NumVerifyNetworkBobi}{72}% BOBILib's inter-fire files passing
\newcommand{\NumSoluOpt}{608}% .solu =opt= lines
\newcommand{\NumSoluInf}{40}% .solu =inf= lines
\newcommand{\NumSoluBest}{695}% .solu =best= lines
\newcommand{\NumSoluBestDual}{194}% .solu =bestdual= lines

% Counts quoted only in this supplement (network attack counts, OR bound sources), written
% by scripts/gen_appendix_tables.py with the tables.
% Generated by paper/scripts/gen_appendix_tables.py from bobilib/results/; do not edit by hand.
% Counts quoted in the appendix text; generated/facts.tex has the paper's.
\newcommand{\SuppNetTrees}{36}% tree instances enumerated (summary.csv)
\newcommand{\SuppNetTreeAttacks}{9{,}340{,}760}% attacks of 0..Phi nodes over the 36 trees
\newcommand{\SuppNetTableTreeAttacks}{5{,}063{,}940}% attacks over the 34 BOBILib-optimal trees
\newcommand{\SuppNetForests}{40}% forests among the inter-fire instances (review.tsv)
\newcommand{\SuppNetRndForests}{4}% rndgraph-25 instances that are forests
\newcommand{\SuppNetRndForestAttacks}{4{,}142}% attacks on those
\newcommand{\SuppNetForestAttacks}{9{,}344{,}902}% attacks on all forests: review DP = TreeDp
\newcommand{\SuppNetFullBudgetAttacks}{1{,}906{,}884}% attacks with exactly Phi nodes, open trees
\newcommand{\SuppNetHighsInstances}{28}% instances with HiGHS at every attack
\newcommand{\SuppNetHighsAttacks}{47{,}496}% attacks solved by HiGHS
\newcommand{\SuppNetHighsChecks}{18{,}072}% HiGHS follower solves in the DP runs (36 trees)
\newcommand{\SuppCoupledExact}{29}% coupled files certified by bobi certify --tight
\newcommand{\SuppCoupledFiles}{29}% coupled files (results/general-coupled)
\newcommand{\SuppGeneralTightenedShown}{26}% points of the table re-derived under the strict tolerance
\newcommand{\SuppOrLbOmega}{166}% one-knapsack lower bounds given by the merged-row online bound
\newcommand{\SuppOrLbMilp}{104}% one-knapsack lower bounds given by the LP bracket only
\newcommand{\SuppOrBetterLbOmega}{166}% improved lower bounds from the online bound
\newcommand{\SuppOrBetterLbMilp}{25}% improved lower bounds from the LP bracket
\newcommand{\SuppOrBracketOptimal}{34}% LP-bracket MILPs solved to optimality (the others: HiGHS's dual bound at the limit)
\newcommand{\SuppOrBracketBetterPoint}{50}% bracket leader choices with a better follower optimum than the file (not written)
\newcommand{\SuppOrNoLb}{143}% OR instances without a nontrivial lower bound (gap n/a)
\newcommand{\SuppOrNoLbPoint}{88}% of them with a verified point here


\hypersetup{pdftitle={Supplementary material for: \papertitle},pdfauthor={Brian J. Connolly}}

\begin{document}

\begin{center}
{\Large\bfseries Supplementary material for:\\[2pt] \papertitle\par}
\vspace{14pt}
{\large Brian J. Connolly\textsuperscript{*}\par}
\vspace{4pt}
Independent Researcher\par
\end{center}
{\renewcommand{\thefootnote}{\fnsymbol{footnote}}%
\footnotetext[1]{E-mail: \authoremail}}
\vspace{6pt}

\noindent This document supplements the paper. It contains the proofs behind
the paper's methods, the verification protocol, a map from each result to its
machine-checked (Lean) proof, a description of the data package, and the
per-instance tables. Section~\ref{app:classes} sorts every result into
three certificate classes (exact certificates, solver-certified numerical
results, benchmark-accepted points). Notation follows the paper. Theorems, lemmas and
equations are numbered by section (Theorem~B.3, equation~(A.1)); tables are
numbered S1, S2, \dots{} The counts quoted in the text are macros
generated from the data package (\texttt{bobilib/results/}), and so are the
tables of Section~\ref{app:tables}. Sections~\ref{app:kip} to~\ref{app:network} say for
each result whether it is formalised in Lean (``Formalised:'') or proved in
prose only (``Not formalised'').

{\small\tableofcontents}

% Notation and conventions used throughout the appendix; the tolerances of
% the checkers. The verification protocol, the certificate classes and the
% threats to validity (verification.tex) follow as the last subsections of
% this section.
\section{Notation and verification protocol}\label{app:notation}

\subsection{Notation and conventions}\label{app:conventions}

\paragraph{Bilevel problems in BOBILib form.}
The leader controls $x\in\mathbf{R}^{n_x}$ and the follower controls
$y\in\mathbf{R}^{n_y}$. As in BOBILib,\cite{thurauf2026bobilib} the leader minimises:
\begin{align}
  \min_{x,y}\quad & c^{\top}x+d^{\top}y \label{eq:bl}\\
  \text{s.t.}\quad & Ax+By\ge a,\qquad x_i\in\mathbf{Z}\ (i\in I_u),\qquad y\in S(x),\notag\\
  S(x)={}&\operatorname*{arg\,min}_{y'}\bigl\{f^{\top}y' : Cx+Dy'\ge b,\ y'_i\in\mathbf{Z}\ (i\in I_l)\bigr\}. \notag
\end{align}
All data are rational. Variable bounds, equality rows and ranged rows of
the instance files are written as pairs of rows of the form ``$\ge$''; when
an argument depends on the side of a row we say so. For a leader choice $x$
the follower's problem is
$P(x)=\min\{f^{\top}y : Cx+Dy\ge b,\ y_i\in\mathbf{Z}\ (i\in I_l)\}$ and its value
is $\varphi(x)\in\mathbf{R}\cup\{-\infty,+\infty\}$ ($+\infty$ when $P(x)$ is
infeasible). So $S(x)\neq\emptyset$ exactly when $\varphi(x)$ is finite and
attained.

An upper-level row is a \emph{coupling row} if it has a nonzero coefficient
on a follower variable. A leader variable is a \emph{linking variable} if its
column in $C$ is nonzero; $x_L$ denotes the linking variables and $C_L$ their
columns, so $Cx=C_Lx_L$, and $P(x)$, $\varphi(x)$ and $S(x)$ depend on $x$ only
through $x_L$.

\begin{definition}[Bilevel feasible point, optimistic convention]\label{def:bilevel-feasible}
A point $(x,y)$ is \emph{bilevel feasible} if $Ax+By\ge a$, $x_i\in\mathbf{Z}$ for
$i\in I_u$, and $y\in S(x)$. The bilevel optimum $\mathrm{OPT}$ is the infimum
of $c^{\top}x+d^{\top}y$ over bilevel feasible points ($+\infty$ if there is
none).
\end{definition}

Problem~\eqref{eq:bl} minimises over the pair $(x,y)$, so when $S(x)$ has
several elements the follower is assumed to pick the one that is best for the
leader. This is the optimistic convention, which BOBILib uses for all of its
reported values. For a fixed $x$ the optimistic choice solves the
\emph{tie-break problem}
\[
  \min\bigl\{c^{\top}x+d^{\top}y : Ax+By\ge a,\ y\ \text{feasible for}\ P(x),\ f^{\top}y\le\varphi(x)\bigr\}.
\]

\paragraph{High point relaxation.}
The \emph{high point relaxation} (HPR) of~\eqref{eq:bl} drops the condition
$y\in S(x)$ and keeps every row, bound and integrality requirement of both
levels and the leader's objective:
\[
  \mathrm{HPR}=\min\bigl\{c^{\top}x+d^{\top}y : Ax+By\ge a,\ Cx+Dy\ge b,\
  x_i\in\mathbf{Z}\ (i\in I_u),\ y_i\in\mathbf{Z}\ (i\in I_l)\bigr\}.
\]

\paragraph{Packing interdiction.}
Items are $N=\{1,\dots,n\}$. The leader chooses $x\in\{0,1\}^n$ subject to
budget rows $Gx\le g$ with $G\in\mathbf{Z}_{\ge 0}^{m_u\times n}$ and
$g\in\mathbf{Z}_{\ge 0}^{m_u}$. The follower then solves
\[
  \Phi(x)=\max\bigl\{p^{\top}y : Wy\le c,\ y\le \mathbf{1}-x,\ y\in\{0,1\}^n\bigr\}
\]
with profits $p\in\mathbf{Z}_{\ge 0}^n$, weights $W\in\mathbf{Z}_{\ge 0}^{m\times n}$
and capacities $c\in\mathbf{Z}_{\ge 0}^m$, and the leader solves
$\mathrm{OPT}=\min\{\Phi(x) : Gx\le g,\ x\in\{0,1\}^n\}$. In BOBILib form the
leader's objective is $p^{\top}y$ and the follower's is $-p^{\top}y$, so the
game is zero-sum: every $y\in S(x)$ gives the leader the same value $\Phi(x)$,
and the optimistic and pessimistic readings coincide. We identify $x$ with the
set $X=\{j : x_j=1\}$ of interdicted items and write $q(Y)=\sum_{j\in Y}q_j$ for
a vector $q$ and a set $Y\subseteq N$. With
$K(T)=\max\{p(Y): Y\subseteq T,\ \sum_{j\in Y}W_{\cdot j}\le c\}$ the best packing
of the items in $T$, we have $\Phi(X)=K(N\setminus X)$.

KIP (the knapsack interdiction problem) is the case $m_u=m=1$. There we write
$\gamma_j=G_{1j}$ for the cost of interdicting item $j$, $\beta=g_1$ for the
leader's budget, $w_j=W_{1j}$ for the follower's weights and $c=c_1$ for its
capacity, so
$\Phi(X)=\max\{p(Y): Y\subseteq N\setminus X,\ w(Y)\le c\}$ and the leader's sets
satisfy $\gamma(X)\le\beta$. MKIP (multidimensional KIP) allows $m_u\ge1$ and
$m\ge1$. In BOBILib's OR set, variant~$-100$ has one follower row ($m=1$) and
several leader rows; variants $-0$ and $-50$ have several follower rows. The
\emph{subset-sum class} of KIP is $p=w$; then $\Phi(X)$ is the largest subset
sum of $\{w_j: j\notin X\}$ that does not exceed $c$.

\paragraph{Clique and network interdiction.}
For edge clique interdiction, $G=(V,E)$ is a simple graph with $n=|V|$ and
$m=|E|$; the leader deletes $D\subseteq E$ with $|D|\le k$, and the follower
takes a maximum clique of $G-D=(V,E\setminus D)$, of size $\omega(G-D)$. In the
attack-protect network game the leader attacks at most $\Phi$ nodes and the
follower then protects at most $\Lambda$ nodes (Section~\ref{app:network}; there
$\Phi$ is the attack budget, as in BOBILib's row name, not a value function).

\paragraph{Letters with two meanings.} Following the paper, a few letters
are reused; the second meaning is local to one section, and the context
always says which is meant: $c$ is the leader's cost vector
in~\eqref{eq:bl} and the follower's capacity in packing interdiction; $D$ is
the follower's matrix in~\eqref{eq:bl} and the deleted edges in
Section~\ref{app:clique}; $G$ is the leader's budget matrix and the graph of
Section~\ref{app:clique}; $m$ is the number of follower rows in
Section~\ref{app:mkip} and the number of edges in Section~\ref{app:clique};
$k$ is an index or a count elsewhere (the depth of a search node in
Section~\ref{app:kip}, a row in Section~\ref{app:mkip}) and the deletion
budget in Section~\ref{app:clique}; $w$ denotes follower weights, the width of a
band table (Section~\ref{app:kip}) and node weights
(Section~\ref{app:network}); $\Phi$ is the follower's value in packing
interdiction and the attack budget in Section~\ref{app:network}. The value of
the online game is $\Omega$ (Sections~\ref{app:kip} and~\ref{app:mkip}), and
$\omega$ is the clique number only; a recession direction of the follower is
$r$ (Section~\ref{app:general}), not to be confused with the cost vector~$d$.

\subsection{Tolerances of the two checkers}\label{app:tolerances}

BOBILib's feasibility checker (described in Section~5.2 of the BOBILib
paper by Th\"urauf, Kleinert, Ljubi\'c, Ralphs and Schmidt,\cite{thurauf2026bobilib}
and distributed with the library as a script) accepts a
solution file in three steps: (i)~the HPR is feasible; (ii)~the file's point
passes the MIPLIB~2017 solution checker\cite{gleixner2021miplib} (version~1.0.3) on the HPR with
tolerance $10^{-4}$ for the linear rows, scaled by the size of the row (a relative test), and
$10^{-4}$ for integrality, and the
stated leader objective matches; (iii)~the follower's problem with the
leader's variables fixed to the file's values is solved to optimality (by
Gurobi with feasibility tolerance $10^{-6}$), giving $\varphi_{\mathrm{check}}$,
and the file's follower value $\varphi_{\mathrm{sol}}$ must satisfy
\begin{equation}\label{eq:bobilib-tol}
  \frac{|\varphi_{\mathrm{sol}}-\varphi_{\mathrm{check}}|}{10^{-10}+|\varphi_{\mathrm{check}}|}\le 10^{-4}.
\end{equation}
(The paper prints the denominator without the absolute value; the script uses
$|\varphi_{\mathrm{check}}|$. The script also contains an optimistic tie-break
test, which its driver does not call.)
Our verifier \texttt{bobi verify} checks (a)~every row, bound and integrality
requirement of the HPR at the file's point with absolute tolerance $10^{-6}$;
(b)~the stated leader objective; (c)~the follower's problem with the leader
fixed, re-solved by an exact knapsack oracle when there is one follower row
and by a HiGHS\cite{huangfu2018parallelizing} MILP otherwise, against test~\eqref{eq:bobilib-tol}, where
differences of at most $10^{-9}$ count as equal. A MILP stopped at a time
limit gives a verdict only when its incumbent is strictly better than the
file's value by more than~\eqref{eq:bobilib-tol} allows (``not optimal'') or
its proven bound passes~\eqref{eq:bobilib-tol} (``optimal''); otherwise the
verdict is ``unproven''. The follower problem contains only the follower's
rows, as in BOBILib's checker; coupling rows constrain the leader. Our
re-certification \texttt{bobi certify} (used for the points with coupling
rows) adds the optimistic tie-break test, and both of its optimality tests
use proven dual bounds and the relative tolerance $10^{-5}$.

The \emph{tight test} \texttt{bobi verify -{}-tight} keeps (a) and (b) but
replaces the relative test~\eqref{eq:bobilib-tol} in (c) by
\begin{equation}\label{eq:tight-tol}
  \varphi_{\mathrm{sol}}-\varphi\le\max\bigl(10^{-6},\,10^{-9}|\varphi|\bigr)
  \quad\text{for every follower optimum $\varphi$ the solver's proof allows,}
\end{equation}
with the follower solved to zero relative gap and absolute gap $10^{-9}$ (for an
integer-valued follower objective the bound is rounded up). A file that passes
BOBILib's test but not~\eqref{eq:tight-tol} uses BOBILib's tolerance on the
follower's value; it is a benchmark-accepted point (Section~\ref{app:classes}).
All \NumTightFiles{} solution files of the data package pass the tight test
(\texttt{results/notes/tight\_verify.md}). \texttt{bobi certify -{}-tight} applies
the same bound with proven dual bounds and also bounds the tie-break: the
leader's objective within $10^{-6}+10^{-9}|F|$ of the proven optimistic
tie-break value.

``Optimal'' in this appendix means: a bilevel feasible point in the sense of
Definition~\ref{def:bilevel-feasible} (checked within the tolerances above)
whose value equals a proven lower bound. How rigorous the point and the bound
are is given by their certificate class (Section~\ref{app:classes}): exact,
solver-certified or benchmark-accepted.

% Verification protocol and threats to validity.
\section{Verification protocol and threats to validity}\label{app:verification}

\subsection{Protocol}

Every claim in the paper rests on one or more of the following checks.

\begin{enumerate}
\item \textbf{Solution files.} Every solution file in the data package is
written in BOBILib's format and passes \texttt{bobi verify}
(Section~\ref{app:notation}): all rows, bounds and integrality requirements of
the HPR within $10^{-6}$; the stated leader objective; and the follower's
problem with the leader fixed, re-solved by the exact knapsack oracle (one
follower row) or by a HiGHS MILP built from the model (otherwise), accepted
by BOBILib's relative test~\eqref{eq:bobilib-tol}. This covers the 583
imported KIP files, the 16 subset-sum files, the 690 OR files (270 with one
follower knapsack, 420 with several), the 64 general files, the 29 files with
coupling rows, the 60 clique files and the 36 network files. The points with
coupling rows also pass \texttt{bobi certify} (proven dual bounds, relative
$10^{-5}$, optimistic tie-break).

\item \textbf{Optimality proofs.} Each ``optimal'' status rests on a lower
bound of one of these kinds: a published proof (cited, with the published
value matched by a verified solution file); the exact subset-sum search
(Section~\ref{app:kip}), whose pruning rules are proved in Lean; a theorem
with a checkable certificate (Mantel, Brouwer and partitions for clique
interdiction, Section~\ref{app:clique}; the HPR bound with an exact
model comparison, Section~\ref{app:general}); or exhaustive enumeration with
checked witnesses (Section~\ref{app:network}). ``No bilevel feasible point''
rests on Theorem~\ref{thm:ray} with an exact rational check (38 instances)
or a floating-point check (2 instances).

\item \textbf{Independent provers for the subset-sum proofs.} Two Python
provers that share no code with the Rust search, one written by the
designer and one by a reviewer who read neither, parse the instance files
themselves and re-prove all 16 lower bounds with only the direct rule and the
per-node rule (Theorem~\ref{thm:node-rule}), in integer arithmetic; both
compare themselves with brute force on random small instances. Two other
designs of the search, built separately, prove
\texttt{interKP-200-100-6-3}, and one of them all four subset-sum instances
with budget index~3.
The reviewer compared every search binary with exhaustive enumeration on
2,990 random small instances (no mismatch).

\item \textbf{Tests of the code against its specification.} Unit tests
compare the subset-sum search with brute force over every deletion, the
threshold tables with the bitset recursion, the node tables with the online
game defined directly and with enumeration of every deletion, and the
table code with a direct transcription of the Lean recursion
$\mathrm{thr}$ (also with costs not ascending, as after restricted copies).
The network DP is compared with brute force over all follower points on
random forests, and the whole enumeration with a brute-force bilevel
optimum.

\item \textbf{Independent reviews.} The clique, network and coupled-row
results were each re-derived by a separate reviewer with its own instance
reader and different solvers or code: two exact maximum-clique searches,
HiGHS and a branch and bound written in C for the part-size cases of
Proposition~\ref{prop:clique-sizes}; a C implementation of the tree DP and a
brute force based on Lemma~\ref{lem:fire-closure}, compared with ours at all
9.3 million attacks; and a second reading of the coupled files with HiGHS's
own MPS reader. The reviewers found only wording errors.

\item \textbf{Exact arithmetic where it is cheap.} The clique inequalities
(Mantel, Brouwer's threshold, the size vectors of
Proposition~\ref{prop:clique-sizes}), the recession directions and the
comparison of each HPR with its MIPLIB model are checked with integers or
exact rationals.

\item \textbf{Machine-checked lemmas.} The Lean~4 project (with Mathlib)
proves the lemmas listed in Section~\ref{app:lean-map}. It contains no
\texttt{sorry}, and \texttt{\#print axioms} shows only Lean's standard axioms
(\texttt{propext}, \texttt{Classical.choice}, \texttt{Quot.sound}) for each
main theorem.
\end{enumerate}

\subsection{Threats to validity}

\begin{enumerate}
\item \textbf{Trusted solvers.} Several verdicts come from MILP or CP solvers:
follower optimality for followers with several rows and for general
instances (HiGHS), the 170 infeasibility verdicts behind $w_3(G)>k$ (each
reproduced by five methods, two of them solver-free), the follower optimum
at the minimising attack in the network proof (HiGHS), and the dual bounds
in the OR bracket (Theorem~\ref{thm:mkip-bracket}). A solver bug that all
methods share would go unnoticed.

\item \textbf{Floating point and tolerances.} Points are feasible within
$10^{-6}$ and follower-optimal within BOBILib's relative $10^{-4}$, not
exactly. Nine coupled points use that slack, and two coupled files are
feasible only within it. The timtab1 point lies $2.2\times10^{-7}$ below its
bound. Two of the 40 ray certificates are checked only in floating point.
The merged-row online bound (Theorem~\ref{thm:mkip-omega}) and the LP
bracket are evaluated in double precision, so their rounding steps are in
the safe direction only up to floating-point error. The window rule of the
subset-sum search uses a floating-point fractional bound; the 16 proofs are
independently confirmed by the integer-only provers above, which do not use
that rule.

\item \textbf{Scope of the Lean proofs.} Lean proves lemmas about abstract
models, not the code. That the Rust loops compute the recursions, the range
shortcuts and early exits of the table code, the bitset and cache code, and
the absence of integer overflow are argued in doc comments and tested, not
formalised. That the instance files have the stated parameters ($n$, $m$, $k$,
budgets) is checked by scripts, not by Lean. Section~\ref{app:lean-map} lists
what is not formalised.

\item \textbf{Results taken from the literature.} Brouwer's theorem was read
in the scanned 1981 report and agrees with later restatements, but it is
neither formalised nor reproved here. Published optima (Weninger and
Fukasawa for KIP, Fischetti et al.\ for imkp, Lozano, Bergman and Cire for
OR, MIPLIB optima for the HPR bounds) and the published lower bounds of the
OR instances with several follower knapsacks are used as reported; for
them only the points are verified here.

\item \textbf{Instance reading.} Every argument depends on reading the MPS
and AUX files correctly. Each family's model is checked row by row against
its expected pattern, and the key results were re-derived with independent
readers (two subset-sum provers, the exact ray and HPR scripts, the
reviewers' parsers).

\item \textbf{Running times.} All times were measured on a shared machine
with four cores that was often oversubscribed; they are indicative only and
do not affect any proof.
\end{enumerate}

% Knapsack interdiction: monotonicity, count formulation, pruning rules,
% online relaxation, robust tables, threshold form, per-node rule, exactness.
\section{Knapsack interdiction proofs}\label{app:kip}

This section proves the rules behind the exact search for subset-sum
knapsack interdiction (KIP with $p=w$, Section~\ref{app:notation}). The search
answers a yes/no question: given an integer $\ell$, does every leader set
within budget leave the follower at least $\ell$? A ``yes'' proves
$\mathrm{OPT}\ge\ell$; a ``no'' comes with a leader set whose follower value is
below~$\ell$. Repeating with $\ell$ equal to the incumbent value proves
optimality (Corollary~\ref{cor:kip-loop}).

\subsection{Monotonicity and the count formulation}

\begin{lemma}[Monotonicity of interdiction]\label{lem:kip-monotone}
Let $\mathcal{F}$ be any family of subsets of $N$ (the follower's feasible
sets) and $\Phi_{\mathcal{F}}(X)=\max\{p(Y): Y\in\mathcal{F},\ Y\cap X=\emptyset\}$
with $p\ge 0$ (a maximum over no set is $0$). If $X\subseteq X'$ then $\Phi_{\mathcal{F}}(X')\le\Phi_{\mathcal{F}}(X)$. If
$\mathcal{F}$ is closed under taking subsets, then
$\Phi_{\mathcal{F}}(X)=\max\{p(Y\setminus X): Y\in\mathcal{F}\}$. The knapsack
follower $\mathcal{F}=\{Y: \sum_{j\in Y}W_{\cdot j}\le c\}$ is closed under taking
subsets because $W\ge 0$; in particular $\Phi(X')\le\Phi(X)$ for $X\subseteq X'$.
\end{lemma}
\begin{proof}
Every $Y\in\mathcal{F}$ that avoids $X'$ also avoids $X$, so the maximum for $X'$
is over a subfamily. For the second claim, a set $Y\in\mathcal{F}$ avoiding $X$
has $p(Y)=p(Y\setminus X)$, which gives ``$\le$''; conversely, for
$Y\in\mathcal{F}$ the set $Y\setminus X$ lies in $\mathcal{F}$ and avoids $X$, which
gives ``$\ge$''. A subset of a packing has no larger weight in any row
because $W\ge 0$.
\end{proof}
{\raggedright\noindent Formalised: \texttt{Bobi/Knapsack.lean}, theorems
\texttt{follower\_antitone}, \texttt{bestAvoiding\_antitone},
\texttt{bestAvoiding\_eq\_sup\_sdiff}, \texttt{packings\_downward\_closed} and
\texttt{K\_eq\_bestAvoiding}.\par}

\begin{lemma}[Cheapest copies first]\label{lem:kip-cheapest}
Consider KIP and a class $Q\subseteq N$ of identical items ($p_i=p_j$ and
$w_i=w_j$ for $i,j\in Q$). For every leader set $X$ there is a leader set $X'$
with $X'\setminus Q=X\setminus Q$, $|X'\cap Q|=|X\cap Q|$,
$\gamma_a\le\gamma_b$ for all $a\in X'\cap Q$ and $b\in Q\setminus X'$,
$\gamma(X')\le\gamma(X)$ and $\Phi(X')=\Phi(X)$.
\end{lemma}
\begin{proof}
First a single swap: if $i\in X\cap Q$, $j\in Q\setminus X$ and
$\gamma_j\le\gamma_i$, let $X^{s}=(X\setminus\{i\})\cup\{j\}$. The transposition of
$i$ and $j$ preserves $p$ and $w$ and maps $N\setminus X$ onto $N\setminus X^{s}$,
so it maps the packings avoiding $X$ bijectively onto those avoiding $X^{s}$
with the same profits; hence $\Phi(X^{s})=\Phi(X)$, and
$\gamma(X^{s})=\gamma(X)-\gamma_i+\gamma_j\le\gamma(X)$. Now induct on the integer
$\gamma(X)$. If some $i\in X\cap Q$ and $j\in Q\setminus X$ have $\gamma_j<\gamma_i$,
the swap lowers the cost strictly, keeps $X\setminus Q$, $|X\cap Q|$ and
$\Phi$, and the induction hypothesis applies to $X^{s}$. Otherwise $X'=X$
works.
\end{proof}
{\raggedright\noindent Formalised: \texttt{Bobi/Knapsack.lean}, theorems
\texttt{K\_map\_equiv}, \texttt{swap\_block} and \texttt{cheapest\_first}.\par}

\begin{definition}[Class instance]\label{def:class-instance}
A \emph{class instance} is a list of classes $1,\dots,L$; class $i$ has a
value $v_i\in\mathbf{Z}_{\ge 1}$ and $n_i$ copies with deletion costs
$\gamma_{i,1},\dots,\gamma_{i,n_i}\in\mathbf{Z}_{\ge 0}$ (in a fixed order, not
necessarily ascending). Let $P_i(r)=\sum_{h\le r}\gamma_{i,h}$ for
$0\le r\le n_i$. A \emph{deletion} is a count vector $r=(r_1,\dots,r_L)$ with
$0\le r_i\le n_i$; its cost is $P(r)=\sum_i P_i(r_i)$ and its survivors form
the multiset $R(r)$ with $n_i-r_i$ copies of $v_i$ for each $i$. For a finite
multiset $M$ of nonnegative integers, $\mathrm{SS}(M)$ is the set of sums of its
sub-multisets ($0\in\mathrm{SS}(M)$). Given integers $0\le lo\le hi$ and a
budget $\beta\ge 0$, the instance is \emph{safe} if every deletion with
$P(r)\le\beta$ has $\mathrm{SS}(R(r))\cap[lo,hi]\neq\emptyset$; a deletion within
budget with $\mathrm{SS}(R(r))\cap[lo,hi]=\emptyset$ is \emph{blocking}.
\end{definition}

\begin{proposition}[Reduction of subset-sum KIP]\label{prop:kip-reduction}
Consider subset-sum KIP with data $\gamma,\beta,w,c$ (all nonnegative
integers). Let $F=\{j:\gamma_j=0\}$, $\mathcal{L}=\{j\notin F: 1\le w_j\le c\}$ and
$g=\gcd\{w_j: j\in\mathcal{L}\}$ ($g=1$ if $\mathcal{L}=\emptyset$). Group
$\mathcal{L}$ into classes by the value $w_j/g$, with each class's costs sorted
ascending, and let $hi=\lfloor c/g\rfloor$. For a deletion $r$ let $X(r)$ be $F$
together with the $r_i$ cheapest items of each class. Then
\begin{enumerate}
\item $\gamma(X(r))=P(r)$ and $\Phi(X(r))=g\cdot\max\bigl(\mathrm{SS}(R(r))\cap[0,hi]\bigr)$;
\item for every leader set $X$ with $\gamma(X)\le\beta$ there is a deletion $r$
with $P(r)\le\gamma(X)$ and $\Phi(X(r))\le\Phi(X)$;
\item for every integer $\ell\ge 0$ with $lo=\lceil\ell/g\rceil\le hi$:
$\mathrm{OPT}\ge\ell$ if and only if the class instance is safe for the band
$[lo,hi]$ and budget $\beta$; and $r$ is blocking if and only if
$P(r)\le\beta$ and $\Phi(X(r))<\ell$.
\end{enumerate}
\end{proposition}
\begin{proof}
(1) Items of $F$ cost nothing. A packing avoiding $X(r)$ contains no item of
$F$; an item outside $F\cup\mathcal{L}$ has $w_j=0$ (then $p_j=0$) or $w_j>c$ (then
it is in no packing), so dropping such items changes neither weight nor
profit. Every weight in $\mathcal{L}$ is $g$ times its class value, so a set
$Y\subseteq\mathcal{L}$ has $w(Y)\le c$ iff its value sum is at most
$\lfloor c/g\rfloor=hi$, and the value sums of packings of $\mathcal{L}\setminus X(r)$
are exactly $\mathrm{SS}(R(r))$.
(2) Let $X_1=X\cup F$: $\gamma(X_1)=\gamma(X)$ and $\Phi(X_1)\le\Phi(X)$ by
Lemma~\ref{lem:kip-monotone}. Items outside $F\cup\mathcal{L}$ do not affect $\Phi$
(as in~(1)) and blocking them costs something, so remove them from $X_1$
(cost does not increase, $\Phi$ unchanged). Apply Lemma~\ref{lem:kip-cheapest}
to each class in turn. The result blocks in each class $i$ some $r_i$ copies,
each no costlier than every unblocked copy, so its cost in class $i$ is the
sum of the $r_i$ smallest costs, $P_i(r_i)$, and its follower value is that
of $X(r)$.
(3) Since $g s\ge\ell$ iff $s\ge lo$ for integers $s$, (1) gives:
$\Phi(X(r))\ge\ell$ iff $\mathrm{SS}(R(r))\cap[lo,hi]\neq\emptyset$. If the instance is
safe, (2) gives $\Phi(X)\ge\Phi(X(r))\ge\ell$ for every $X$ within budget. If a
deletion $r$ is blocking, $X(r)$ is within budget with $\Phi(X(r))<\ell$, so
$\mathrm{OPT}<\ell$.
\end{proof}
\noindent Not formalised as a whole; its two ingredients are
Lemmas~\ref{lem:kip-monotone} and~\ref{lem:kip-cheapest}.

If $\lceil\ell/g\rceil>hi$, the question is trivial: the leader set $F=X(0)$
costs nothing and, by part~(1), $\Phi(F)\le g\cdot hi<\ell$, so
$\mathrm{OPT}<\ell$ and the deletion $r=0$ is a witness. From now on we work with a class instance, a band
$[lo,hi]$ with $lo\le hi$ and a budget~$\beta$. All results below hold for
arbitrary nonnegative costs in the count formulation; where ascending costs
are needed (deletions of arbitrary copies) we say so.

\subsection{Windows of reachable sums and the stall test}

\begin{lemma}[Window extension]\label{lem:window-extend}
Let $M$ be a multiset with $[a,b]\subseteq\mathrm{SS}(M)$, $a\le b$. If a value $v$
satisfies $v\le b-a+1$ then $[a,b+v]\subseteq\mathrm{SS}(M\uplus\{v\})$. More
generally, if values $v_1,\dots,v_q$ satisfy
$v_h\le (b-a+1)+v_1+\dots+v_{h-1}$ for every $h$, then
$[a,\,b+\sum_h v_h]\subseteq\mathrm{SS}(M\uplus\{v_1,\dots,v_q\})$.
\end{lemma}
\begin{proof}
A target $t\in[a,b]$ is a sum of $M$. A target $t\in(b,b+v]$ has
$a\le t-v\le b$, because $t>b$ and $v\le b-a+1$; so $t$ is $v$ plus a sum of
$M$. The
general statement follows by induction on $q$: after adding $v_1,\dots,v_{h-1}$
the window is $[a,b+v_1+\dots+v_{h-1}]$, of width
$(b-a+1)+v_1+\dots+v_{h-1}\ge v_h$.
\end{proof}
{\raggedright\noindent Formalised: \texttt{Bobi/SubsetSum.lean}, theorems
\texttt{window\_extend} and \texttt{window\_extend\_all}.\par}

\begin{lemma}[Survivors that fit]\label{lem:window-survivors}
Let $[a,b]\subseteq\mathrm{SS}(M)$ with $a\le b$ and $\theta=b-a+1$. Let
$(v_1,\eta_1),\dots,(v_q,\eta_q)$ be a list of copies with values $v_h$ and
flags $\eta_h\in\{\text{kept},\text{deleted}\}$. Write $\mu_h=\sum_{h'<h}v_{h'}$
for the mass before copy $h$ and $\delta_h$ for the deleted mass before copy
$h$. If $\delta_h+v_h\le\theta+\mu_h$ for every $h$, then
$[a,\,b+\sigma]\subseteq\mathrm{SS}(M\uplus K)$, where $K$ is the multiset of kept
copies and $\sigma$ its total.
\end{lemma}
\begin{proof}
Let $h$ be a kept copy. The kept mass before it is $\mu_h-\delta_h$, and the
hypothesis says $v_h\le\theta+(\mu_h-\delta_h)$. So the kept copies, in list
order, satisfy the hypothesis of Lemma~\ref{lem:window-extend}.
\end{proof}
{\raggedright\noindent Formalised: \texttt{Bobi/SubsetSum.lean}, theorems
\texttt{fits\_survivors} and \texttt{window\_survivors}
(\texttt{SubsetSum.perm}: subset sums do not depend on the order).\par}

\begin{lemma}[Fractional deletion bound]\label{lem:fractional}
Let copies $h$ have values $v_h>0$ and costs $\gamma_h\ge 0$, sorted so that
$\gamma_1/v_1\le\gamma_2/v_2\le\cdots$. For $d>0$ let
$\Psi(d)=\min\{\sum_h\gamma_h z_h:\sum_h v_hz_h\ge d,\ 0\le z\le 1\}$ ($+\infty$
if $\sum_h v_h<d$). The greedy solution, which fills copies in this order
until mass $d$ is reached, is optimal, and every set $H$ of copies with
$\sum_{h\in H}v_h\ge d$ has $\sum_{h\in H}\gamma_h\ge\Psi(d)$.
\end{lemma}
\begin{proof}
The indicator vector of $H$ is feasible, which proves the last claim. Let
$m_h=v_hz_h$ and $\rho_h=\gamma_h/v_h$, so the cost is $\sum_h\rho_hm_h$ with
$0\le m_h\le v_h$ and $\sum_h m_h\ge d$. Let $m^{*}$ be the greedy allocation,
$\lambda$ the ratio of the copy where it stops, and $m$ any feasible
allocation. Then
$\sum_h\rho_h(m_h-m^{*}_h)=\sum_h(\rho_h-\lambda)(m_h-m^{*}_h)+\lambda\sum_h(m_h-m^{*}_h)$.
Each term of the first sum is nonnegative ($m^{*}_h=v_h\ge m_h$ when
$\rho_h<\lambda$, $m^{*}_h=0\le m_h$ when $\rho_h>\lambda$), and the last term is
nonnegative because $\sum_hm^{*}_h=d\le\sum_hm_h$ and $\lambda\ge 0$.
\end{proof}
\noindent Not formalised (a standard linear programming fact).

The next rule is stated for a search node (Section~\ref{app:kip-exact}): the
decided classes have left a multiset $M$ of survivors, $S=\mathrm{SS}(M)\cap[0,hi]$,
and the remaining (undecided) classes may be deleted with residual budget
$\beta'$. A \emph{completion} is a count vector for the undecided classes of
cost at most~$\beta'$.

\begin{proposition}[Window rule with stall test]\label{prop:window-rule}
Order the undecided classes by ascending value, $u_1,\dots,u_q$ (the proof
works for any order; ascending value is the code's choice), and let
$M_i=\sum_{i'<i}n_{u_{i'}}v_{u_{i'}}$ be the mass of the classes before $u_i$.
Let $\Delta$ be the largest mass that a completion deletes and $\sigma_0$ the
total undecided mass minus $\Delta$. Let $[a,b]\subseteq S$ be a run of
consecutive integers with $b<lo\le b+\sigma_0$ and width $\theta=b-a+1$, and
put $D_i=M_i-v_{u_i}+\theta+1$. Suppose that for every $i$: $D_i>0$, and either
$i=1$ or the fractional bound of Lemma~\ref{lem:fractional} over the copies
of $u_1,\dots,u_{i-1}$ satisfies $\Psi_i(D_i)>\beta'$. Then every completion
leaves survivors $R$ with $lo\in\mathrm{SS}(M\uplus R)$; the node is safe.
\end{proposition}
\begin{proof}
Fix a completion and list the undecided copies class by class in the order
$u_1,\dots,u_q$, flagging the deleted ones. Let $e_i$ be the mass the completion
deletes from $u_1,\dots,u_{i-1}$. Those deleted copies form a set of total cost
at most $\beta'$, so $e_i\ge D_i$ would contradict $\Psi_i(D_i)>\beta'$ by
Lemma~\ref{lem:fractional}; for $i=1$, $e_1=0<D_1$. Hence $e_i\le D_i-1$. At
the copy number $t$ ($t=0,1,\dots$) of class $u_i$, the mass before it is
$M_i+tv_{u_i}$ and the deleted mass before it is at most $e_i+tv_{u_i}$, so
deleted mass plus $v_{u_i}$ is at most
$D_i-1+(t+1)v_{u_i}=\theta+M_i+tv_{u_i}$. Lemma~\ref{lem:window-survivors}
gives $[a,b+\sigma]\subseteq\mathrm{SS}(M\uplus R)$ with $\sigma\ge\sigma_0$ the
surviving mass, and $a\le b<lo\le b+\sigma_0$ puts $lo$ in that window.
\end{proof}
{\raggedright\noindent The per-copy hypothesis is formalised
(\texttt{Bobi/SubsetSum.lean}, theorem \texttt{window\_survivors}); the
reduction to one test per class and the use of Lemma~\ref{lem:fractional} are
argued here and in the doc comment of the code, not formalised. $\Delta$ is
computed exactly by a bounded-knapsack recursion over the undecided classes.\par}

\subsection{The online relaxation}

The online relaxation of Weninger and Fukasawa\cite{weninger2023fast,weninger2025fast}
(Theorem~2 in Section~3 of their preprint) applies to general KIP. Fix a list
$o=(j_1,\dots,j_L)$ of distinct items, any subset of $N$ in any order. The
\emph{online game} on $o$ has $2L$ rounds: at item $j_h$ the leader first
decides whether to block it (allowed if its cost fits the remaining budget);
if the leader does not block it, the follower decides whether to take it
(allowed if its weight fits the remaining capacity). The follower collects
the profits of the items it takes. Both players see all earlier moves. The
value of this finite game of perfect information is given by backward
induction: $\Omega(\emptyset,\beta',\kappa)=0$ and, for $o=(j,o')$,
\begin{align*}
  \Omega(o,\beta',\kappa)&=
  \begin{cases}\min\{\Omega(o',\beta'-\gamma_j,\kappa),\ T\} & \text{if }\gamma_j\le\beta',\\
  T & \text{otherwise,}\end{cases}\\
  T&=\begin{cases}\max\{\Omega(o',\beta',\kappa),\ \Omega(o',\beta',\kappa-w_j)+p_j\} & \text{if }w_j\le\kappa,\\
  \Omega(o',\beta',\kappa) & \text{otherwise.}\end{cases}
\end{align*}

\begin{theorem}[Online lower bound, any play order]\label{thm:online}
For every list $o$ of distinct items and every $X\subseteq N$ with
$\gamma(X)\le\beta$,
$\Omega(o,\beta,c)\le K(o\setminus X)\le\Phi(X)$, where $K(T)$ is the best
packing of the items of $T$ with capacity $c$. Hence $\Omega(o,\beta,c)\le\mathrm{OPT}$.
\end{theorem}
\begin{proof}
In the online game the leader may adapt to the follower's earlier moves; the
real game corresponds to the leader strategy ``block exactly the items of
$X$'', which does not adapt. Formally, we prove by induction on the length of
$o$, for all $\beta',\kappa\ge 0$ and all $X\subseteq o$ with $\gamma(X)\le\beta'$,
that $\Omega(o,\beta',\kappa)\le K_\kappa(o\setminus X)$ (best packing with
capacity $\kappa$). For $o=\emptyset$ both sides are $0$. Let $o=(j,o')$. If
$j\in X$, then $\gamma_j\le\gamma(X)\le\beta'$, so the blocking branch is offered and
$\Omega(o,\beta',\kappa)\le\Omega(o',\beta'-\gamma_j,\kappa)\le K_\kappa(o'\setminus X)=K_\kappa(o\setminus X)$
by induction applied to $X\setminus\{j\}$. If $j\notin X$, then $X\subseteq o'$ and
$\Omega(o,\beta',\kappa)\le T$. By induction $\Omega(o',\beta',\kappa)\le K_\kappa(o'\setminus X)$
and, if $w_j\le\kappa$,
$\Omega(o',\beta',\kappa-w_j)+p_j\le K_{\kappa-w_j}(o'\setminus X)+p_j\le K_\kappa((o'\setminus X)\cup\{j\})$
(add $j$ to an optimal packing). So in both cases of $T$ we get
$T\le K_\kappa((o'\setminus X)\cup\{j\})=K_\kappa(o\setminus X)$. For an arbitrary $X$ within
budget apply this to $X\cap o$ (its cost is no larger and $o\setminus(X\cap o)=o\setminus X$).
Finally $K(o\setminus X)\le K(N\setminus X)=\Phi(X)$ because more items can only help.
\end{proof}
{\raggedright\noindent Formalised: \texttt{Bobi/Online.lean}, theorems
\texttt{online\_le\_K}, \texttt{online\_le\_K\_any} and
\texttt{online\_lower\_bound}.\par}

\begin{remark}
The bound need not be tight. In Weninger and Fukasawa's Table~1 (profits
$(1,1,2)$, leader costs $(1,1,2)$, follower weights $(1,1,3)$, $\beta=2$, $c=3$)
the online game in the play order $(1,2,3)$ is worth $1$ and the bilevel optimum is $2$ (Lean:
\texttt{table1\_online}, \texttt{table1\_opt}, \texttt{table1\_block}, checked by
evaluation).
\end{remark}

\subsection{Robust tables with an end set}

For subset-sum instances the same idea is applied class by class, in the
count formulation. Let $Q=(q_1,\dots,q_L)$ be a list of classes (the
\emph{play order}) and $A\subseteq\mathbf{Z}_{\ge 0}$ an \emph{end set}. The game
$\Gamma_A(Q,\beta',t)$ starts with residual $t$ and budget $\beta'$; at class $q$
the leader announces $r\le n_q$ with $P_q(r)$ at most the remaining budget,
then the follower takes $j\le n_q-r$ copies with $jv_q$ at most the current
residual, which decreases by $jv_q$; the follower wins if the final residual
lies in~$A$.

\begin{definition}[Robust tables]\label{def:tables}
$G_A(\emptyset,\beta')=A$ and, for $Q=(q,Q')$,
\[
  t\in G_A(Q,\beta')\iff \forall r\le n_q \text{ with } P_q(r)\le\beta'\ \
  \exists j\le n_q-r:\ jv_q\le t\ \text{and}\ t-jv_q\in G_A(Q',\beta'-P_q(r)).
\]
So $G_A(Q,\beta')$ is the set of residuals from which the follower wins
$\Gamma_A(Q,\beta',\cdot)$. The robust table is $H=G_{\{0\}}$ and the band table of
width $w\ge 1$ is $G_w=G_{[0,w-1]}$.
\end{definition}

\begin{theorem}[Soundness of the tables]\label{thm:tables-sound}
If $t\in G_A(Q,\beta')$ and $r$ is a deletion of the classes of $Q$ with
$P(r)\le\beta'$, then there are $s\in\mathrm{SS}(R(r))$ and $a\in A$ with $s+a=t$. In
particular $t\in H(Q,\beta')$ gives $t\in\mathrm{SS}(R(r))$, and $t\in G_w(Q,\beta')$ gives
some $s\in\mathrm{SS}(R(r))$ with $s\le t<s+w$. If each class's costs are ascending,
the same holds for a deletion of arbitrary copies of total cost at most
$\beta'$.
\end{theorem}
\begin{proof}
Induction on $L$. For $L=0$ take $s=0$, $a=t$. For $Q=(q,Q')$:
$P_q(r_1)\le P(r)\le\beta'$, so the definition gives $j\le n_q-r_1$ with
$t-jv_q\in G_A(Q',\beta'-P_q(r_1))$, and the rest of $r$ costs at most
$\beta'-P_q(r_1)$. By induction $t-jv_q=s'+a$ with $s'$ a subset sum of the
survivors of $Q'$; then $s=jv_q+s'$ uses $j$ of the $n_q-r_1$ survivors of
$q$. For deletions of arbitrary copies with ascending costs, deleting $r_i$
copies of class $i$ costs at least $P_i(r_i)$, and the survivors are again
$n_i-r_i$ copies of $v_i$.
\end{proof}
{\raggedright\noindent Formalised: \texttt{Bobi/EndTable.lean}, theorems
\texttt{GA\_sound}, \texttt{GA\_sound\_any}, \texttt{G\_iff\_GA},
\texttt{H\_iff\_GA}, \texttt{G\_sound\_of\_GA}, \texttt{H\_sound\_of\_GA};
\texttt{Bobi/RobustTable.lean}, \texttt{H\_sound}, \texttt{H\_sound\_any},
\texttt{take\_sum\_le}, \texttt{counts\_of\_deletion};
\texttt{Bobi/BandTable.lean}, \texttt{G\_sound}, \texttt{G\_sound\_any},
\texttt{G\_one}, \texttt{G\_sound\_one}.\par}

\begin{proposition}[Any play order]\label{prop:play-order}
Let $Q$ be the undecided classes in the search order and $\pi(Q)$ any
permutation of them. If $t\in G_A(\pi(Q),\beta')$, the conclusion of
Theorem~\ref{thm:tables-sound} holds for every deletion $r$ of $Q$ with
$P(r)\le\beta'$.
\end{proposition}
\begin{proof}
Permuting the pairs (class, count) permutes the survivors and leaves the cost
unchanged, and subset sums depend only on the multiset. Apply
Theorem~\ref{thm:tables-sound} to the permuted deletion.
\end{proof}
{\raggedright\noindent Formalised: \texttt{Bobi/EndTable.lean}, theorems
\texttt{counts\_perm}, \texttt{remaining\_eq\_flatten},
\texttt{totalCost\_eq\_sum}, \texttt{remaining\_perm},
\texttt{totalCost\_perm}, \texttt{GA\_sound\_perm} and
\texttt{GA\_sound\_perm\_any}.\par}

\begin{corollary}[Root-table rules]\label{cor:root-rules}
Let $M$ be the decided survivors and $Q$ the undecided classes (any play
order). (i) If $s\in\mathrm{SS}(M)$ and $t\in H(Q,\beta')$, every completion leaves
$s+t\in\mathrm{SS}(M\uplus R)$. (ii) If $[a,b]\subseteq\mathrm{SS}(M)$, $t\in G_w(Q,\beta')$,
$a+t\le\tau$ and $\tau+w\le b+t+1$, every completion leaves
$\tau\in\mathrm{SS}(M\uplus R)$. The search prunes when $s+t$, respectively $\tau$,
lies in $[lo,hi]$.
\end{corollary}
\begin{proof}
(i) Theorem~\ref{thm:tables-sound} gives $t\in\mathrm{SS}(R)$. (ii) It gives
$s'\in\mathrm{SS}(R)$ with $s'\le t<s'+w$; then $a\le\tau-t\le\tau-s'$ and
$\tau-s'\le\tau-t+w-1\le b$, so $\tau-s'\in\mathrm{SS}(M)$.
\end{proof}
{\raggedright\noindent Formalised: \texttt{Bobi/RobustTable.lean}, theorem
\texttt{covers\_sound}; \texttt{Bobi/BandTable.lean},
\texttt{band\_covers\_sound}, \texttt{covers\_sound\_of\_band};
\texttt{Bobi/Threshold.lean}, \texttt{covers\_threshold\_sound},
\texttt{covers\_threshold\_sound\_perm}, \texttt{band\_covers\_threshold\_sound},
\texttt{band\_covers\_threshold\_sound\_perm} (the same on stored thresholds).\par}

\begin{lemma}[Antitone in the budget]\label{lem:antitone}
If $\beta''\le\beta'$ then $G_A(Q,\beta')\subseteq G_A(Q,\beta'')$. If
$A\subseteq A'$ then $G_A(Q,\beta')\subseteq G_{A'}(Q,\beta')$.
\end{lemma}
\begin{proof}
Induction on $L$; the base case does not depend on the budget. A deletion
$r$ with $P_q(r)\le\beta''$ also has $P_q(r)\le\beta'$; the answer $j$ for
budget $\beta'$ leads to $G_A(Q',\beta'-P_q(r))\subseteq G_A(Q',\beta''-P_q(r))$ by
induction. The second claim is the same induction with a larger base set.
\end{proof}
{\raggedright\noindent Formalised: \texttt{Bobi/EndTable.lean}, theorems
\texttt{GA\_antitone}, \texttt{G\_antitone}, \texttt{H\_antitone} and
\texttt{GA\_mono}.\par}

\subsection{Threshold form}

By Lemma~\ref{lem:antitone}, for each residual $t$ the budgets at which
$t\in G_A(Q,\cdot)$ form an initial segment, so one number per $t$ describes
the table. Fix a cap $\mathrm{cap}\ge 0$ (the root budget $\beta$).

\begin{definition}[Threshold]\label{def:thr}
$\mathrm{thr}(\emptyset,t)=\mathrm{cap}$ if $t\in A$ and $-1$ otherwise; for
$Q=(q,Q')$ with value $v$ and $n$ copies,
\[
  \mathrm{thr}(Q,t)=\min\Bigl(\mathrm{cap},\ \min_{\substack{r\le n\\ P_q(r)\le\mathrm{cap}}}
  \Bigl[P_q(r)+\max\Bigl(-1,\ \max_{\substack{j\le n-r\\ jv\le t}}\mathrm{thr}(Q',t-jv)\Bigr)\Bigr]\Bigr).
\]
\end{definition}

\begin{theorem}[Threshold form]\label{thm:thr}
For every integer $B\ge 0$: $B\le\mathrm{thr}(Q,t)$ if and only if $B\le\mathrm{cap}$
and $t\in G_A(Q,B)$. Consequently $-1\le\mathrm{thr}(Q,t)\le\mathrm{cap}$;
$\mathrm{thr}(Q,t)=-1$ iff $t\notin G_A(Q,0)$; and otherwise $\mathrm{thr}(Q,t)$ is the
largest $B\le\mathrm{cap}$ with $t\in G_A(Q,B)$.
\end{theorem}
\begin{proof}
Induction on $L$. For $L=0$: if $t\in A$ the claim reads $B\le\mathrm{cap}$; if
$t\notin A$ both sides are false since $B\ge 0>-1$. For $Q=(q,Q')$ write
$m_r(t)$ for the inner maximum (at least $-1$). Then $B\le\mathrm{thr}(Q,t)$ iff
$B\le\mathrm{cap}$ and $B\le P_q(r)+m_r(t)$ for all $r$ with $P_q(r)\le\mathrm{cap}$.
Assume $B\le\mathrm{cap}$. For $r$ with $P_q(r)>B$ the inequality holds
automatically, as $B\le P_q(r)-1$. For $r$ with $P_q(r)\le B$ (these have
$P_q(r)\le\mathrm{cap}$): $B\le P_q(r)+m_r(t)$ iff $0\le B-P_q(r)\le m_r(t)$ iff some
$j\le n-r$ with $jv\le t$ has $B-P_q(r)\le\mathrm{thr}(Q',t-jv)$ iff, by induction
(note $B-P_q(r)\le\mathrm{cap}$), $t-jv\in G_A(Q',B-P_q(r))$. This is the
definition of $t\in G_A(Q,B)$. The consequences follow by taking $B=0$ and
$B=\mathrm{thr}(Q,t)$.
\end{proof}
{\raggedright\noindent Formalised: \texttt{Bobi/Threshold.lean}, theorems
\texttt{thr\_iff}, \texttt{GA\_iff\_thr}, \texttt{thr\_bounds},
\texttt{thr\_eq\_neg\_one\_iff}, \texttt{thr\_greatest}, \texttt{G\_iff\_thr}
and \texttt{H\_iff\_thr} (definition \texttt{thr}).\par}

\begin{remark}[What the code computes]
The recursion at residual $t$ reads only residuals $\le t$, so arrays indexed
$0,\dots,hi$ hold the exact thresholds of those residuals. With nonnegative
costs $P_q$ is nondecreasing, so stopping at the first $r$ with
$P_q(r)>\mathrm{cap}$ implements the filter. That the code's loops compute
$\mathrm{thr}$ (including the shortcuts that compute only residuals in a range
that can matter, and an early exit once every computed threshold is below the
budget) is argued in the code's doc comments and tested against a direct
transcription of Definition~\ref{def:thr} and against brute force; it is not
formalised.
\end{remark}

\subsection{The per-node rule}

At a search node the decided survivors $M$ reach the sums $S=\mathrm{SS}(M)\cap[0,hi]$.
The follower wins with any total in $[lo,hi]$, so it need not fix its target
in advance. The node's end set is
\[
  A_{\mathrm{node}}=\{a\in[0,hi] : \exists s\in S,\ s\le a\le s+(hi-lo)\},
\]
that is, $(S+[0,hi-lo])\cap[0,hi]$, and the node table asks whether
$hi\in G_{A_{\mathrm{node}}}(Q,\beta')$.

\begin{theorem}[Per-node rule]\label{thm:node-rule}
Let $S\subseteq\mathrm{SS}(M)$, $lo\le hi$, $A$ any set whose elements $a$ each have
some $s\in S$ with $s\le a\le s+(hi-lo)$, and $\pi(Q)$ any permutation of the
undecided classes $Q$. If $hi\in G_A(\pi(Q),\beta')$, then every completion
(a deletion $r$ of $Q$ with $P(r)\le\beta'$, costs in any order) leaves some
$u\in\mathrm{SS}(M\uplus R(r))$ with $lo\le u\le hi$. In threshold form: if
$\beta'\le\mathrm{cap}$ and $\beta'\le\mathrm{thr}(\pi(Q),hi)$ computed with end set
$A_{\mathrm{node}}$, the node is safe at every budget up to $\mathrm{thr}(\pi(Q),hi)$.
\end{theorem}
\begin{proof}
Proposition~\ref{prop:play-order} gives $x\in\mathrm{SS}(R(r))$ and $a\in A$ with
$x+a=hi$. Pick $s\in S$ with $s\le a\le s+(hi-lo)$. Then $u=s+x$ is a subset sum
of $M\uplus R(r)$ and $u=s+hi-a\in[hi-(hi-lo),hi]=[lo,hi]$. The threshold form
follows from Theorem~\ref{thm:thr} with $B=\beta'$ (and for smaller budgets,
since $B\le\mathrm{thr}$ is the same condition). The hypothesis $lo\le hi$ is
needed: with $lo>hi$ there is no target, and the code returns $-1$.
\end{proof}
{\raggedright\noindent Formalised: \texttt{Bobi/EndTable.lean}, theorems
\texttt{node\_rule}, \texttt{node\_table\_sound}, \texttt{node\_table\_sound\_perm}
and \texttt{node\_table\_sound\_any} (definition \texttt{nodeEnd});
\texttt{Bobi/Threshold.lean}, \texttt{node\_table\_threshold\_sound} and
\texttt{node\_table\_threshold\_sound\_counts} (the form the search uses:
count formulation, any play order, any cost order).\par}

\begin{remark}[Strength]
The per-node rule contains the rule of $H$ in the same play order: if $s\in S$,
$s\le\tau$, $\tau\in[lo,hi]$ and $\tau-s\in H(Q,\beta')$, then
$hi\in G_{A_{\mathrm{node}}}(Q,\beta')$ (shift the follower's strategy up by
$hi-\tau+s$). Formalised: \texttt{Bobi/EndTable.lean}, theorem
\texttt{node\_table\_of\_H}. The analogous containment of the band rules is
argued in the code's doc comment, not formalised; soundness does not need
either. The play order matters only for strength; the search plays the
classes in ascending order of deletion cost per unit of mass.
\end{remark}

\subsection{Restricted copies}

\begin{proposition}[Restricted follower]\label{prop:restrict}
For $k\ge 1$ let $I^{(k)}$ be the class instance in which every class $i$ with
$n_i>k$ is replaced by a class with the same value, $k$ copies and costs
$(\gamma_{i,1}+\dots+\gamma_{i,n_i-k+1},\ \gamma_{i,n_i-k+2},\dots,\gamma_{i,n_i})$. If
$I^{(k)}$ is safe (same band and budget), then so is the original instance.
\end{proposition}
\begin{proof}
Map a deletion $r$ of the original to $r'$ with
$r'_i=\max(0,r_i-(n_i-k))$ for the restricted classes and $r'_i=r_i$ otherwise.
If $r'_i=0$ its cost is $0\le P_i(r_i)$; otherwise it is
$\gamma_{i,1}+\dots+\gamma_{i,n_i-k+r'_i}=P_i(r_i)$. So $P'(r')\le P(r)$. The
restricted survivors number $k-r'_i=\min(k,n_i-r_i)\le n_i-r_i$, so
$R'(r')\subseteq R(r)$ as multisets and $\mathrm{SS}(R'(r'))\subseteq\mathrm{SS}(R(r))$. Safety
of $I^{(k)}$ applied to $r'$ gives a sum in $[lo,hi]$ for $r$.
\end{proof}
\noindent Not formalised. The restricted costs are not ascending, which is
why the count-formulation results above make no assumption on the order of
costs. (The option is off by default; no run recorded in the repository
uses it.)

\subsection{Why the search is exact}\label{app:kip-exact}

The search fixes an order of the classes (any order). A node at depth $k$ has
decided counts $r_1,\dots,r_k$ with residual budget $\beta'=\beta-\sum_{i\le k}P_i(r_i)\ge 0$,
decided survivors $M_k$ and $S_k=\mathrm{SS}(M_k)\cap[0,hi]$. The node is \emph{safe}
if every completion $(r_{k+1},\dots,r_L)$ with $\sum_{i>k}P_i(r_i)\le\beta'$ has
$\mathrm{SS}(M_k\uplus R_{>k})\cap[lo,hi]\neq\emptyset$. The root ($k=0$, $S_0=\{0\}$) is safe
iff the class instance is safe.

\begin{lemma}[Node state]\label{lem:node-state}
Whether a node is safe depends only on $(k,S_k,\beta')$, and a node safe at
$\beta'$ is safe at every $\beta''\le\beta'$.
\end{lemma}
\begin{proof}
$\mathrm{SS}(M\uplus R)=\mathrm{SS}(M)+\mathrm{SS}(R)$, and all values are nonnegative, so a sum in
$[lo,hi]$ has its $\mathrm{SS}(M)$ part in $[0,hi]$. Hence the condition reads
$(S_k+\mathrm{SS}(R_{>k}))\cap[lo,hi]\neq\emptyset$, and the completions depend only on $k$
and $\beta'$, with fewer completions for smaller budgets.
\end{proof}

\begin{theorem}[Sound pruning]\label{thm:kip-rules}
Each of the following implies that the node $(k,S_k,\beta')$ is safe.
\begin{enumerate}
\item Direct: $S_k\cap[lo,hi]\neq\emptyset$.
\item Guaranteed copies: with $d_i=\max\{r\le n_i: P_i(r)\le\beta'\}$, the
multiset $U$ with $n_i-d_i$ copies of $v_i$ for each undecided $i$ has
$(S_k+\mathrm{SS}(U))\cap[lo,hi]\neq\emptyset$.
\item Window: the hypothesis of Proposition~\ref{prop:window-rule} holds for a
run of $S_k$.
\item Root tables: Corollary~\ref{cor:root-rules} gives a target in $[lo,hi]$,
for a table built once at the root (any play order).
\item Node table: $\beta'\le\mathrm{thr}(\pi(Q),hi)$ with end set $A_{\mathrm{node}}$
(Theorem~\ref{thm:node-rule}).
\item Cache: an entry $(k,S_k)\mapsto\beta^{*}$ with $\beta'\le\beta^{*}$, where an
entry is stored only after $(k,S_k)$ was shown safe at $\beta^{*}$ by one of
the rules or by exhausting its children; rule~5 stores
$\beta^{*}=\mathrm{thr}(\pi(Q),hi)$.
\end{enumerate}
\end{theorem}
\begin{proof}
(1) The decided survivors are present in every completion. (2) Costs are
nonnegative, so a completion within $\beta'$ has $P_i(r_i)\le\beta'$ and, $P_i$
being nondecreasing, $r_i\le d_i$; at least $n_i-d_i$ copies of $v_i$ survive,
and a subset sum of a sub-multiset is a subset sum of the multiset. (3)--(5)
are Proposition~\ref{prop:window-rule}, Corollary~\ref{cor:root-rules} and
Theorem~\ref{thm:node-rule}. (6) is Lemma~\ref{lem:node-state} together with
Theorem~\ref{thm:node-rule}, which holds at every budget up to the stored
threshold.
\end{proof}
\noindent Not formalised as a whole: rules~(3)--(5) are the results cited in
the proof (rule~(3) only in part, see Proposition~\ref{prop:window-rule}), and
rules~(1), (2) and~(6) and Lemma~\ref{lem:node-state} are argued here.

\begin{theorem}[Exactness]\label{thm:kip-exact}
The search visits the root; at a node it returns ``safe'' if a rule of
Theorem~\ref{thm:kip-rules} applies or if all children
$r_{k+1}\in\{r\le n_{k+1}:P_{k+1}(r)\le\beta'\}$ return ``safe''; at a leaf
($k=L$) to which no rule applies it returns ``blocking'' with the counts on
the path. If the search ends without reaching its node limit, then ``safe''
at the root means the class instance is safe, and ``blocking'' returns a
blocking deletion.
\end{theorem}
\begin{proof}
By induction from the leaves, every node that returns ``safe'' is safe: for
the rules this is Theorem~\ref{thm:kip-rules}; for a node whose children are
all safe, every completion within $\beta'$ starts with some $r_{k+1}$ with
$P_{k+1}(r_{k+1})\le\beta'$ (costs are nonnegative), and the rest of it is a
completion of that child within $\beta'-P_{k+1}(r_{k+1})$. A leaf returning
``blocking'' has $S_L\cap[lo,hi]=\emptyset$ (rule~1 did not apply), and $M_L=R(r)$
for the path's deletion $r$, whose cost is $\beta-\beta'\le\beta$; so $r$ is
blocking. A node returns ``blocking'' only when a child does.
\end{proof}
\noindent Not formalised (the search code; see Section~\ref{app:lean-map}).
The rules it relies on are formalised as listed above, except the fractional
bound of Lemma~\ref{lem:fractional}.

\begin{corollary}[Exact solve]\label{cor:kip-loop}
Start from any leader set $X_0$ within budget with $u_0=\Phi(X_0)$ (the code
uses Weninger and Fukasawa's greedy leader). For $h=0,1,\dots$ run the search
with $\ell=u_h$: if it returns ``safe'', $\mathrm{OPT}=u_h$; if it returns a
blocking deletion $r$, set $X_{h+1}=X(r)$ and $u_{h+1}=\Phi(X(r))<u_h$. If no
search hits its node limit, this ends after at most $u_0+1$ searches with a
proof of optimality and an optimal leader set.
\end{corollary}
\begin{proof}
By Proposition~\ref{prop:kip-reduction}(3) and Theorem~\ref{thm:kip-exact},
``safe'' proves $\mathrm{OPT}\ge u_h=\Phi(X_h)\ge\mathrm{OPT}$, and ``blocking''
yields $X(r)$ within budget with $\Phi(X(r))<u_h$. (If
$\lceil u_h/g\rceil>hi$, no search is needed: $r=0$ is blocking, as noted
after Proposition~\ref{prop:kip-reduction}.) The values $u_h$ are
nonnegative integers that strictly decrease.
\end{proof}

For each of the \NumNewKip{} instances that were open in the literature the loop ended
with ``safe'' at BOBILib's own value, so BOBILib's point was already optimal
and the proof of optimality is new. These proofs are computational proofs
supported by formally verified lemmas, not formally verified solutions: Lean
proves the pruning rules (Section~\ref{app:lean-map}), and the search that
applies them, like the two integer-only provers that repeat it
(Section~\ref{app:verification}), is an exhaustive computation in integer
arithmetic whose code is tested but not verified. In the terms of
Section~\ref{app:classes} they are exact certificates (class~(i)).

% Multidimensional knapsack interdiction: the dualised follower LP, the
% integrality-gap bracket, the merged-row online bound, certified points.
\section{Multidimensional knapsack interdiction proofs}\label{app:mkip}

We use the packing interdiction notation of Section~\ref{app:notation}: leader
rows $Gx\le g$, follower rows $Wy\le c$ with $m$ rows, $p,G,g,W,c\ge 0$ integral,
and $\Phi(x)$ the follower's optimum. Let
\[
  U(x)=\max\bigl\{p^{\top}y : Wy\le c,\ 0\le y\le\mathbf{1}-x\bigr\}
\]
be the value of the follower's linear programming (LP) relaxation, and
$p_{\max}=\max_jp_j$.

\subsection{The dualised follower LP}

\begin{lemma}[Dual form of $U$]\label{lem:mkip-dual}
For $x\in\{0,1\}^n$,
\[
  U(x)=\min_{\lambda\in\mathbf{R}^m_{\ge 0}}D(x,\lambda),\qquad
  D(x,\lambda)=c^{\top}\lambda+\sum_{j=1}^n\max\bigl(0,\,p_j-(W^{\top}\lambda)_j\bigr)(1-x_j),
\]
and the minimum is attained by some $\lambda$ with $0\le\lambda_k\le p_{\max}/w^{\min}_k$
for every row $k$, where $w^{\min}_k$ is the smallest positive entry of row $k$
of $W$ ($w^{\min}_k=1$ if the row is zero).
\end{lemma}
\begin{proof}
The LP is feasible ($y=0$, since $c\ge 0$) and bounded ($0\le y\le\mathbf{1}$), so by
LP duality $U(x)=\min\{c^{\top}\lambda+\sum_j(1-x_j)\nu_j : W^{\top}\lambda+\nu\ge p,\ \lambda,\nu\ge 0\}$,
and the minimum is attained. For fixed $\lambda$ the best feasible $\nu_j$ with
$x_j=0$ is $\max(0,p_j-(W^{\top}\lambda)_j)$; for $x_j=1$ the term has coefficient
$0$. This gives $U(x)=\min_\lambda D(x,\lambda)$. For the box: if
$\lambda_k>p_{\max}/w^{\min}_k$, lower $\lambda_k$ to $p_{\max}/w^{\min}_k$. Every
item $j$ with $W_{kj}>0$ still has $(W^{\top}\lambda)_j\ge\lambda_kW_{kj}\ge p_{\max}\ge p_j$,
so its term stays $0$; items with $W_{kj}=0$ are unaffected; and $c_k\lambda_k$
does not increase because $c_k\ge 0$. So $D$ does not increase.
\end{proof}
\noindent Not formalised.

\begin{proposition}[Exact linearisation]\label{prop:mkip-milp}
Let $U^{*}=\min\{U(x): Gx\le g,\ x\in\{0,1\}^n\}$. Then $U^{*}$ is the optimal
value of the MILP
\begin{align*}
  \min\quad & c^{\top}\lambda+\sum_{j:\,p_j>0}v_j\\
  \text{s.t.}\quad & v_j\ge p_j-(W^{\top}\lambda)_j-p_jx_j,\quad v_j\ge 0\quad (p_j>0),\\
  & 0\le\lambda_k\le p_{\max}/w^{\min}_k\quad (k=1,\dots,m),\qquad Gx\le g,\quad x\in\{0,1\}^n.
\end{align*}
\end{proposition}
\begin{proof}
For $x_j\in\{0,1\}$ and $(W^{\top}\lambda)_j\ge 0$,
\[
  \max\bigl(0,p_j-(W^{\top}\lambda)_j\bigr)(1-x_j)=\max\bigl(0,\,p_j-(W^{\top}\lambda)_j-p_jx_j\bigr):
\]
both sides equal $\max(0,p_j-(W^{\top}\lambda)_j)$ when $x_j=0$, and both are $0$
when $x_j=1$.
For fixed $(x,\lambda)$ the smallest feasible $v_j$ is this maximum, and items
with $p_j=0$ contribute $0$. So the MILP's value is the minimum of
$D(x,\lambda)$ over leader-feasible $x$ and $\lambda$ in the box, which is
$\min_xU(x)=U^{*}$ by Lemma~\ref{lem:mkip-dual}.
\end{proof}
\noindent Not formalised.

\subsection{The integrality-gap bracket}

\begin{lemma}[Integrality gap of a packing problem with $m$ rows]\label{lem:mkip-gap}
For every $x\in\{0,1\}^n$, with $p_{\max}(x)=\max\{p_j: x_j=0\}$ ($0$ if there is
no such $j$),
\[
  U(x)-m\,p_{\max}(x)\ \le\ \Phi(x)\ \le\ U(x).
\]
\end{lemma}
\begin{proof}
$\Phi(x)\le U(x)$ because the LP is a relaxation. The LP's feasible region
$P=\{y: Wy\le c,\ 0\le y_j\le 1-x_j\}$ is a nonempty polytope, so the LP has
an optimal vertex $y^{*}$. A vertex of $P\subseteq\mathbf{R}^n$ satisfies $n$
linearly independent constraints of $P$ with equality. Let
$F=\{j: 0<y^{*}_j<1\}$; these coordinates have $x_j=0$ and neither bound
tight. The bound constraints of coordinate $j$ are $y_j\ge 0$ and
$y_j\le 1-x_j$, with normals $\pm e_j$, so a linearly independent family of
tight constraints contains at most one bound constraint per coordinate
outside $F$ and none for coordinates in $F$. Hence at least $|F|$ of the $n$
independent tight constraints are rows of $W$, and $|F|\le m$. Let
$\tilde y_j=1$ if $y^{*}_j=1$ and $\tilde y_j=0$ otherwise. Then
$\tilde y\le y^{*}$, so $W\tilde y\le Wy^{*}\le c$ (as $W\ge 0$) and
$\tilde y\le\mathbf{1}-x$; $\tilde y$ is a binary follower response, and
$\Phi(x)\ge p^{\top}\tilde y=U(x)-\sum_{j\in F}p_jy^{*}_j\ge U(x)-m\,p_{\max}(x)$.
\end{proof}
\noindent Not formalised.

\begin{theorem}[Bracket]\label{thm:mkip-bracket}
Let $L\le U^{*}$ be any lower bound on $U^{*}$ (for example a MILP solver's
proven dual bound for the MILP of Proposition~\ref{prop:mkip-milp}) and let
$\hat x$ be any leader-feasible choice. Then
\[
  \lceil L-m\,p_{\max}\rceil\ \le\ \mathrm{OPT}\ \le\ \Phi(\hat x)\ \le\ U(\hat x).
\]
\end{theorem}
\begin{proof}
For every leader-feasible $x$, Lemma~\ref{lem:mkip-gap} gives
$\Phi(x)\ge U(x)-m\,p_{\max}(x)\ge U^{*}-m\,p_{\max}\ge L-m\,p_{\max}$, so
$\mathrm{OPT}\ge L-m\,p_{\max}$, and $\mathrm{OPT}$ is an integer because $p$ is.
The upper bounds hold because the leader may play $\hat x$ and the game is
zero-sum.
\end{proof}
\noindent Not formalised.

For the OR instances with one follower knapsack (variant $-100$, $m=1$) the
lower bound is $\lceil L-p_{\max}\rceil$. With several follower rows the width
$m\,p_{\max}$ makes the lower bound weak, and only the upper bounds are used.

\subsection{The merged-row online bound}

Theorem~\ref{thm:online} needs one leader row and one follower row. For one
follower row ($m=1$, weights $w_j=W_{1j}$, capacity $c=c_1$) and several leader
rows, choose multipliers $\mu\in\mathbf{R}^{m_u}_{\ge 0}$ and scales
$s_u,s_l>0$, and define the integer data
\[
  \gamma'_j=\Bigl\lfloor s_u\sum_i\mu_iG_{ij}\Bigr\rfloor,\quad
  \beta'=\Bigl\lceil s_u\sum_i\mu_ig_i\Bigr\rceil,\quad
  w'_j=\lceil s_lw_j\rceil,\quad
  c'=\lfloor s_lc\rfloor .
\]
Let $\omega'$ be the value of the online game of Theorem~\ref{thm:online}
with profits $p$, leader costs $\gamma'$, budget $\beta'$, follower weights $w'$
and capacity $c'$, on any list of distinct items.

\begin{theorem}[Merged-row online bound]\label{thm:mkip-omega}
For all $\mu\ge 0$ and $s_u,s_l>0$: $\omega'\le\mathrm{OPT}$.
\end{theorem}
\begin{proof}
\emph{The leader is relaxed.} If $Gx\le g$ then, as $x\ge 0$ and
$\gamma'_j\le s_u\sum_i\mu_iG_{ij}$,
\[
  \sum_j\gamma'_jx_j\le s_u\sum_i\mu_i(Gx)_i\le s_u\sum_i\mu_ig_i\le\beta'.
\]
So every leader-feasible set is feasible for the merged row $(\gamma',\beta')$.
\emph{The follower is restricted.} If $\sum_{j\in Y}w'_j\le c'$ then
$s_l\,w(Y)\le\sum_{j\in Y}w'_j\le c'\le s_lc$, so $w(Y)\le c$. Hence every
packing for $(w',c')$ is a packing for $(w,c)$, and the best packing
$K'(T)$ of a set $T$ under $(w',c')$ satisfies $K'(T)\le K(T)$.
\emph{Combine.} For every leader-feasible $X$, Theorem~\ref{thm:online} for
the data $(p,\gamma',\beta',w',c')$ gives
$\omega'\le K'(N\setminus X)\le K(N\setminus X)=\Phi(X)$. Minimising over $X$ gives
$\omega'\le\mathrm{OPT}$.
\end{proof}
\noindent The online bound itself is formalised (Theorem~\ref{thm:online});
the merging, scaling and rounding steps are not.

Lowering leader costs and raising the leader's budget only enlarge the
leader's options in the relaxed game, and raising follower weights and
lowering its capacity only shrink the follower's options; both moves can only
lower the value of the game, which is why the result stays a valid lower
bound for any multipliers and scales. The choice of $\mu$ (coordinate ascent
on a coarse table) and of the scales (which control the table size,
$(\beta'+1)(c'+1)$ entries per item) affects only the strength of the bound.
The code plays the items with positive profit in decreasing order of
$p_j/w_j$.

\subsection{Certified points}

\begin{proposition}[Certified point]\label{prop:mkip-point}
Let $\hat x$ satisfy $G\hat x\le g$, and let $\hat y$ be a follower response at
$\hat x$ proven optimal (by the exact knapsack oracle when $m=1$, or by a MILP
solver whose dual bound meets $p^{\top}\hat y$). Then $(\hat x,\hat y)$ is bilevel
feasible with value $\Phi(\hat x)$, so $\mathrm{OPT}\le\Phi(\hat x)$. Without an
exact follower solve, $\mathrm{OPT}\le\lfloor U(\hat x)\rfloor$ still holds.
\end{proposition}
\begin{proof}
$\hat y\in S(\hat x)$ by definition; the game is zero-sum, so every element of
$S(\hat x)$ gives the leader the value $\Phi(\hat x)$. The last claim is
Lemma~\ref{lem:mkip-gap} and integrality of $\Phi$.
\end{proof}
\noindent Not formalised.

For several follower knapsacks the leader choice $\hat x$ is the minimiser
found for the MILP of Proposition~\ref{prop:mkip-milp}; when HiGHS proves the
follower's optimum at $\hat x$, the point is certified. For one follower
knapsack, a first-improvement local search on the leader's choice (add an
item the follower uses, or swap a blocked item for one it uses, keeping
every budget row) evaluates each candidate with the exact knapsack oracle, so
every incumbent it reports is a certified point. Every point reported is
checked again from its file by \texttt{bobi verify}
(Section~\ref{app:verification}).

% Edge clique interdiction: the model, Mantel, partitions, Brouwer, the
% part-size reduction, and the three open instances.
\section{Clique interdiction proofs}\label{app:clique}

\subsection{The model}

BOBILib's instance \texttt{normalClique$n$-$d$-$i$} has a graph $G=(V,E)$ with
$n=|V|$, $m=|E|$, and budget $k=\lfloor m/4\rfloor$. All variables are binary.
The leader has $w_e$ for $e\in E$ and the single row $\sum_ew_e\le k$. The
follower has $x_v$ ($v\in V$) and $y_e$ ($e\in E$) and the rows
\[
  x_u+x_v\le 1\ \ (uv\notin E),\qquad
  x_u+x_v-y_{uv}\le 1\ \ (uv\in E),\qquad
  w_{uv}+y_{uv}\le 1\ \ (uv\in E).
\]
The leader minimises $\sum_vx_v$ and the follower minimises $-\sum_vx_v$. (Every row,
coefficient and objective of the 60 files was checked against this pattern,
once by the instance reader of the clique scripts and once by the parser of a
separate review session, which classifies rows by their coefficients.)

\begin{lemma}[Follower's problem]\label{lem:clique-model}
Fix the leader's choice and let $D=\{e: w_e=1\}$. A vector $x\in\{0,1\}^V$
extends to a feasible follower point if and only if $\{v: x_v=1\}$ is a clique
of $G-D$. Hence the follower's optimum is $\omega(G-D)$, the game is zero-sum,
and the bilevel optimum is $\min\{\omega(G-D): D\subseteq E,\ |D|\le k\}$.
\end{lemma}
\begin{proof}
Let $x$ be part of a feasible point and $x_u=x_v=1$. The first family of rows
forces $uv\in E$. If $uv\in D$ then $w_{uv}=1$ forces $y_{uv}=0$, and the
second family gives $x_u+x_v\le 1$, a contradiction. So $uv\in E\setminus D$.
Conversely, for a clique of $G-D$ set $y_{uv}=1-w_{uv}$: the third family holds
with equality; the second holds for $uv\notin D$ because $y_{uv}=1$, and for
$uv\in D$ because a clique of $G-D$ does not contain both $u$ and $v$; the
first holds because a clique contains no non-edge of $G$. The leader's
objective is minus the follower's.
\end{proof}
\noindent Not formalised (the Lean file states the bounds directly in terms of
$\omega(G-D)$).

\subsection{Lower bound 3: Mantel's theorem}

\begin{theorem}[Mantel 1907, deletion form]\label{thm:mantel}
A graph on $n$ vertices with more than $\lfloor n^2/4\rfloor$ edges contains a
triangle. Consequently, if $m-k>\lfloor n^2/4\rfloor$, then $G-D$ contains a
triangle for every set $D$ of at most $k$ edges, and the bilevel optimum is
at least~$3$.
\end{theorem}
\begin{proof}
Mantel's theorem\cite{mantel1907problem} is Tur\'an's theorem\cite{turan1941extremal}
for $K_3$. After deleting at most $k$
edges at least $m-k$ edges remain, and $m-k>\lfloor n^2/4\rfloor$.
\end{proof}
{\raggedright\noindent Formalised: \texttt{Bobi/Clique.lean}, theorems \texttt{mantel},
\texttt{mantel'} and \texttt{triangle\_survives}; the instance arithmetic in
\texttt{mantel\_40}, \texttt{mantel\_50}, \texttt{mantel\_60},
\texttt{clique40}, \texttt{clique50} and \texttt{clique60}.\par}

All 30 density-0.7 instances have $(n,m,k)\in\{(40,546,136),(50,857,214),(60,1239,309)\}$,
so $m-k=410,643,930$ exceeds $\lfloor n^2/4\rfloor=400,625,900$. The density-0.9
instances satisfy the same inequality.

\subsection{Upper bounds: partitions}

\begin{proposition}[Partition upper bound]\label{prop:clique-partition}
Let $f:V\to\{1,\dots,r\}$ and let $D$ contain every edge of $G$ whose ends have
the same colour. Then $G-D$ has no clique on $r+1$ vertices. If at most $k$
edges of $G$ lie inside the colour classes, deleting exactly those edges is a
leader choice, and the bilevel optimum is at most~$r$.
\end{proposition}
\begin{proof}
$f$ is a proper colouring of $G-D$, and a clique on $r+1$ vertices needs
$r+1$ colours.
\end{proof}
{\raggedright\noindent Formalised: \texttt{Bobi/Clique.lean}, theorem
\texttt{cliqueFree\_of\_partition}.\par}

For a partition $\mathcal{P}$ of $V$ into $r$ parts (parts may be empty) let
$\mathrm{in}_G(\mathcal{P})$ be the number of edges of $G$ inside parts, and let
$w_r(G)=\min_{\mathcal{P}}\mathrm{in}_G(\mathcal{P})$. Each solution file deletes the inside edges
of a 3- or 4-partition found by tabu search and states a maximum clique of the
remaining graph, computed exactly (the review re-checks it with two exact
clique searches\cite{bron1973algorithm,tomita2006worst,carraghan1990exact});
the files pass the checks of Section~\ref{app:verification}.

\subsection{Lower bound 4: Brouwer's theorem}

Let $t_r(n)$ be the number of edges of the balanced complete $r$-partite graph
on $n$ vertices (Tur\'an's number), and
$T(n,r)=\binom n2-t_r(n)$, the number of edges of the disjoint union of $r$
cliques of balanced sizes.

\begin{theorem}[Brouwer 1981, Theorem~2\cite{brouwer1981lotto}; quoted, with his $G$ renamed $H$ and his $m:=\lfloor n/t\rfloor$ written out]\label{thm:brouwer}
Let $H$ be a graph on $n$ points such that
$\alpha(H)\le t$, and $n>2t$. If $H$ has no more than $T(n,t)+\lfloor n/t\rfloor-2$
edges, then $V(H)$ can be written as a union of $t$ cliques. Moreover, this
bound is best possible.
\end{theorem}

Here $\alpha(H)$ is the independence number. The theorem is restated by
Tyomkyn and Uzzell\cite{tyomkyn2015strong} (their equation~(1) and
Theorem~1.1), who credit Brouwer with the first explicit proof, and, in the
complemented form below for $t\ge 2$, by Kor\'andi, Roberts and
Scott\cite{korandi2021exact} (their Theorem~1.1). Brouwer's report is not
refereed, but the statement is: Tyomkyn and Uzzell give a new proof of it in
their Section~2, Hanson and Toft and Kang and Pikhurko found it
again,\cite{hanson1991ksaturated,kang2005maximum} and Amin, Faudree, Gould and
Sidorowicz\cite{amin2013non} determine the largest number of edges of a maximal
$K_p$-free graph that is not $(p-1)$-partite, which gives the same threshold
(a $K_p$-free graph that is not $(p-1)$-partite stays so when edges are added
until it is maximal). We use the theorem in complemented form.

\begin{corollary}[Complement form]\label{cor:brouwer}
Let $t\ge 1$ and $n\ge 2t+1$. Every $K_{t+1}$-free graph $F$ on $n$ vertices with
$e(F)\ge t_t(n)-\lfloor n/t\rfloor+2$ is $t$-partite (its vertex set is the
union of $t$ independent sets, some possibly empty).
\end{corollary}
\begin{proof}
Apply Theorem~\ref{thm:brouwer} to the complement $H=\overline F$:
$\alpha(H)=\omega(F)\le t$, $n>2t$, and
$e(H)=\binom n2-e(F)\le\binom n2-t_t(n)+\lfloor n/t\rfloor-2=T(n,t)+\lfloor n/t\rfloor-2$.
So $V$ is a union of $t$ cliques of $H$, that is, of $t$ independent sets of
$F$. Assigning each vertex to one of the sets containing it gives a partition
into independent sets.
\end{proof}
\noindent Not formalised: Brouwer's theorem is not in Mathlib, and it is
neither proved nor assumed in the Lean files. The statement above was read
in the scanned report, and the translation was checked by computer (for
$t=3$: $T(n,3)+t_3(n)=\binom n2$).

\begin{theorem}[Density 0.9: value at least 4]\label{thm:clique-lb4}
Let $n\ge 7$, $m-k\ge t_3(n)-\lfloor n/3\rfloor+2$ and $w_3(G)>k$. Then $G-D$
contains a $K_4$ for every $D\subseteq E$ with $|D|\le k$, so the bilevel optimum is
at least~$4$.
\end{theorem}
\begin{proof}
Suppose $|D|\le k$ and $G-D$ is $K_4$-free. It has $m-|D|\ge m-k$ edges, so by
Corollary~\ref{cor:brouwer} with $t=3$ ($n\ge 7=2t+1$) it is 3-partite with some
partition $\mathcal{P}$. No edge of $G-D$ lies inside a part, so every edge of $G$
inside a part lies in $D$, and $|D|\ge\mathrm{in}_G(\mathcal{P})\ge w_3(G)>k$, a
contradiction.
\end{proof}
\noindent Not formalised (it rests on Theorem~\ref{thm:brouwer}).

For the density-0.9 instances $(n,m,k)\in\{(40,702,175),(50,1102,275),(60,1593,398)\}$,
so $m-k=527,827,1195$ and $t_3(n)-\lfloor n/3\rfloor+2=522,819,1182$
($t_3(n)=533,833,1200$): the hypothesis holds with margins $5$, $8$ and $13$,
while $m-k<t_3(n)$, so Tur\'an's theorem alone does not give a $K_4$. It
remains to show $w_3(G)>k$.

\subsection{Reducing $w_3(G)>k$ to finitely many part-size cases}

Let $\overline E$ be the edge set of the complement $\overline G$, so
$|\overline E|=\binom n2-m$, and for a partition $\mathcal{P}$ let
$\mathrm{cut}_{\overline G}(\mathcal{P})$ be the number of edges of $\overline G$ joining
different parts.

\begin{proposition}[Part-size reduction]\label{prop:clique-sizes}
For a partition $\mathcal{P}=(P_1,P_2,P_3)$ of $V$ with part sizes $s_1,s_2,s_3$,
\[
  \mathrm{in}_G(\mathcal{P})=\sum_{i=1}^3\binom{s_i}2-|\overline E|+\mathrm{cut}_{\overline G}(\mathcal{P}).
\]
Hence, with $b(s)=k+|\overline E|-\sum_i\binom{s_i}2$, $w_3(G)>k$ holds if and only
if, for every size vector $s_1\ge s_2\ge s_3\ge 0$ with $s_1+s_2+s_3=n$ and
$b(s)\ge 0$, no partition with these part sizes has
$\mathrm{cut}_{\overline G}(\mathcal{P})\le b(s)$.
\end{proposition}
\begin{proof}
Each of the $\sum_i\binom{s_i}2$ pairs inside a part is an edge of $G$ or of
$\overline G$, and the edges of $\overline G$ inside parts number
$|\overline E|-\mathrm{cut}_{\overline G}(\mathcal{P})$. So $\mathrm{in}_G(\mathcal{P})\le k$ iff
$\mathrm{cut}_{\overline G}(\mathcal{P})\le b(s)$. Relabelling the parts, every partition
has sorted sizes, and since $\mathrm{cut}_{\overline G}\ge 0$, size vectors with
$b(s)<0$ admit no partition with $\mathrm{in}_G\le k$.
\end{proof}
\noindent Not formalised (checked in exact integer arithmetic by script, and
on random partitions of every instance).

Since $n$, $m$ and $k$ are the same for all ten graphs of a group, so are the
size vectors with $b(s)\ge 0$:
\begin{center}
\small
\begin{tabular}{lll}
\hline
$n$ & $|\overline E|$ & size vectors $s$ with $b(s)$ in brackets\\
\hline
40 & 78 & (14,13,13)\,[6], (14,14,12)\,[5], (15,13,12)\,[4], (15,14,11)\,[2], (16,12,12)\,[1], (16,13,11)\,[0]\\
50 & 123 & (17,17,16)\,[6], (18,16,16)\,[5], (18,17,15)\,[4], (19,16,15)\,[2], (18,18,14)\,[1], (19,17,14)\,[0]\\
60 & 177 & (20,20,20)\,[5], (21,20,19)\,[4], (21,21,18)\,[2], (22,19,19)\,[2], (22,20,18)\,[1]\\
\hline
\end{tabular}
\end{center}
That is $10\cdot 6+10\cdot 6+10\cdot 5=170$ cases (instance, size vector), each
asking for a partition of $V$ with prescribed part sizes and at most $b(s)\le 6$
edges of the sparse complement between parts. All 170 were decided
infeasible by five methods: CP-SAT\cite{perron2026cpsat} and
SCIP\cite{hojny2025scip} on two different models, an exact
solver-free depth-first search (tested against brute force), and, in a
separate review session with its own code, HiGHS on a third model with an explicit budget row and
a separate exact branch and bound written in C (tested against brute force
on 450 random graphs and on planted instances). The three solver verdicts are
solver-certified (class~(ii), Section~\ref{app:classes}). For an exact
certificate, each case also has a branch-and-bound tree
(\texttt{results/clique/certificates/}), checked by the integer checker
\texttt{scripts/check\_clique\_certificates.py}: it reads the instance files
itself, re-derives the graph, $n$, $m$ and $k$ from the rows, checks the
hypothesis $m-k\ge t_3(n)-\lfloor n/3\rfloor+2$, enumerates the size vectors
with $b(s)\ge 0$, and verifies at every leaf of each tree a lower bound on
$\mathrm{cut}_{\overline G}$ above $b(s)$, using only Python's integers and
no project code (\CliqueCheckerLines{} lines of code). It passes on all
\CliqueCertCases{} cases of the \CliqueCertInstances{} instances
(\texttt{check.log}): the trees have \CliqueCertNodes{} nodes and
\CliqueCertLeaves{} leaves in all, the largest \CliqueCertMaxNodes{} nodes, and
the check takes about \CliqueCertSeconds{}~s. With Theorem~\ref{thm:clique-lb4} and
Proposition~\ref{prop:clique-partition} (4-partitions with at most $k$ inside
edges), all 30 density-0.9 instances have value exactly~$4$: an exact
certificate (class~(i)) resting on Brouwer's published theorem, which is
neither formalised nor reproved here. The case computations are not
formalised proofs.

\subsection{The three open instances}

For normalClique60-70-2, -3 and -5 ($n=60$, $m=1239$, $k=309$) Mantel gives
value $\ge 3$ and a 4-partition gives value $\le 4$. Let $\tau_4(G)$ be the
fewest edges meeting every $K_4$ of $G$. By Lemma~\ref{lem:clique-model} the
value is $3$ if and only if $\tau_4(G)\le k$. Deleting the inside edges of a
3-partition destroys every $K_4$ (two of any four vertices share a part), so
$\tau_4(G)\le w_3(G)$, and $w_3(G)\le k$ would give value~$3$; the best
3-partitions found have $311$, $310$ and $310$ inside edges, against $k=309$.

The converse argument of Theorem~\ref{thm:clique-lb4} is not available
here. It would need every $K_4$-free subgraph with at least $m-k=930$ edges
to be 3-partite, but $930<t_3(60)-20+2=1182$, so Corollary~\ref{cor:brouwer}
does not apply, and its conclusion is false at this edge count: the disjoint
union of the wheel $W_5$ (a 5-cycle plus a vertex joined to all of it; 6
vertices, 10 edges, $K_4$-free, chromatic number~4) and any 3-partite graph
with $920$ edges on the other 54 vertices ($920\le t_3(54)=972$) is a
$K_4$-free graph on 60 vertices with 930 edges that is not 3-partite. So a
proof of value~4 must bound $\tau_4(G)$ itself from below. Linear
programming bounds on $\tau_4$ (rows for every $K_4$, for maximal cliques and
for vertex sets whose induced subgraph must lose a Tur\'an excess of edges)
stop near $296$, below the $310$ needed, and searches for a transversal with
at most $309$ edges that are not restricted to 3-partite remainders stop at
the best 3-partition found. The three instances stay open with
$3\le\text{value}\le 4$.

% General MIPLIB-derived instances: recession directions, the HPR bound,
% the HPR -> follower -> tie-break pipeline, coupling rows, tolerances.
\section{General instances: rays, HPR bound and pipelines}\label{app:general}

Notation is that of~\eqref{eq:bl}. In this section a follower row is written
$\ell_i\le C_ix+D_iy\le h_i$ with $\ell_i\in\mathbf{R}\cup\{-\infty\}$ and
$h_i\in\mathbf{R}\cup\{+\infty\}$, and the follower's variable bounds as
$\ell^y_j\le y_j\le h^y_j$, because the direction test below depends on which
sides are finite.

\subsection{Unbounded followers: no bilevel feasible point}

\begin{theorem}[Recession direction]\label{thm:ray}
Let $r\in\mathbf{Q}^{n_y}$ satisfy, for every follower row $i$, $D_ir\ge 0$ if
$\ell_i>-\infty$ and $D_ir\le 0$ if $h_i<+\infty$; for every follower variable
$j$, $r_j\ge 0$ if $\ell^y_j>-\infty$ and $r_j\le 0$ if $h^y_j<+\infty$; and
$f^{\top}r<0$. Then for every leader choice $x$ (real or integral) the
follower's problem $P(x)$ is infeasible or unbounded, so $S(x)=\emptyset$ and
the bilevel problem has no feasible point.
\end{theorem}
\begin{proof}
Let $\nu$ be a positive integer such that $\nu r_j\in\mathbf{Z}$ for every $j\in I_l$
($r$ is rational). Suppose $y$ is feasible for $P(x)$. For every integer
$t\ge 0$, $y+t\nu r$ satisfies every follower row, since $C_ix$ is a constant
and $D_i(y+t\nu r)$ moves only in the permitted direction; it satisfies every
bound for the same reason; and its integer-constrained entries stay integral.
Its objective $f^{\top}y+t\nu f^{\top}r$ tends to $-\infty$. So $P(x)$ has no
optimal solution whenever it is feasible.
\end{proof}
\noindent Not formalised.

\begin{remark}[Meyer's theorem]
The theorem only needs an explicit rational direction. It is the case we
need of a classical result of Meyer:\cite{meyer1974existence} with rational data, the convex hull of the feasible set
of a mixed-integer program is a rational polyhedron with the same recession
cone as the LP relaxation, so a feasible mixed-integer program whose LP
relaxation is unbounded is itself unbounded. Rationality matters: without it,
a direction of the LP relaxation need not be usable by integer points. The
directions here are found by an LP solver in floating point, written as
decimals (which are rationals) and checked exactly.
\end{remark}

\paragraph{The exact check.} For each certificate, a script that reads the
instance files independently of the solver code converts every entry of $r$
to an exact rational and checks the conditions of Theorem~\ref{thm:ray} in
rational arithmetic, counting only follower variables in each follower row
(leader variables enter those rows as constants). When rounding leaves tiny
residuals, the direction is first projected onto the exact null space of its
tight rows. This succeeds for \NumRayExact{} of the \NumNoFeasible{} open instances, whose
certificates are therefore exact (class~(i), Section~\ref{app:classes}); for
\texttt{supportcase42\_0\_100} and \texttt{supportcase42\_50\_50} residuals of
about $10^{-12}$ (relative) remain, so those two certificates rest on floating
point (class~(ii)).
% ROUND2: update when results/rays has exact certificates for supportcase42_0_100 and
% supportcase42_50_50 (exact.ray.json checked by scripts/verify_ray.py): then all
% \NumNoFeasible{} are class (i) and the last clause goes.

\subsection{The HPR bound and the MIPLIB optimum}

\begin{lemma}[HPR bound]\label{lem:hpr}
Every bilevel feasible point is feasible for the HPR, so
$\mathrm{HPR}\le\mathrm{OPT}$. If a bilevel feasible point has value
$\mathrm{HPR}$, it is optimal.
\end{lemma}
\begin{proof}
A bilevel feasible point satisfies $Ax+By\ge a$, the integrality conditions,
and, as $y\in S(x)$, also $Cx+Dy\ge b$. The HPR minimises the same objective
over a superset.
\end{proof}
\noindent Not formalised.

BOBILib builds its general instances from a MIPLIB problem by assigning
variables and rows to the two levels and keeping the problem's objective as
the leader's. If nothing else changes, dropping the condition $y\in S(x)$
gives back the MIPLIB problem: the HPR \emph{is} the MIPLIB model, and its
published optimum\cite{gleixner2021miplib} $z^{*}$ is a lower bound on $\mathrm{OPT}$ by
Lemma~\ref{lem:hpr}. This is not assumed: a script reads the bilevel
instance and the MIPLIB~2017 model and checks in exact rational arithmetic
that they have the same variables, rows, coefficients, sides, bounds,
integrality requirements and objective (up to the objective's sense). For
the five instances where it holds and a verified bilevel feasible point
attains the bound (\texttt{pigeon-10\_0\_100}, \texttt{pigeon-10\_50\_50},
\texttt{timtab1\_0\_100}, \texttt{noswot\_0\_100}, \texttt{noswot\_50\_50}),
the point is optimal by accepted benchmark conventions. The values $z^{*}$
are the optima in MIPLIB~2017's solution file. They were computed by MILP
solvers in floating point and are solver-certified numerical results
(class~(ii), Section~\ref{app:classes}), not exact-arithmetic proofs, so the
five optima are no more rigorous than these published values:
\begin{itemize}
\item \texttt{noswot}: $z^{*}=-41.00000885$. The objective has integer
coefficients on integer variables only, so the HPR optimum is an integer, and
a valid lower bound $L$ gives the lower bound $\lceil L\rceil$. The bound
$\lceil z^{*}\rceil=-41$ therefore holds unless the published optimum
overstates the HPR optimum by $0.99999$ or more (the HPR optimum would have to
be $-42$ or less); BOBILib's own points attain it. Integrality rounding makes
this the safest of the five.
\item \texttt{pigeon-10}: $z^{*}=-9{,}000$, attained by both points; the bound
inherits MIPLIB's optimal status for the instance.
\item \texttt{timtab1}: $z^{*}=764{,}771.99999978$ (printed as $764{,}772$ on
the MIPLIB web pages); the point attains it with row violations below
$10^{-12}$, so the agreement of point and bound is numerical.
\end{itemize}
``Optimal'' for these five therefore means equal to the published optimum to
its printed precision (after rounding up for \texttt{noswot}), as BOBILib and
MIPLIB use the word; no exact-arithmetic proof of the bound is claimed. The
model comparison is exact; the points are class~(ii) (Section~\ref{app:classes}).

\subsection{Feasible points without coupling rows}

\begin{proposition}[HPR, follower, tie-break]\label{prop:general-pipeline}
Assume there are no coupling rows ($B=0$). Let $(\hat x,\hat y)$ be feasible for
the HPR, and suppose $P(\hat x)$ has an optimal solution, with value
$\varphi(\hat x)$. Then the tie-break problem at $\hat x$ (the optimistic choice
among the follower's optima, Section~\ref{app:notation}) is feasible, and for
any of its feasible points $y^{*}$, the point $(\hat x,y^{*})$ is bilevel
feasible. If $y^{*}$ is optimal for the tie-break problem, $(\hat x,y^{*})$ is the
best bilevel feasible point with leader part~$\hat x$.
\end{proposition}
\begin{proof}
$A\hat x\ge a$ holds because $(\hat x,\hat y)$ is HPR feasible and $B=0$, and
$\hat x$ satisfies its integrality conditions. Any optimal solution of
$P(\hat x)$ is feasible for the tie-break problem, whose rows are those of
$P(\hat x)$, $f^{\top}y\le\varphi(\hat x)$, and $A\hat x+By\ge a$, which does not involve
$y$. A feasible point $y^{*}$ of it is feasible for $P(\hat x)$ with
$f^{\top}y^{*}\le\varphi(\hat x)$, so $y^{*}\in S(\hat x)$. The bilevel feasible points
with leader part $\hat x$ are exactly the pairs $(\hat x,y)$ with $y\in S(\hat x)$, the
feasible set of the tie-break problem.
\end{proof}
\noindent Not formalised. The existence of an optimal solution of $P(\hat x)$
is what the follower solve establishes; the HPR point $\hat y$ shows that
$P(\hat x)$ is feasible. The computations accept $f^{\top}y\le\varphi(\hat x)+10^{-7}(1+|\varphi(\hat x)|)$
in the tie-break, well inside BOBILib's relative tolerance~\eqref{eq:bobilib-tol}.
The pipeline handles continuous linking variables, which neither of
BOBILib's two solvers accepts.

\begin{sloppypar}
The tie-break uses this slack. For this supplement the \NumGeneralFiles{} files
of \texttt{results/general} were verified again with \texttt{bobi verify
--generic} (HiGHS as the follower) and the follower value of each file
compared with the follower optimum at the file's leader values. In 19 files
the follower value lies more than $10^{-8}$ (relative) above that optimum:
about $10^{-7}$ in 18 of them (\texttt{binkar10\_1}, \texttt{cmflsp50-24-8-8},
\texttt{core2536-691}, \texttt{cost266-UUE}, \texttt{exp-1-500-5-5},
\texttt{ic97\_potential}, \texttt{n4-3}, \texttt{n5-3}, \texttt{n9-3},
\texttt{neos-2978193-inde}, \texttt{neos-5188808-nattai}, \texttt{neos-873061},
\texttt{neos17}, \texttt{pg}, \texttt{pg5\_34}, \texttt{seymour1},
\texttt{tr12-30}, \texttt{uct-subprob}, all \texttt{\_0\_100}) and one unit,
$5\times10^{-8}$ relative, in \texttt{netdiversion\_0\_100}. These 19 files
are benchmark-accepted points (class~(iii), Section~\ref{app:classes}). The
instances themselves have bilevel feasible points with the same leader values,
since without coupling rows every optimal follower response is feasible for
the leader's rows (Proposition~\ref{prop:general-pipeline}); but the leader
values stated in the 19 files are not shown to be attainable with an exactly
optimal follower. The other files, among them the two optimal points of
Table~\ref{tab:hpr}, have the follower within $10^{-8}$ of its optimum
(class~(ii)).
\end{sloppypar}
% ROUND2: update when results/general records the follower error of each file
% (or the 19 files are re-solved with an exactly optimal follower). The 19 and
% the list come from a re-verification run for this supplement (bobi verify
% --generic, time limit 120 s, on all 64 files), not from the data package.

\subsection{Coupling rows}

With coupling rows ($B\neq 0$) the follower's optima at the HPR's leader
values may all violate them. Recall that $P(x)$, $\varphi(x)$ and $S(x)$ depend
on $x$ only through the linking variables $x_L$; write $\varphi(x_L)$. For fixed
values $\bar x_L$ define the \emph{best response}
\[
  \mathrm{BR}(\bar x_L)=\min\bigl\{c^{\top}x+d^{\top}y:\ x_L=\bar x_L,\ Ax+By\ge a,\
  Cx+Dy\ge b,\ f^{\top}y\le\varphi(\bar x_L),\ \text{integrality}\bigr\},
\]
in which the other leader variables are free.

\begin{proposition}[Best response and blocked leader choices]\label{prop:coupled}
Suppose $P(\bar x_L)$ has an optimal solution. Then the feasible points of
$\mathrm{BR}(\bar x_L)$ are exactly the bilevel feasible points with linking part
$\bar x_L$, and $\mathrm{BR}(\bar x_L)$ is the best value among them. If
$\mathrm{BR}(\bar x_L)$ is infeasible, no bilevel feasible point has linking part
$\bar x_L$ (the choice $\bar x_L$ is \emph{blocked}). Consequently
$\mathrm{OPT}=\inf_{\bar x_L}\mathrm{BR}(\bar x_L)$ over linking values for which
$P(\bar x_L)$ has an optimal solution.
\end{proposition}
\begin{proof}
$(x,y)$ is bilevel feasible iff it satisfies the upper-level rows and
integrality and $y\in S(x)$, and $y\in S(x)$ iff $y$ is feasible for $P(x_L)$ and
$f^{\top}y\le\varphi(x_L)$. These are the constraints of $\mathrm{BR}(x_L)$. If
$P(\bar x_L)$ has no optimal solution, $S(x)=\emptyset$ for every $x$ with these
linking values.
\end{proof}
\noindent Not formalised.

\begin{lemma}[Value-function cuts\cite{mitsos2010global,lozano2017value}]\label{lem:vcut}
Let $\hat y$ be feasible for $P(\hat x_L)$ with $f^{\top}\hat y=\hat f$. Every bilevel
feasible $(x,y)$ satisfies: $\hat y$ is infeasible for $P(x_L)$, or
$f^{\top}y\le\hat f$.
\end{lemma}
\begin{proof}
If $\hat y$ is feasible for $P(x_L)$ then $\varphi(x_L)\le\hat f$, and $y\in S(x)$ gives
$f^{\top}y=\varphi(x_L)$.
\end{proof}
\noindent Not formalised.

The search for an unblocked linking choice tries, in turn within a time
budget: the HPR's own $x_L$; best-response dynamics; the HPR with objective
$c^{\top}x+d^{\top}y+\rho f^{\top}y$ for increasing $\rho$; a penalty alternating
direction method after Kleinert and Schmidt;\cite{kleinert2021computing} moves of continuous linking
variables towards the interior; the HPR with the cuts of
Lemma~\ref{lem:vcut} (modelled with binary indicators, exact when every
follower row's linking part takes integer values); and no-good cuts on
integer linking variables. These are heuristics: none of them is needed for
correctness, and the ``blocked'' verdicts (solver tolerances included) only
steer the search; no instance is claimed to have no feasible point on that
basis. What is certified is each returned point, re-checked from scratch: all
rows, bounds and integrality within $10^{-6}$; a fresh follower solve whose
proven dual bound puts the point's follower value within relative $10^{-5}$ of
the optimum; and the optimistic tie-break at the point's leader values within
relative $10^{-5}$ by a proven bound. Each file then passes \texttt{bobi
verify}.

% ROUND2: update when the 9 coupled files in results/general-coupled are repaired or
% re-classified (they are class (iii) until then).
In \NumCoupledTolerance{} of the \NumCoupledFiles{} files with coupling rows the follower's value lies about
$10^{-7}$ (relative) above its optimum, which BOBILib's
test~\eqref{eq:bobilib-tol} accepts; these files are benchmark-accepted points
(class~(iii), Section~\ref{app:classes}), and the other files with coupling
rows are class~(ii). For seven of them a point with the
follower at its optimum (within $10^{-9}$ relative) exists at the same leader
values with a slightly worse leader objective. For \texttt{pg\_50\_50} and
\texttt{neos-2978193-inde\_50\_50} the coupling rows force the follower about
$10^{-7}$ above its optimum at the files' leader values; with the non-linking
leader variables free, exactly optimal responses exist, so both instances do
have bilevel feasible points, but the two files are feasible only within the
tolerance.

\paragraph{Few binary linking variables.} If the linking variables are
binary and few, Proposition~\ref{prop:coupled} gives the exact optimum by
enumerating all $\bar x_L$, solving $P(\bar x_L)$ and $\mathrm{BR}(\bar x_L)$ for
each. This reproduces the published optimum $299{,}577$ of Lozano, Bergman
and Cire\cite{lozano2026network} for \texttt{general30-30-10-20-20-8} (10
linking bits); with 20 bits it is too slow.

\subsection{A tolerance observation}

\begin{remark}[\texttt{general20-30-20-20-20-2}]\label{rem:tolerance}
Test~\eqref{eq:bobilib-tol} is relative. In the generalized interdiction
instance \texttt{general20-30-20-20-20-2} the follower's objective values are
near $10^8$, so a response up to about $10^4$ units from the follower's optimum
passes. In BOBILib's solution file for this instance the follower's value is
$106{,}899{,}984.14$ against the follower's optimum $106{,}899{,}953$ at the same
leader choice, $31$ units above it (relative error about $2.9\times10^{-7}$), and
rows are violated by up to $3.3\times10^{-4}$ (our verifier's record is in
\texttt{results/notes/}). The row violations pass the $10^{-4}$ tolerance of
the MIPLIB checker inside BOBILib's checker (Section~\ref{app:tolerances}) but
not the $10^{-6}$ of \texttt{bobi verify}, and the stated objective
\GenTolObjFile{} differs from the \GenTolObjComputed{} of the point, so
\texttt{bobi verify} rejects the file. Such a point is
not bilevel feasible in the sense of Definition~\ref{def:bilevel-feasible}:
with an exactly optimal follower response the leader's value at the same
leader choice can differ. A value accepted only under a relative tolerance
is therefore not proven attainable with an exactly optimal follower
response: the file is the example of a benchmark-accepted point
(class~(iii), Section~\ref{app:classes}).
\end{remark}

% Attack-protect network interdiction (inter-fire): model, closure lemma,
% monotonicity, tree dynamic program, exhaustive enumeration.
\section{Network interdiction}\label{app:network}

\subsection{The model as encoded in the files}

BOBILib's \texttt{inter-fire} instances (from Baggio, Carvalho, Lodi and
Tramontani,\cite{baggio2021multilevel} with the vaccinated nodes already removed) have a graph
$(V,E)$ and three binary variables per node $v$: $y_v$ (the leader attacks
$v$), $x_v$ (the follower protects $v$) and $a_v$ (node $v$ is saved). The rows
are
\begin{gather*}
  \sum_vy_v\le\Phi,\qquad
  \sum_vx_v\le\Lambda,\qquad
  a_v+y_v\le 1\ \ (v\in V),\\
  a_v-a_u-x_v\le 0\ \ \text{for both orientations }(u,v)\text{ of every edge }uv\in E.
\end{gather*}
The first row is the leader's; the others are the follower's. The leader
minimises $\sum_vw_va_v$ and the follower minimises $-\sum_vw_va_v$, with
$w_v=1$ in every instance (the arguments below only use $w\ge 0$). The reader of
\texttt{bobi fire} checks every row, bound and objective coefficient of each file against this
pattern and rejects anything else; an independent reviewer re-derived the
pattern with a separate parser. The game is zero-sum, so the optimistic
tie-break plays no part. In this section $\Phi$ is the attack budget.

\subsection{The follower's optimum as a closure}

For an attack $Y\subseteq V$ and a protection set $X\subseteq V$, call
$B\subseteq V$ \emph{closed} if $Y\subseteq B$ and, for every edge and each
ordering $(u,v)$ of its two ends, $u\in B$ and $v\notin X$ imply $v\in B$. The intersection of closed sets is
closed and $V$ is closed, so there is a smallest closed set $B(Y,X)$: the
nodes reached from $Y$ by fire spreading through unprotected nodes.

\begin{lemma}[Closure]\label{lem:fire-closure}
For fixed $y=\mathbf{1}_Y$ and $x=\mathbf{1}_X$, a vector $a\in\{0,1\}^V$ satisfies the
follower's rows if and only if $Z(a)=\{v: a_v=0\}$ is closed. Hence the
follower's best value for $X$ is $w(V\setminus B(Y,X))$, and its optimum against
the attack $Y$ is
\[
  \mathrm{sv}(Y)=\max\bigl\{w(V\setminus B(Y,X)) : X\subseteq V,\ |X|\le\Lambda\bigr\}.
\]
\end{lemma}
\begin{proof}
The rows $a_v+y_v\le 1$ say $Y\subseteq Z(a)$. The row $a_v\le a_u+x_v$ for the
orientation $(u,v)$ says: if $a_u=0$ and $x_v=0$ then $a_v=0$, that is, $u\in Z(a)$
and $v\notin X$ imply $v\in Z(a)$. So the feasible $a$ are exactly the indicator
vectors of complements of closed sets. Every closed set contains $B(Y,X)$,
and $B(Y,X)$ is closed, so with $w\ge 0$ the best choice is
$a=\mathbf{1}_{V\setminus B(Y,X)}$.
\end{proof}
\noindent Not formalised (checked against brute force over all $(a,x)$ on 400
random graphs with up to 6 nodes, and used by the reviewer's brute force).

\begin{corollary}[Monotonicity in the attack]\label{cor:fire-monotone}
If $Y\subseteq Y'$ then $\mathrm{sv}(Y')\le\mathrm{sv}(Y)$. Hence, when $|V|\ge\Phi$,
the minimum of $\mathrm{sv}$ over attacks with $|Y|\le\Phi$ is attained by an attack
with exactly $\Phi$ nodes.
\end{corollary}
\begin{proof}
$B(Y',X)$ is closed for $Y$ (it contains $Y'\supseteq Y$), so
$B(Y,X)\subseteq B(Y',X)$ and $w(V\setminus B(Y',X))\le w(V\setminus B(Y,X))$ for every
$X$. Any attack with fewer than $\Phi$ nodes can be enlarged.
\end{proof}
\noindent Not formalised. The enumeration below covers every attack size
from $0$ to $\Phi$, so the certificate does not depend on this corollary.
It explains why the least saved count falls with the attack size: for
the instance \texttt{tree-50\_1-5-5\_006} the minima for attack sizes
$0,\dots,5$ are $49,48,44,38,32,26$.

\subsection{The tree dynamic program}

Suppose the graph is a forest. Root every component and hang the roots below
a virtual root whose edges carry no constraint. Give every node $v$ a state
$\sigma_v=(x_v,a_v)\in\{0,1\}^2$; state $\sigma$ is \emph{allowed} at $v$ unless
$v\in Y$ and $a_v=1$; states $\sigma_p,\sigma_c$ of the ends of an edge $pc$ are
\emph{compatible} if $a_c\le a_p+x_c$ and $a_p\le a_c+x_p$. For a node $v$ with
children $c_1,\dots,c_d$, a state $\sigma$ and $0\le q\le\Lambda$, let
$T_v(\sigma,q)$ be the largest $\sum_uw_ua_u$ over the subtree of $v$ among
assignments of allowed states with $v$ in state $\sigma$, compatible states on
every edge of the subtree and exactly $q$ protected nodes ($-\infty$ if there
is none). With $T^{(0)}_v(\sigma,q)=w_va(\sigma)$ if $\sigma$ is allowed and
$q=x(\sigma)$, and $-\infty$ otherwise,
\[
  T^{(i)}_v(\sigma,q)=\max_{0\le q'\le q}\Bigl[T^{(i-1)}_v(\sigma,q-q')+
  \max_{\sigma'\ \text{compatible with}\ \sigma}T_{c_i}(\sigma',q')\Bigr],\qquad
  T_v=T^{(d)}_v .
\]

\begin{proposition}[Tree DP]\label{prop:fire-dp}
$T_v$ satisfies the definition above for every node, and the follower's
optimum against $Y$ is the best value at the virtual root over $q\le\Lambda$.
\end{proposition}
\begin{proof}
Induction over the nodes in post-order. The follower's rows other than the
budget row involve one node ($a_v+y_v\le 1$) or the two ends of one edge (the
two orientations of the arc row), so an assignment of the subtree of $v$ is
feasible iff all states are allowed and every edge is compatible. The edges
of the subtree of $v$ are the edges $vc_i$ and the edges inside the subtrees of
the $c_i$, which share no node. Hence a feasible assignment of the subtree of
$v$ with $v$ in state $\sigma$ and $q$ protections is the same as: a state for
$v$, and for each $i$ a feasible assignment of the subtree of $c_i$ whose root
state is compatible with $\sigma$, with protection counts adding up to $q$; and
its value is the sum of the parts. The recursion maximises over exactly these
choices, one child at a time. At the virtual root every child state is
compatible, and the budget row is $q\le\Lambda$.
\end{proof}
\noindent Not formalised (tested against brute force over all $(a,x)$ on 300
random forests, and compared at every attack with a separate C implementation
by the reviewer: 9,344,902 attacks on the 40 forests).

\subsection{Exhaustive enumeration and the certificate}

\begin{theorem}[Certified optimum]\label{thm:fire-cert}
Suppose that for every attack $Y$ with $|Y|\le\Phi$ we have a follower point
$(x^Y,a^Y)$ that satisfies all follower rows at $Y$ (checked against the raw
model), with value $\mathrm{dp}(Y)=w(a^Y)$. Let $Y^{*}$ minimise $\mathrm{dp}$, and
suppose the follower's optimum at $Y^{*}$ is proven to equal $\mathrm{dp}(Y^{*})$.
Then the bilevel optimum is $\mathrm{dp}(Y^{*})$, attained by $Y^{*}$ and
$(x^{Y^{*}},a^{Y^{*}})$.
\end{theorem}
\begin{proof}
For every attack, $\mathrm{sv}(Y)\ge\mathrm{dp}(Y)\ge\mathrm{dp}(Y^{*})$, since the
follower may play the checked point. The leader's feasible set is exactly
$\{Y:|Y|\le\Phi\}$, so the optimum is at least $\mathrm{dp}(Y^{*})$; and
$\mathrm{sv}(Y^{*})=\mathrm{dp}(Y^{*})$ shows it is attained.
\end{proof}
\noindent Not formalised.

The certificate does not rest on the DP being optimal at any attack other
than $Y^{*}$, where HiGHS re-solves the follower MILP built from the raw
model. For \texttt{tree-50\_1-5-5\_006} and \texttt{\_010} the graph has 49
nodes and $\Phi=\Lambda=5$; the enumeration visits every $k$-subset for
$k=0,\dots,5$ (the count per size is checked against $\binom{49}k$),
$\sum_{k=0}^5\binom{49}k=2{,}138{,}410$ attacks per instance, and both optima are
$26$. The reviewer obtained the same minima, the same number of optimal
attacks and the same per-size minima with independent code.

% Map from appendix results to Lean theorems and to the Rust code that relies
% on them, and the list of what is not formalised.
\section{Formal proofs: map and scope}\label{app:lean-map}

The Lean~4 project \texttt{bobilib/lean} (with Mathlib) proves the lemmas in
the tables below. Lean file names are relative to \texttt{bobilib/lean}; Rust
names are in \texttt{bobilib/src} (\texttt{Search}, \texttt{RobustTable} and
\texttt{play\_order} are in \texttt{robust\_ss.rs}; \texttt{Kip} is in
\texttt{kip.rs}). The Lean statements are about abstract models (items,
weight classes, graphs); the Rust code is not verified, and the right-hand
column says which code relies on each statement.

\begin{center}
\footnotesize
\setlength{\tabcolsep}{4pt}
\begin{tabular}{p{1.35in}p{2.65in}p{2.0in}}
\hline
Result & Lean theorems (file) & Rust code that relies on it\tabularnewline
\hline
\multicolumn{3}{l}{\emph{Knapsack interdiction (Section~\ref{app:kip})}}\tabularnewline[3pt]
\raggedright Lemma~\ref{lem:kip-monotone} (monotonicity) &
\raggedright \texttt{Bobi/Knapsack.lean}: \texttt{follower\_antitone}, \texttt{bestAvoiding\_antitone}, \texttt{bestAvoiding\_eq\_sup\_sdiff}, \texttt{packings\_downward\_closed}, \texttt{K\_eq\_bestAvoiding} &
\raggedright \texttt{Kip::free\_items}, \texttt{Kip::greedy}, \texttt{Kip::subset\_sum\_classes}; guaranteed copies in \texttt{Search::node}\tabularnewline[3pt]
\raggedright Lemma~\ref{lem:kip-cheapest} (count formulation) &
\raggedright \texttt{Bobi/Knapsack.lean}: \texttt{K\_map\_equiv}, \texttt{swap\_block}, \texttt{cheapest\_first} &
\raggedright \texttt{Kip::subset\_sum\_classes}, \texttt{Kip::x\_from\_deletions}; branching on counts in \texttt{Search::node}\tabularnewline[3pt]
\raggedright Lemmas~\ref{lem:window-extend}, \ref{lem:window-survivors} (windows, stall test) &
\raggedright \texttt{Bobi/SubsetSum.lean}: \texttt{window\_extend}, \texttt{window\_extend\_all}, \texttt{fits\_survivors}, \texttt{window\_survivors}, \texttt{SubsetSum.perm}, \texttt{Window.perm} &
\raggedright \texttt{Search::window\_safe}\tabularnewline[3pt]
\raggedright Theorem~\ref{thm:online} (online bound) &
\raggedright \texttt{Bobi/Online.lean}: \texttt{online\_le\_K}, \texttt{online\_le\_K\_any}, \texttt{online\_lower\_bound}; example: \texttt{table1\_online}, \texttt{table1\_block}, \texttt{table1\_opt} &
\raggedright \texttt{mkip::omega\_root}\tabularnewline[3pt]
\raggedright Theorem~\ref{thm:tables-sound} (table soundness) &
\raggedright \texttt{Bobi/RobustTable.lean}: \texttt{H\_sound}, \texttt{H\_sound\_any}, \texttt{take\_sum\_le}, \texttt{counts\_of\_deletion}; \texttt{Bobi/BandTable.lean}: \texttt{G\_sound}, \texttt{G\_sound\_any}, \texttt{G\_one}, \texttt{G\_sound\_one}; \texttt{Bobi/EndTable.lean}: \texttt{GA\_sound}, \texttt{GA\_sound\_any}, \texttt{G\_iff\_GA}, \texttt{H\_iff\_GA}, \texttt{G\_sound\_of\_GA}, \texttt{H\_sound\_of\_GA} &
\raggedright \texttt{RobustTable::build}, \texttt{RobustTable::build\_band}\tabularnewline
\hline
\end{tabular}
\end{center}

\begin{center}
\footnotesize
\setlength{\tabcolsep}{4pt}
\begin{tabular}{p{1.35in}p{2.65in}p{2.0in}}
\hline
Result & Lean theorems (file) & Rust code that relies on it\tabularnewline
\hline
\multicolumn{3}{l}{\emph{Knapsack interdiction, continued}}\tabularnewline[3pt]
\raggedright Proposition~\ref{prop:play-order} (play orders) &
\raggedright \texttt{Bobi/EndTable.lean}: \texttt{counts\_perm}, \texttt{remaining\_eq\_flatten}, \texttt{totalCost\_eq\_sum}, \texttt{remaining\_perm}, \texttt{totalCost\_perm}, \texttt{GA\_sound\_perm}, \texttt{GA\_sound\_perm\_any} &
\raggedright \texttt{play\_order}, \texttt{RobustTable::build\_suffix}, \texttt{Search::node\_table\_threshold}\tabularnewline[3pt]
\raggedright Corollary~\ref{cor:root-rules} (root-table rules) &
\raggedright \texttt{Bobi/RobustTable.lean}: \texttt{covers\_sound}; \texttt{Bobi/BandTable.lean}: \texttt{band\_covers\_sound}, \texttt{covers\_sound\_of\_band}; \texttt{Bobi/Threshold.lean}: \texttt{covers\_threshold\_sound}, \texttt{covers\_threshold\_sound\_perm}, \texttt{band\_covers\_threshold\_sound}, \texttt{band\_covers\_threshold\_sound\_perm} &
\raggedright \texttt{RobustTable::covers}, \texttt{RobustTable::covers\_run}, \texttt{Search::tables\_cover}\tabularnewline[3pt]
\raggedright Lemma~\ref{lem:antitone} (antitone in the budget) &
\raggedright \texttt{Bobi/EndTable.lean}: \texttt{GA\_antitone}, \texttt{G\_antitone}, \texttt{H\_antitone}, \texttt{GA\_mono} &
\raggedright threshold storage in \texttt{RobustTable}; cache entries of \texttt{Search::remember}\tabularnewline[3pt]
\raggedright Theorem~\ref{thm:thr} (threshold form) &
\raggedright \texttt{Bobi/Threshold.lean}: \texttt{thr} (definition), \texttt{thr\_iff}, \texttt{GA\_iff\_thr}, \texttt{thr\_bounds}, \texttt{thr\_eq\_neg\_one\_iff}, \texttt{thr\_greatest}, \texttt{G\_iff\_thr}, \texttt{H\_iff\_thr}, \texttt{le\_minOver}, \texttt{le\_maxOver} &
\raggedright \texttt{RobustTable::step}, \texttt{RobustTable::base}, \texttt{RobustTable::threshold}, \texttt{RobustTable::contains}, \texttt{RobustTable::step\_range}\tabularnewline[3pt]
\raggedright Theorem~\ref{thm:node-rule} (per-node rule) &
\raggedright \texttt{Bobi/EndTable.lean}: \texttt{nodeEnd} (definition), \texttt{node\_rule}, \texttt{node\_table\_sound}, \texttt{node\_table\_sound\_perm}, \texttt{node\_table\_sound\_any}; \texttt{Bobi/Threshold.lean}: \texttt{node\_table\_threshold\_sound}, \texttt{node\_table\_threshold\_sound\_counts} &
\raggedright \texttt{Search::node\_table\_threshold}, the prune and cache entry in \texttt{Search::node}, \texttt{Search::remember}\tabularnewline[3pt]
\raggedright Remark after Theorem~\ref{thm:node-rule} (strength) &
\raggedright \texttt{Bobi/EndTable.lean}: \texttt{node\_table\_of\_H} &
\raggedright none (strength, not soundness)\tabularnewline
\hline
\end{tabular}
\end{center}

\begin{center}
\footnotesize
\setlength{\tabcolsep}{4pt}
\begin{tabular}{p{1.35in}p{2.65in}p{2.0in}}
\hline
Result & Lean theorems (file) & Rust code that relies on it\tabularnewline
\hline
\multicolumn{3}{l}{\emph{Edge clique interdiction (Section~\ref{app:clique})}}\tabularnewline[3pt]
\raggedright Theorem~\ref{thm:mantel} (Mantel, deletion form) &
\raggedright \texttt{Bobi/Clique.lean}: \texttt{mantel}, \texttt{mantel'}, \texttt{triangle\_survives}, \texttt{mantel\_40}, \texttt{mantel\_50}, \texttt{mantel\_60}, \texttt{clique40}, \texttt{clique50}, \texttt{clique60} &
\raggedright none; the arithmetic is also checked by \texttt{scripts/clique/lb.py}\tabularnewline[3pt]
\raggedright Proposition~\ref{prop:clique-partition} (partition bound) &
\raggedright \texttt{Bobi/Clique.lean}: \texttt{cliqueFree\_of\_partition} &
\raggedright none; the certificates are written by \texttt{scripts/clique/ub.py}\tabularnewline
\hline
\end{tabular}
\end{center}

All names in the middle column are in namespace \texttt{Bobi}. There is no
\texttt{sorry}; \texttt{lean/Check.lean} prints the axioms of 72 main
theorems, each of which uses at most \texttt{propext},
\texttt{Classical.choice} and \texttt{Quot.sound}.

\paragraph{Not formalised.}
{\raggedright
\begin{itemize}
\item \emph{Subset-sum search.} The reduction of
Proposition~\ref{prop:kip-reduction} as a whole (free and unusable items, the
gcd scaling; \texttt{Kip::subset\_sum\_classes}); the fractional deletion
bound, Lemma~\ref{lem:fractional} (\texttt{Cumulative::min\_cost}), the exact
largest deletable mass (\texttt{Search::build\_max\_del}) and the reduction of
the stall test to one test per class in Proposition~\ref{prop:window-rule};
restricted copies, Proposition~\ref{prop:restrict}
(\texttt{restrict\_copies}); Lemma~\ref{lem:node-state},
Theorems~\ref{thm:kip-rules} and~\ref{thm:kip-exact} and
Corollary~\ref{cor:kip-loop} as statements about the code
(\texttt{Search::node}, \texttt{Search::run}, \texttt{Kip::solve\_subset\_sum});
that the bitset holds exactly the decided subset sums up to $hi$; that the
Rust loops compute $\mathrm{thr}$, including the range shortcuts and early
exit of \texttt{RobustTable::step\_range} and \texttt{Search::node\_table\_threshold},
the fallback to \texttt{RobustTable::step} for classes with more than 63
copies, the \texttt{row\_max} and \texttt{blockmax} skips, and the absence of
overflow in the \texttt{i32} arithmetic (all tested against a direct
transcription of $\mathrm{thr}$ and against brute force); and the containment
of the band rules in the per-node rule.
\item \emph{Multidimensional knapsack interdiction.}
Lemmas~\ref{lem:mkip-dual} and~\ref{lem:mkip-gap},
Proposition~\ref{prop:mkip-milp} and Theorem~\ref{thm:mkip-bracket}
(\texttt{PackingInterdiction::lp\_follower\_bound} in \texttt{mkip.rs}); the
merging, scaling and rounding of Theorem~\ref{thm:mkip-omega}
(\texttt{mkip::omega\_root}, \texttt{mkip::tune\_mu},
\texttt{mkip::omega\_scales}); the local search; all floating-point
evaluation.
\item \emph{Clique interdiction.} Lemma~\ref{lem:clique-model}; Brouwer's
theorem (Theorem~\ref{thm:brouwer}) and Corollary~\ref{cor:brouwer}, which are
neither proved nor assumed in Lean; Theorem~\ref{thm:clique-lb4};
Proposition~\ref{prop:clique-sizes} and the 170 solver verdicts; that the
instance files have the stated $n$, $m$ and $k$ and that each certificate
partition has at most $k$ inside edges (Lean takes these numbers as
hypotheses).
\item \emph{Network interdiction.} Lemma~\ref{lem:fire-closure},
Corollary~\ref{cor:fire-monotone}, Proposition~\ref{prop:fire-dp}
(\texttt{fire::TreeDp::solve}) and Theorem~\ref{thm:fire-cert}
(\texttt{fire::enumerate}).
\item \emph{General instances.} Theorem~\ref{thm:ray}
(\texttt{general::follower\_unbounded\_ray}, checked exactly by
\texttt{scripts/verify\_ray.py}) and Meyer's theorem; Lemma~\ref{lem:hpr} and
the model comparison (\texttt{scripts/compare\_hpr.py});
Proposition~\ref{prop:general-pipeline} (\texttt{general::feasible\_point},
\texttt{general::optimal\_response}); Proposition~\ref{prop:coupled} and
Lemma~\ref{lem:vcut} (\texttt{coupled::Search::evaluate},
\texttt{coupled::Search::best\_response}, \texttt{coupled::certify}).
\end{itemize}
\par}

% The data package bobilib/results/: layout, instances.csv, the .solu file,
% solution and certificate formats, checksums. Source: results/README.md.
\section{Data package and formats}\label{app:data}

Every result in the paper is backed by a file in the data package, the
folder \texttt{bobilib/results/}. The package's index,
\texttt{results/README.md}, follows the repository style of the INFORMS
Journal on Computing (cite, description, data, file formats, verification,
replication). The index, \texttt{instances.csv},
\texttt{bobilib\_update.solu} and \texttt{SHA256SUMS} are written by
\texttt{paper/scripts/build\_data\_package.py}; \texttt{make data} in
\texttt{bobilib/paper/} regenerates them, together with the macros of every
number quoted in the paper and the tables of Section~\ref{app:tables}. BOBILib's
statuses are those of its metadata snapshot of 18~October~2026.

\subsection{Layout}

\begin{center}
\small
\begin{tabular}{@{}p{1.5in}p{4.8in}@{}}
\toprule
Path & Contents\\
\midrule
\texttt{instances.csv} & one row per BOBILib instance (\NumInstances{} rows); Section~\ref{app:data-csv}\\
\texttt{bobilib\_update.solu} & status changes and improved values in the MIPLIB/SCIP \texttt{.solu} format; Section~\ref{app:data-solu}\\
\texttt{SHA256SUMS} & SHA-256 checksum of every other file of the package\\
\texttt{summary.json} & per-instance outcome and per-family counts\\
\texttt{kip/} & \NumKipFiles{} knapsack interdiction solution files at the published optima of Weninger and Fukasawa\cite{weninger2025fast}\\
\texttt{ss/} & the \NumNewKip{} subset-sum optima proven here; \texttt{proofs.csv} (optimum, nodes, CPU seconds, nodes of the independent prover)\\
\texttt{or/} & \NumOrPoints{} verified points for the \NumOr{} OR instances; the bounds of each instance; \texttt{compare.csv} (published bounds, BOBILib's value, ours, gaps)\\
\texttt{clique/} & \NumCliqueSet{} partition certificates, solver logs of the part-size proofs, the independent review\\
\texttt{network/} & \NumVerifyNetwork{} attack-protect trees certified by enumeration; \texttt{summary.csv}; the review\\
\texttt{rays/} & \NumRayFiles{} recession directions (the \NumNoFeasible{} open instances and one that BOBILib lists as infeasible); \texttt{ray\_verify.jsonl} (exact checks)\\
\texttt{general/} & \NumGeneralFiles{} points for MIPLIB-derived instances without coupling rows (nine files over 1\,MB in \texttt{large\_points.tar.gz}); \texttt{hpr\_optimal.csv}\\
\texttt{general-coupled/} & \NumCoupledFiles{} points for MIPLIB-derived instances with coupling rows; \texttt{summary.csv}\\
\bottomrule
\end{tabular}
\end{center}

\noindent Each folder has its own \texttt{README.md} with the method, the
columns of its tables, the commands that produced it and, where one was done,
its independent review.

\subsection{The instance table \texttt{instances.csv}}\label{app:data-csv}

The file is UTF-8, comma separated, with one header row and one row per
BOBILib instance. Its columns:

\begin{center}
\small
\begin{tabular}{@{}p{1.5in}p{4.8in}@{}}
\toprule
Column & Meaning\\
\midrule
\texttt{instance} & BOBILib instance name\\
\texttt{collection} & BOBILib sub-collection (last directory of the instance path), e.g.\ \texttt{kp}, \texttt{or}, \texttt{miplib2017}\\
\texttt{family} & \texttt{knapsack}, \texttt{multidim-knapsack}, \texttt{clique}, \texttt{network}, \texttt{generalized}, \texttt{assignment} (interdiction), \texttt{general-pure}, \texttt{general-mixed}\\
\texttt{bobilib\_status}, \texttt{bobilib\_value} & BOBILib's status and best known objective value (empty: none)\\
\texttt{new\_status} & \texttt{optimal} (proven, in published work or here), \texttt{infeasible\_bilevel} (no bilevel feasible point), \texttt{feasible} (first verified point; BOBILib had none), \texttt{unchanged}, \texttt{bobilib\_optimal}, \texttt{bobilib\_infeasible}\\
\texttt{primal\_value} & best upper-level value backed by a verified solution file (of this package or of BOBILib) or by a published optimum\\
\texttt{dual\_bound} & best proven lower bound (the leader minimises), if any\\
\texttt{source} & citation key of the published result (Table~\ref{tab:source-keys}), \texttt{this work}, or \texttt{ThuraufEtAl2026} (BOBILib's own value)\\
\texttt{method} & how the status or value was obtained\\
\texttt{certificate} & path, relative to \texttt{results/}, of the solution or certificate file; \texttt{ARCHIVE.tar.gz:FILE} names a member of an archive\\
\texttt{verified} & \texttt{yes}: the certificate passed its check (\texttt{bobi verify}, or the exact ray check); \texttt{no}: checked in floating point only (the two \texttt{supportcase42} directions); \texttt{na}: no certificate in this package\\
\texttt{notes} & free text, separated by semicolons\\
\bottomrule
\end{tabular}
\end{center}

\noindent Rows per \texttt{new\_status}: \NumBobiOptimal{} \texttt{bobilib\_optimal},
\NumBobiInfeasible{} \texttt{bobilib\_infeasible}, \NumOptimalNow{} \texttt{optimal},
\NumNoFeasible{} \texttt{infeasible\_bilevel}, \NumFirstPoint{} \texttt{feasible}
and \NumUnchanged{} \texttt{unchanged}. The \NumOptimalNow{} optimal rows are the
\NumLitOptimal{} published optima and the \NumNewOptimal{} proven here; together
with the \NumNoFeasible{} rows without a bilevel feasible point they are the
\NumSettled{} settled instances. The keys of the \texttt{source} column map to
the references as follows.

\begin{table}[ht]
\centering
\small
\caption{Citation keys of the \texttt{source} column of \texttt{instances.csv}.}\label{tab:source-keys}
\begin{tabular}{@{}lll@{}}
\toprule
Key in \texttt{source} & Reference & Optima\\
\midrule
\texttt{ThuraufEtAl2026} & BOBILib's own value\cite{thurauf2026bobilib} & \\
\texttt{WeningerFukasawa2025} & Weninger and Fukasawa\cite{weninger2025fast} & \NumLitKip{}\\
\texttt{LozanoBergmanCire2026} & Lozano, Bergman and Cire\cite{lozano2026network} & \NumLitLozano{}\\
\texttt{FischettiEtAl2019} & Fischetti, Ljubi\'c, Monaci and Sinnl\cite{fischetti2019interdiction} & \NumLitImkp{}\\
\texttt{FischettiEtAl2017} & Fischetti, Ljubi\'c, Monaci and Sinnl\cite{fischetti2017new} & \NumLitFischettiSeventeen{}\\
\texttt{VasquezLozanoVanHoeve2025} & V\'asquez, Lozano and van Hoeve\cite{vasquez2025single} & \NumLitVasquez{}\\
\texttt{BOBILibEarlierRelease} & an earlier release of BOBILib's metadata\cite{bobilib_website} & \NumLitEarlierRelease{}\\
\texttt{this work} & a certificate in this package & \NumNewOptimal{}\\
\bottomrule
\end{tabular}
\end{table}

\subsection{The status file \texttt{bobilib\_update.solu}}\label{app:data-solu}

The file follows the MIPLIB/SCIP convention, one line per instance whose
status differs from BOBILib's metadata or whose best value improves:
\begin{center}
\small
\begin{tabular}{@{}lll@{}}
\texttt{=opt=} & \emph{name value} & proven optimal value (published, or proven here)\\
\texttt{=inf=} & \emph{name} & no bilevel feasible point (the follower is unbounded for every leader choice)\\
\texttt{=best=} & \emph{name value} & best verified bilevel feasible point (a solution file here), not proven optimal\\
\texttt{=bestdual=} & \emph{name value} & best proven lower bound, where this work proves a better one than any published\\
\end{tabular}
\end{center}
It has \NumSoluOpt{} \texttt{=opt=} lines (\NumLitOptimal{} published,
\NumNewOptimal{} from this work), \NumSoluInf{} \texttt{=inf=}, \NumSoluBest{}
\texttt{=best=} and \NumSoluBestDual{} \texttt{=bestdual=} lines. The
\texttt{=bestdual=} lines are the \NumOrBetterLb{} OR instances with one
follower knapsack and the \NumCliqueOpen{} open clique instances; an instance
can have both a \texttt{=best=} and a \texttt{=bestdual=} line. Every value is
the upper-level objective $c^{\top}x+d^{\top}y$ of the minimising leader under
the optimistic convention, printed at full precision (integers without a
decimal point, other values as the shortest decimal that reads back as the
same double). Comment lines start with \texttt{\#}.

\subsection{Solution files}

Solution files (\texttt{*.res.json}) are in BOBILib's own format, so BOBILib's
checker and \texttt{bobi verify} read them unchanged. Each is a JSON object
with the fields \texttt{name}, \texttt{bilevel\_type} (always
\texttt{optimistic}), \texttt{status} (\texttt{optimal} or \texttt{feasible}),
\texttt{difficulty}, \texttt{objective\_value} (the upper-level objective,
which BOBILib minimises), and \texttt{upper\_level\_decisions} and
\texttt{lower\_level\_decisions}, which map every variable name of the
instance to its value. A file states \texttt{optimal} only when the instance's
row in \texttt{instances.csv} has \texttt{new\_status} \texttt{optimal}.

\subsection{Certificates}

\paragraph{Recession directions.} \texttt{rays/*.ray.json} maps follower
variable names to the entries of a direction $d$ (its support only), as
decimal numbers; \texttt{*.exact.ray.json} holds exact rationals written as
\texttt{"p/q"}. The certificate is valid if $d$ satisfies the conditions of
Theorem~\ref{thm:ray}: for every follower row, with the leader terms and
right-hand side dropped, $a^{\top}d\le 0$ for a row $a^{\top}y\le\cdot$,
$a^{\top}d\ge 0$ for a row $a^{\top}y\ge\cdot$ and $a^{\top}d=0$ for an
equation; $d_j\ge 0$ for a finite lower bound and $d_j\le 0$ for a finite upper
bound; and $f^{\top}d<0$. \texttt{scripts/verify\_ray.py} checks this in rational
arithmetic and writes \texttt{rays/ray\_verify.jsonl}.

\paragraph{Partitions.} The clique solution files are the certificates: the
leader deletes exactly the edges inside the parts of an $r$-partition (at most
$k$ of them), so the remainder is $r$-partite
(Proposition~\ref{prop:clique-partition}), and the follower part is a maximum
clique of the remainder. \texttt{clique/summary.csv} gives per instance $n$, the
density, $m$, $k$, BOBILib's value, the lower bound with its argument (Mantel,
or Brouwer's theorem with the part-size proofs) and the upper bound with its
partition. The logs of the solver runs behind the \NumCliquePartCases{}
part-size proofs are in \texttt{clique/logs/lb\_solver\_logs.tar.gz}.

\paragraph{Subset-sum proofs.} A proof of $\mathrm{OPT}\ge\ell$ is a
computation, not a file. \texttt{ss/proofs.csv} records, per instance, the
optimum, the nodes and CPU seconds of the search and the nodes of the
independent prover \texttt{scripts/ss\_check.py}, which re-proves the bound
from the instance files (Section~\ref{app:verification}).

\paragraph{Enumerations and HPR bounds.} \texttt{network/summary.csv} gives
per tree the attacks enumerated, the optimum, the number of optimal attacks
and the cross-checks (Theorem~\ref{thm:fire-cert}).
\texttt{general/hpr\_optimal.csv} lists the points that meet the MIPLIB
optimum of the HPR; \texttt{scripts/compare\_hpr.py} re-checks that the HPR is
the MIPLIB model (Lemma~\ref{lem:hpr}).

\paragraph{OR bounds.} \texttt{or/bounds\_one\_knapsack.jsonl} gives, per
instance with one follower knapsack, the lower bound with its two parts (merged-row
online bound, Theorem~\ref{thm:mkip-omega}, and LP bracket,
Theorem~\ref{thm:mkip-bracket}) and the upper bound;
\texttt{or/points\_several\_knapsacks.jsonl} gives, per instance with several
follower knapsacks, the bracket minimiser's LP bound and the certified
follower value.

\subsection{Checksums and how to check the package}

\texttt{SHA256SUMS} has one line per file of the package other than itself,
sorted by path, in the format of \texttt{sha256sum}. With BOBILib's instance
files in \texttt{data/collection} (from \texttt{bobilib.org}) and the
\texttt{bobi} crate built, the following commands, run from \texttt{bobilib/},
repeat the checks of Section~\ref{app:verification}:
\begin{center}
\begin{minipage}{0.96\linewidth}
\footnotesize
\begin{verbatim}
(cd results && sha256sum -c SHA256SUMS)
D=data/collection
./target/release/bobi verify --instances $D results/kip results/ss results/network results/or
./target/release/bobi verify --instances $D --generic results/clique results/general \
    results/general-coupled
python3 -I scripts/verify_ray.py INSTANCE.mps.gz INSTANCE.aux results/rays/INSTANCE.ray.json
python3 -I scripts/ss_check.py INSTANCE.mps.gz INSTANCE.aux - OPTIMUM
python3 -I scripts/compare_hpr.py INSTANCE.mps.gz INSTANCE.aux MIPLIB.mps.gz
\end{verbatim}
\end{minipage}
\end{center}
The nine large general points are checked after unpacking
\texttt{general/large\_points.tar.gz}.

% Per-instance tables. The tables themselves are generated by
% paper/scripts/gen_appendix_tables.py from bobilib/results/ (make data) and
% must not be edited by hand.
\section{Per-instance tables}\label{app:tables}

The tables below are generated from the data package by
\texttt{paper/scripts/gen\_appendix\_tables.py}. Values are upper-level
objectives of the minimising leader; full-precision values and the
certificate of every row are in \texttt{instances.csv}
(Section~\ref{app:data-csv}). The order follows the paper:
published optima (Table~\ref{tab:literature}); knapsack interdiction
(Table~\ref{tab:kip}); the OR set (Table~\ref{tab:or}); clique
interdiction (Table~\ref{tab:clique}); network interdiction
(Tables~\ref{tab:network} and~\ref{tab:network-trees}); general instances
without a feasible point (Table~\ref{tab:rays}), with an HPR-bound optimum
(Table~\ref{tab:hpr}) and with a first point (Table~\ref{tab:general-points});
and the instances whose status is unchanged (Table~\ref{tab:unchanged}).

\subsection{Published optima}
% Generated by paper/scripts/gen_appendix_tables.py from bobilib/results/; do not edit by hand.
% Needs booktabs, longtable, siunitx (v3) and multicol in the preamble.

{\footnotesize\sisetup{group-separator={,},group-minimum-digits=4,group-digits=integer}\setlength{\tabcolsep}{4pt}\setlength{\LTcapwidth}{\linewidth}
\begin{longtable}{@{}p{0.6\linewidth} l S[table-format=3.0]@{}}
\caption{Published optima for instances BOBILib lists as feasible or open, by source (sub-collection and count). The instance names follow.}\label{tab:literature}\\
\toprule
Source & Collections & {Instances} \\
\midrule
\endfirsthead
\multicolumn{3}{l}{\emph{(continued)}}\\
\toprule
Source & Collections & {Instances} \\
\midrule
\endhead
\midrule
\multicolumn{3}{r}{\emph{continued on the next page}}\\
\endfoot
\bottomrule
\endlastfoot
Weninger and Fukasawa (2025), combinatorial branch and bound & kp 336, inter-kp 25, cclw 15 & 376 \\
Lozano, Bergman and Cire (2026), decision-diagram relaxations & or 88, miplib3 5, generalized 4 & 97 \\
Fischetti, Ljubi\'c, Monaci and Sinnl (2019), interdiction cuts & imkp 52 & 52 \\
earlier BOBILib release (listed optimal there) & miplib2017 1 & 1 \\
Fischetti, Ljubi\'c, Monaci and Sinnl (2017), MIX++ & miplib3 1 & 1 \\
V\'asquez, Lozano and van Hoeve (2025), decision diagrams & miplib3 1 & 1 \\
\midrule
Total &  & 528 \\
\end{longtable}
}

\par\medskip\noindent\textbf{Weninger and Fukasawa (2025), combinatorial branch and bound (376)}\par\nopagebreak\smallskip
{\setlength{\multicolsep}{3pt plus 1pt}\begin{multicols}{4}
\footnotesize\raggedright\setlength{\parindent}{0pt}\setlength{\parskip}{0pt}
interdiction40-5\par
interdiction40-6\par
interdiction45-4\par
interdiction45-5\par
interdiction45-6\par
interdiction50-3\par
interdiction50-4\par
interdiction50-5\par
interdiction50-7\par
interdiction55-3\par
interdiction55-4\par
interdiction55-5\par
interdiction55-6\par
interdiction55-7\par
interdiction55-8\par
interKP-100-100-1-1\par
interKP-100-100-1-2\par
interKP-100-100-1-3\par
interKP-100-100-1-4\par
interKP-100-100-1-5\par
interKP-100-100-1-6\par
interKP-100-100-1-7\par
interKP-100-100-1-8\par
interKP-100-100-2-1\par
interKP-100-100-2-2\par
interKP-100-100-2-3\par
interKP-100-100-2-4\par
interKP-100-100-2-5\par
interKP-100-100-2-6\par
interKP-100-100-2-7\par
interKP-100-100-3-1\par
interKP-100-100-3-2\par
interKP-100-100-3-3\par
interKP-100-100-3-4\par
interKP-100-100-3-5\par
interKP-100-100-3-6\par
interKP-100-100-3-7\par
interKP-100-100-3-8\par
interKP-100-100-3-9\par
interKP-100-100-4-2\par
interKP-100-100-4-3\par
interKP-100-100-4-4\par
interKP-100-100-4-5\par
interKP-100-100-4-6\par
interKP-100-100-4-7\par
interKP-100-100-4-8\par
interKP-100-100-5-1\par
interKP-100-100-5-2\par
interKP-100-100-5-3\par
interKP-100-100-5-4\par
interKP-100-100-5-5\par
interKP-100-100-5-6\par
interKP-100-100-5-7\par
interKP-100-100-5-8\par
interKP-100-100-6-2\par
interKP-100-100-6-3\par
interKP-100-100-6-4\par
interKP-100-100-6-5\par
interKP-100-100-6-6\par
interKP-100-100-6-7\par
interKP-100-100-6-8\par
interKP-100-100-7-2\par
interKP-100-100-7-3\par
interKP-100-100-7-4\par
interKP-100-100-7-5\par
interKP-100-100-7-6\par
interKP-100-100-7-7\par
interKP-100-100-7-8\par
interKP-100-100-8-1\par
interKP-100-100-8-2\par
interKP-100-100-8-3\par
interKP-100-100-8-4\par
interKP-100-100-8-5\par
interKP-100-100-8-6\par
interKP-100-100-8-7\par
interKP-100-100-8-8\par
interKP-100-100-8-9\par
interKP-200-100-1-1\par
interKP-200-100-1-2\par
interKP-200-100-1-3\par
interKP-200-100-1-4\par
interKP-200-100-1-5\par
interKP-200-100-1-6\par
interKP-200-100-1-7\par
interKP-200-100-1-8\par
interKP-200-100-1-9\par
interKP-200-100-2-1\par
interKP-200-100-2-2\par
interKP-200-100-2-3\par
interKP-200-100-2-4\par
interKP-200-100-2-5\par
interKP-200-100-2-6\par
interKP-200-100-2-7\par
interKP-200-100-2-8\par
interKP-200-100-3-1\par
interKP-200-100-3-2\par
interKP-200-100-3-3\par
interKP-200-100-3-4\par
interKP-200-100-3-5\par
interKP-200-100-3-6\par
interKP-200-100-3-7\par
interKP-200-100-3-8\par
interKP-200-100-3-9\par
interKP-200-100-4-1\par
interKP-200-100-4-2\par
interKP-200-100-4-3\par
interKP-200-100-4-4\par
interKP-200-100-4-5\par
interKP-200-100-4-6\par
interKP-200-100-4-7\par
interKP-200-100-5-1\par
interKP-200-100-5-2\par
interKP-200-100-5-3\par
interKP-200-100-5-4\par
interKP-200-100-5-5\par
interKP-200-100-5-6\par
interKP-200-100-5-7\par
interKP-200-100-5-8\par
interKP-200-100-5-9\par
interKP-200-100-6-1\par
interKP-200-100-6-4\par
interKP-200-100-6-5\par
interKP-200-100-6-6\par
interKP-200-100-6-7\par
interKP-200-100-6-8\par
interKP-200-100-6-9\par
interKP-200-100-7-1\par
interKP-200-100-7-4\par
interKP-200-100-7-5\par
interKP-200-100-7-6\par
interKP-200-100-7-7\par
interKP-200-100-7-8\par
interKP-200-100-7-9\par
interKP-200-100-8-1\par
interKP-200-100-8-2\par
interKP-200-100-8-3\par
interKP-200-100-8-4\par
interKP-200-100-8-5\par
interKP-200-100-8-6\par
interKP-200-100-8-7\par
interKP-200-100-8-8\par
interKP-200-100-8-9\par
interKP-300-100-1-1\par
interKP-300-100-1-2\par
interKP-300-100-1-3\par
interKP-300-100-1-4\par
interKP-300-100-1-5\par
interKP-300-100-1-6\par
interKP-300-100-1-7\par
interKP-300-100-1-8\par
interKP-300-100-1-9\par
interKP-300-100-2-1\par
interKP-300-100-2-2\par
interKP-300-100-2-3\par
interKP-300-100-2-4\par
interKP-300-100-2-5\par
interKP-300-100-2-6\par
interKP-300-100-2-7\par
interKP-300-100-2-8\par
interKP-300-100-2-9\par
interKP-300-100-3-1\par
interKP-300-100-3-2\par
interKP-300-100-3-3\par
interKP-300-100-3-4\par
interKP-300-100-3-5\par
interKP-300-100-3-6\par
interKP-300-100-3-7\par
interKP-300-100-3-8\par
interKP-300-100-3-9\par
interKP-300-100-4-1\par
interKP-300-100-4-2\par
interKP-300-100-4-3\par
interKP-300-100-4-4\par
interKP-300-100-4-5\par
interKP-300-100-4-6\par
interKP-300-100-4-7\par
interKP-300-100-5-1\par
interKP-300-100-5-2\par
interKP-300-100-5-3\par
interKP-300-100-5-4\par
interKP-300-100-5-5\par
interKP-300-100-5-6\par
interKP-300-100-5-7\par
interKP-300-100-5-8\par
interKP-300-100-5-9\par
interKP-300-100-6-1\par
interKP-300-100-6-4\par
interKP-300-100-6-5\par
interKP-300-100-6-6\par
interKP-300-100-6-7\par
interKP-300-100-6-8\par
interKP-300-100-6-9\par
interKP-300-100-7-1\par
interKP-300-100-7-4\par
interKP-300-100-7-5\par
interKP-300-100-7-6\par
interKP-300-100-7-7\par
interKP-300-100-7-8\par
interKP-300-100-8-1\par
interKP-300-100-8-2\par
interKP-300-100-8-3\par
interKP-300-100-8-4\par
interKP-300-100-8-5\par
interKP-300-100-8-6\par
interKP-300-100-8-7\par
interKP-300-100-8-8\par
interKP-300-100-8-9\par
interKP-400-100-1-1\par
interKP-400-100-1-2\par
interKP-400-100-1-3\par
interKP-400-100-1-4\par
interKP-400-100-1-5\par
interKP-400-100-1-6\par
interKP-400-100-1-7\par
interKP-400-100-1-8\par
interKP-400-100-1-9\par
interKP-400-100-1-10\par
interKP-400-100-2-1\par
interKP-400-100-2-2\par
interKP-400-100-2-3\par
interKP-400-100-2-4\par
interKP-400-100-2-5\par
interKP-400-100-2-6\par
interKP-400-100-2-7\par
interKP-400-100-2-8\par
interKP-400-100-2-9\par
interKP-400-100-2-10\par
interKP-400-100-3-1\par
interKP-400-100-3-2\par
interKP-400-100-3-3\par
interKP-400-100-3-4\par
interKP-400-100-3-5\par
interKP-400-100-3-6\par
interKP-400-100-3-7\par
interKP-400-100-3-8\par
interKP-400-100-3-9\par
interKP-400-100-4-1\par
interKP-400-100-4-2\par
interKP-400-100-4-3\par
interKP-400-100-4-4\par
interKP-400-100-4-5\par
interKP-400-100-4-6\par
interKP-400-100-4-7\par
interKP-400-100-4-8\par
interKP-400-100-5-1\par
interKP-400-100-5-2\par
interKP-400-100-5-3\par
interKP-400-100-5-4\par
interKP-400-100-5-5\par
interKP-400-100-5-6\par
interKP-400-100-5-7\par
interKP-400-100-5-8\par
interKP-400-100-5-9\par
interKP-400-100-5-10\par
interKP-400-100-6-1\par
interKP-400-100-6-4\par
interKP-400-100-6-5\par
interKP-400-100-6-6\par
interKP-400-100-6-7\par
interKP-400-100-6-8\par
interKP-400-100-6-9\par
interKP-400-100-7-1\par
interKP-400-100-7-4\par
interKP-400-100-7-5\par
interKP-400-100-7-6\par
interKP-400-100-7-7\par
interKP-400-100-7-8\par
interKP-400-100-7-9\par
interKP-400-100-8-1\par
interKP-400-100-8-2\par
interKP-400-100-8-3\par
interKP-400-100-8-4\par
interKP-400-100-8-5\par
interKP-400-100-8-6\par
interKP-400-100-8-7\par
interKP-400-100-8-8\par
interKP-400-100-8-9\par
interKP-500-100-1-1\par
interKP-500-100-1-2\par
interKP-500-100-1-3\par
interKP-500-100-1-4\par
interKP-500-100-1-5\par
interKP-500-100-1-6\par
interKP-500-100-1-7\par
interKP-500-100-1-8\par
interKP-500-100-1-9\par
interKP-500-100-1-10\par
interKP-500-100-2-1\par
interKP-500-100-2-2\par
interKP-500-100-2-3\par
interKP-500-100-2-4\par
interKP-500-100-2-5\par
interKP-500-100-2-6\par
interKP-500-100-2-7\par
interKP-500-100-2-8\par
interKP-500-100-2-9\par
interKP-500-100-2-10\par
interKP-500-100-3-1\par
interKP-500-100-3-2\par
interKP-500-100-3-3\par
interKP-500-100-3-4\par
interKP-500-100-3-5\par
interKP-500-100-3-6\par
interKP-500-100-3-7\par
interKP-500-100-3-8\par
interKP-500-100-3-9\par
interKP-500-100-3-10\par
interKP-500-100-4-1\par
interKP-500-100-4-2\par
interKP-500-100-4-3\par
interKP-500-100-4-4\par
interKP-500-100-4-5\par
interKP-500-100-4-6\par
interKP-500-100-4-7\par
interKP-500-100-4-8\par
interKP-500-100-5-1\par
interKP-500-100-5-2\par
interKP-500-100-5-3\par
interKP-500-100-5-4\par
interKP-500-100-5-5\par
interKP-500-100-5-6\par
interKP-500-100-5-7\par
interKP-500-100-5-8\par
interKP-500-100-5-9\par
interKP-500-100-5-10\par
interKP-500-100-6-1\par
interKP-500-100-6-4\par
interKP-500-100-6-5\par
interKP-500-100-6-6\par
interKP-500-100-6-7\par
interKP-500-100-6-8\par
interKP-500-100-6-9\par
interKP-500-100-6-10\par
interKP-500-100-7-1\par
interKP-500-100-7-4\par
interKP-500-100-7-5\par
interKP-500-100-7-6\par
interKP-500-100-7-7\par
interKP-500-100-7-8\par
interKP-500-100-7-9\par
interKP-500-100-7-10\par
interKP-500-100-8-1\par
interKP-500-100-8-2\par
interKP-500-100-8-3\par
interKP-500-100-8-4\par
interKP-500-100-8-5\par
interKP-500-100-8-6\par
interKP-500-100-8-7\par
interKP-500-100-8-8\par
interKP-500-100-8-9\par
interKP-500-100-8-10\par
K5040W01.KNP\par
K5040W05.KNP\par
K5040W08.KNP\par
K5040W15.KNP\par
K5040W19.KNP\par
K5050W01.KNP\par
K5050W02.KNP\par
K5050W03.KNP\par
K5050W04.KNP\par
K5050W05.KNP\par
K5050W06.KNP\par
K5050W07.KNP\par
K5050W08.KNP\par
K5050W09.KNP\par
K5050W10.KNP\par
K5050W11.KNP\par
K5050W12.KNP\par
K5050W13.KNP\par
K5050W14.KNP\par
K5050W15.KNP\par
K5050W16.KNP\par
K5050W17.KNP\par
K5050W18.KNP\par
K5050W19.KNP\par
K5050W20.KNP\par
\end{multicols}}

\par\medskip\noindent\textbf{Lozano, Bergman and Cire (2026), decision-diagram relaxations (97)}\par\nopagebreak\smallskip
{\setlength{\multicolsep}{3pt plus 1pt}\begin{multicols}{4}
\footnotesize\raggedright\setlength{\parindent}{0pt}\setlength{\parskip}{0pt}
general30-30-10-20-20-8\par
general30-30-20-20-20-9\par
general30-30-30-20-20-9\par
general30-30-30-20-20-10\par
lseu-0.500000\par
OR5x100-0.75\_1.dat-0\par
OR5x100-0.75\_1.dat-50\par
OR5x100-0.75\_1.dat-100\par
OR5x100-0.75\_2.dat-0\par
OR5x100-0.75\_2.dat-50\par
OR5x100-0.75\_2.dat-100\par
OR5x100-0.75\_3.dat-0\par
OR5x100-0.75\_3.dat-50\par
OR5x100-0.75\_3.dat-100\par
OR5x100-0.75\_4.dat-0\par
OR5x100-0.75\_4.dat-50\par
OR5x100-0.75\_4.dat-100\par
OR5x100-0.75\_5.dat-0\par
OR5x100-0.75\_5.dat-50\par
OR5x100-0.75\_5.dat-100\par
OR5x100-0.75\_6.dat-0\par
OR5x100-0.75\_6.dat-50\par
OR5x100-0.75\_6.dat-100\par
OR5x100-0.75\_7.dat-0\par
OR5x100-0.75\_7.dat-50\par
OR5x100-0.75\_7.dat-100\par
OR5x100-0.75\_8.dat-0\par
OR5x100-0.75\_8.dat-50\par
OR5x100-0.75\_8.dat-100\par
OR5x100-0.75\_9.dat-0\par
OR5x100-0.75\_9.dat-50\par
OR5x100-0.75\_9.dat-100\par
OR5x100-0.75\_10.dat-0\par
OR5x100-0.75\_10.dat-50\par
OR5x100-0.75\_10.dat-100\par
OR5x250-0.75\_1.dat-0\par
OR5x250-0.75\_1.dat-50\par
OR5x250-0.75\_1.dat-100\par
OR5x250-0.75\_2.dat-0\par
OR5x250-0.75\_2.dat-50\par
OR5x250-0.75\_2.dat-100\par
OR5x250-0.75\_3.dat-0\par
OR5x250-0.75\_3.dat-50\par
OR5x250-0.75\_3.dat-100\par
OR5x250-0.75\_4.dat-0\par
OR5x250-0.75\_4.dat-50\par
OR5x250-0.75\_4.dat-100\par
OR5x250-0.75\_5.dat-0\par
OR5x250-0.75\_5.dat-50\par
OR5x250-0.75\_5.dat-100\par
OR5x250-0.75\_6.dat-0\par
OR5x250-0.75\_6.dat-50\par
OR5x250-0.75\_6.dat-100\par
OR5x250-0.75\_7.dat-0\par
OR5x250-0.75\_7.dat-50\par
OR5x250-0.75\_7.dat-100\par
OR5x250-0.75\_8.dat-0\par
OR5x250-0.75\_8.dat-50\par
OR5x250-0.75\_8.dat-100\par
OR5x250-0.75\_9.dat-0\par
OR5x250-0.75\_9.dat-50\par
OR5x250-0.75\_9.dat-100\par
OR5x250-0.75\_10.dat-0\par
OR5x250-0.75\_10.dat-50\par
OR5x250-0.75\_10.dat-100\par
OR10x100-0.75\_1.dat-100\par
OR10x100-0.75\_2.dat-0\par
OR10x100-0.75\_2.dat-100\par
OR10x100-0.75\_3.dat-100\par
OR10x100-0.75\_4.dat-0\par
OR10x100-0.75\_4.dat-100\par
OR10x100-0.75\_5.dat-100\par
OR10x100-0.75\_6.dat-0\par
OR10x100-0.75\_6.dat-100\par
OR10x100-0.75\_7.dat-0\par
OR10x100-0.75\_7.dat-100\par
OR10x100-0.75\_8.dat-0\par
OR10x100-0.75\_8.dat-100\par
OR10x100-0.75\_9.dat-0\par
OR10x100-0.75\_9.dat-100\par
OR10x100-0.75\_10.dat-100\par
OR10x250-0.75\_1.dat-100\par
OR10x250-0.75\_6.dat-100\par
OR10x250-0.75\_7.dat-100\par
OR30x100-0.75\_1.dat-100\par
OR30x100-0.75\_2.dat-100\par
OR30x100-0.75\_3.dat-100\par
OR30x100-0.75\_4.dat-100\par
OR30x100-0.75\_5.dat-0\par
OR30x100-0.75\_5.dat-100\par
OR30x100-0.75\_6.dat-100\par
OR30x100-0.75\_7.dat-100\par
OR30x100-0.75\_8.dat-100\par
p0201-0.100000\par
p0282-0.900000\par
p0548-0.100000\par
p0548-0.900000\par
\end{multicols}}

\par\medskip\noindent\textbf{Fischetti, Ljubi\'c, Monaci and Sinnl (2019), interdiction cuts (52)}\par\nopagebreak\smallskip
{\setlength{\multicolsep}{3pt plus 1pt}\begin{multicols}{5}
\footnotesize\raggedright\setlength{\parindent}{0pt}\setlength{\parskip}{0pt}
pet7.dat-0\par
pet7.dat-50\par
pet7.dat-100\par
sento2.dat-50\par
sento2.dat-100\par
weish14.dat-0\par
weish14.dat-50\par
weish14.dat-100\par
weish16.dat-0\par
weish16.dat-50\par
weish16.dat-100\par
weish17.dat-50\par
weish17.dat-100\par
weish18.dat-0\par
weish18.dat-50\par
weish18.dat-100\par
weish19.dat-0\par
weish19.dat-50\par
weish19.dat-100\par
weish20.dat-0\par
weish20.dat-50\par
weish20.dat-100\par
weish21.dat-0\par
weish21.dat-50\par
weish21.dat-100\par
weish22.dat-0\par
weish22.dat-50\par
weish22.dat-100\par
weish23.dat-0\par
weish23.dat-50\par
weish23.dat-100\par
weish24.dat-0\par
weish24.dat-50\par
weish24.dat-100\par
weish25.dat-0\par
weish25.dat-50\par
weish25.dat-100\par
weish26.dat-0\par
weish26.dat-50\par
weish26.dat-100\par
weish27.dat-0\par
weish27.dat-50\par
weish27.dat-100\par
weish28.dat-0\par
weish28.dat-50\par
weish28.dat-100\par
weish29.dat-0\par
weish29.dat-50\par
weish29.dat-100\par
weish30.dat-0\par
weish30.dat-50\par
weish30.dat-100\par
\end{multicols}}

\par\medskip\noindent\textbf{Earlier BOBILib release (listed optimal there) (1)}\par\nopagebreak\smallskip
{\setlength{\multicolsep}{3pt plus 1pt}\begin{multicols}{5}
\footnotesize\raggedright\setlength{\parindent}{0pt}\setlength{\parskip}{0pt}
neos-950242\_0\_100\par
\end{multicols}}

\par\medskip\noindent\textbf{Fischetti, Ljubi\'c, Monaci and Sinnl (2017), MIX++ (1)}\par\nopagebreak\smallskip
{\setlength{\multicolsep}{3pt plus 1pt}\begin{multicols}{5}
\footnotesize\raggedright\setlength{\parindent}{0pt}\setlength{\parskip}{0pt}
air03-0.100000\par
\end{multicols}}

\par\medskip\noindent\textbf{V\'asquez, Lozano and van Hoeve (2025), decision diagrams (1)}\par\nopagebreak\smallskip
{\setlength{\multicolsep}{3pt plus 1pt}\begin{multicols}{5}
\footnotesize\raggedright\setlength{\parindent}{0pt}\setlength{\parskip}{0pt}
p0282-0.500000\par
\end{multicols}}


\subsection{Knapsack interdiction}
% Generated by paper/scripts/gen_appendix_tables.py from bobilib/results/; do not edit by hand.
% Needs booktabs, longtable, siunitx (v3) and multicol in the preamble.

{\footnotesize\sisetup{group-separator={,},group-minimum-digits=4,group-digits=integer}\setlength{\tabcolsep}{4pt}\setlength{\LTcapwidth}{\linewidth}
\begin{longtable}{@{}l S[table-format=3.0] S[table-format=1.0] S[table-format=1.0] S[table-format=4.0] S[table-format=4.0] S[table-format=5.0] S[table-format=2.2] S[table-format=5.0]@{}}
\caption{The 16 knapsack interdiction instances without a published optimum, all proven optimal here by the robust subset-sum search. $n$: items; class 6: subset sum ($p = w$), class 7: even weights and odd capacity; budget index as in the instance name. Nodes and CPU seconds of \texttt{bobi ss-solve} on one thread of a shared machine; independent prover: nodes of \texttt{scripts/ss\_check.py}, which shares no code with the search. BOBILib's value was already optimal in every case.}\label{tab:kip}\\
\toprule
Instance & {$n$} & {Class} & {Budget} & {BOBILib} & {Optimum} & {Nodes} & {CPU (s)} & {\shortstack{Indep.\\nodes}} \\
\midrule
\endfirsthead
\multicolumn{9}{l}{\emph{(continued)}}\\
\toprule
Instance & {$n$} & {Class} & {Budget} & {BOBILib} & {Optimum} & {Nodes} & {CPU (s)} & {\shortstack{Indep.\\nodes}} \\
\midrule
\endhead
\midrule
\multicolumn{9}{r}{\emph{continued on the next page}}\\
\endfoot
\bottomrule
\endlastfoot
interKP-200-100-6-2 & 200 & 6 & 2 & 1855 & 1855 & 303 & 0.10 & 303 \\
interKP-300-100-6-2 & 300 & 6 & 2 & 2868 & 2868 & 366 & 0.16 & 366 \\
interKP-400-100-6-2 & 400 & 6 & 2 & 3832 & 3832 & 257 & 0.15 & 257 \\
interKP-500-100-6-2 & 500 & 6 & 2 & 4651 & 4651 & 190 & 0.26 & 190 \\
interKP-200-100-6-3 & 200 & 6 & 3 & 2782 & 2782 & 13498 & 3.93 & 13502 \\
interKP-300-100-6-3 & 300 & 6 & 3 & 4302 & 4302 & 19405 & 8.06 & 19537 \\
interKP-400-100-6-3 & 400 & 6 & 3 & 5748 & 5748 & 13430 & 9.01 & 13583 \\
interKP-500-100-6-3 & 500 & 6 & 3 & 6976 & 6976 & 10079 & 10.04 & 10116 \\
interKP-200-100-7-2 & 200 & 7 & 2 & 1872 & 1872 & 22 & 0.06 & 22 \\
interKP-300-100-7-2 & 300 & 7 & 2 & 2894 & 2894 & 27 & 0.04 & 27 \\
interKP-400-100-7-2 & 400 & 7 & 2 & 3866 & 3866 & 18 & 0.05 & 18 \\
interKP-500-100-7-2 & 500 & 7 & 2 & 4694 & 4694 & 1 & 0.05 & 1 \\
interKP-200-100-7-3 & 200 & 7 & 3 & 2808 & 2808 & 651 & 0.10 & 651 \\
interKP-300-100-7-3 & 300 & 7 & 3 & 4340 & 4340 & 470 & 0.20 & 470 \\
interKP-400-100-7-3 & 400 & 7 & 3 & 5800 & 5800 & 378 & 0.17 & 378 \\
interKP-500-100-7-3 & 500 & 7 & 3 & 7042 & 7042 & 256 & 0.15 & 256 \\
\end{longtable}
}


\subsection{The OR set}
% Generated by paper/scripts/gen_appendix_tables.py from bobilib/results/; do not edit by hand.
% Needs booktabs, longtable, siunitx (v3) and multicol in the preamble.

{\footnotesize\sisetup{group-separator={,},group-minimum-digits=4,group-digits=integer}\setlength{\tabcolsep}{3.5pt}\setlength{\LTcapwidth}{\linewidth}
\begin{longtable}{@{}l S[table-format=2.0] S[table-format=2.0] S[table-format=6.0] S[table-format=6.0] S[table-format=6.0] S[table-format=6.0] S[table-format=6.0] S[table-format=3.2]@{}}
\caption{All 810 multidimensional knapsack interdiction instances of the OR set (\texttt{OR}\emph{rows}\texttt{x}\emph{items}\texttt{-}\emph{budget}\texttt{\_}\emph{i}\texttt{.dat-}\emph{variant}). Budget: leader budget in percent of the item weights; follower rows: knapsack rows of the follower (variant -100: one; -0 and -50: several). Published LB: Lozano, Bergman and Cire (2026); published UB: best published point; BOBILib: BOBILib's best known value; our LB: dualised follower LP bracket or merged-row online bound (one follower knapsack only); our point: verified solution file; gap: (best upper bound $-$ best lower bound)/best upper bound in percent, where the best upper bound may be the follower LP bound at our leader choice (100: no lower bound known). Blank: none. $^{\dagger}$: the published lower bound exceeds the published point (inconsistent sources); the gap then leaves the published point out, and OR10x100-0.75\_4.dat-0 is closed at 8,813, the authors' branch-and-cut bound.}\label{tab:or}\\
\toprule
Instance & {Budget} & {\shortstack{Foll.\\rows}} & {\shortstack{Publ.\\LB}} & {\shortstack{Publ.\\UB}} & {BOBILib} & {\shortstack{Our\\LB}} & {\shortstack{Our\\point}} & {\shortstack{Gap\\(\%)}} \\
\midrule
\endfirsthead
\multicolumn{9}{l}{\emph{(continued)}}\\
\toprule
Instance & {Budget} & {\shortstack{Foll.\\rows}} & {\shortstack{Publ.\\LB}} & {\shortstack{Publ.\\UB}} & {BOBILib} & {\shortstack{Our\\LB}} & {\shortstack{Our\\point}} & {\shortstack{Gap\\(\%)}} \\
\midrule
\endhead
\midrule
\multicolumn{9}{r}{\emph{continued on the next page}}\\
\endfoot
\bottomrule
\endlastfoot
OR5x100-0.25\_1.dat-0 & 25 & 4 & 16014 & 20321 & 21879 &  & 19956 & 19.75 \\
OR5x100-0.25\_2.dat-0 & 25 & 4 & 15984 & 20174 & 22491 &  & 19943 & 19.85 \\
OR5x100-0.25\_3.dat-0 & 25 & 4 & 15180 & 19538 & 21738 &  & 19312 & 21.40 \\
OR5x100-0.25\_4.dat-0 & 25 & 4 & 16092 & 20786 & 22263 &  & 20353 & 20.94 \\
OR5x100-0.25\_5.dat-0 & 25 & 4 & 15645 & 20356 & 21965 &  & 19957 & 21.61 \\
OR5x100-0.25\_6.dat-0 & 25 & 4 & 16312 & 20774 & 22601 &  & 20235 & 19.39 \\
OR5x100-0.25\_7.dat-0 & 25 & 4 & 16257 & 20848 & 22905 &  & 20484 & 20.64 \\
OR5x100-0.25\_8.dat-0$^{\dagger}$ & 25 & 4 & 15874 & 10859 & 22343 &  & 19849 & 20.03 \\
OR5x100-0.25\_9.dat-0 & 25 & 4 & 16103 & 20661 & 22294 &  & 20382 & 20.99 \\
OR5x100-0.25\_10.dat-0 & 25 & 4 & 16095 & 20912 & 22698 &  & 20967 & 23.03 \\
OR5x100-0.25\_1.dat-50 & 25 & 2 & 17464 & 21887 & 23044 &  & 21569 & 19.03 \\
OR5x100-0.25\_2.dat-50 & 25 & 2 & 17711 & 21621 & 23112 &  & 21219 & 16.53 \\
OR5x100-0.25\_3.dat-50 & 25 & 2 & 17082 & 21405 & 22880 &  & 20930 & 18.39 \\
OR5x100-0.25\_4.dat-50 & 25 & 2 & 16472 & 21595 & 22781 &  & 21439 & 23.17 \\
OR5x100-0.25\_5.dat-50 & 25 & 2 & 16576 & 21031 & 21738 &  & 20621 & 19.62 \\
OR5x100-0.25\_6.dat-50 & 25 & 2 & 18143 & 22689 & 24519 &  & 22527 & 19.46 \\
OR5x100-0.25\_7.dat-50$^{\dagger}$ & 25 & 2 & 17106 & 16951 & 23302 &  & 22125 & 22.68 \\
OR5x100-0.25\_8.dat-50 & 25 & 2 & 17155 & 21697 & 22739 &  & 21330 & 19.57 \\
OR5x100-0.25\_9.dat-50 & 25 & 2 & 17061 & 21121 & 22661 &  & 20882 & 18.30 \\
OR5x100-0.25\_10.dat-50 & 25 & 2 & 18212 & 22642 & 23881 &  & 22543 & 19.21 \\
OR5x100-0.25\_1.dat-100 & 25 & 1 & 20026 & 24034 & 24489 & 23062 & 23482 & 1.79 \\
OR5x100-0.25\_2.dat-100 & 25 & 1 & 18986 & 22045 & 22393 & 21401 & 21750 & 1.60 \\
OR5x100-0.25\_3.dat-100 & 25 & 1 & 19494 & 22455 & 23303 & 21768 & 22110 & 1.55 \\
OR5x100-0.25\_4.dat-100 & 25 & 1 & 20713 & 24754 & 25202 & 24026 & 24260 & 0.96 \\
OR5x100-0.25\_5.dat-100 & 25 & 1 & 19603 & 23666 & 24344 & 22599 & 23026 & 1.85 \\
OR5x100-0.25\_6.dat-100 & 25 & 1 & 19923 & 23962 & 24483 & 23103 & 23431 & 1.40 \\
OR5x100-0.25\_7.dat-100 & 25 & 1 & 21157 & 25569 & 26905 & 24642 & 24958 & 1.27 \\
OR5x100-0.25\_8.dat-100 & 25 & 1 & 20006 & 23937 & 25286 & 23135 & 23514 & 1.61 \\
OR5x100-0.25\_9.dat-100 & 25 & 1 & 20298 & 24653 & 25199 & 23862 & 24184 & 1.33 \\
OR5x100-0.25\_10.dat-100 & 25 & 1 & 19389 & 23831 & 24840 & 22970 & 23242 & 1.17 \\
OR5x100-0.50\_1.dat-0 & 50 & 4 & 18975 & 23672 & 27489 &  & 19785 & 4.09 \\
OR5x100-0.50\_2.dat-0 & 50 & 4 & 19757 & 22821 & 27526 &  & 21265 & 7.09 \\
OR5x100-0.50\_3.dat-0 & 50 & 4 & 19204 & 23929 & 25707 &  & 21226 & 9.53 \\
OR5x100-0.50\_4.dat-0 & 50 & 4 & 20704 & 24066 & 30878 &  & 21779 & 4.94 \\
OR5x100-0.50\_5.dat-0 & 50 & 4 & 18942 & 22742 & 26240 &  & 21172 & 10.53 \\
OR5x100-0.50\_6.dat-0 & 50 & 4 & 17608 & 23765 & 30318 &  & 22019 & 20.03 \\
OR5x100-0.50\_7.dat-0 & 50 & 4 & 18963 & 22721 & 29946 &  & 21030 & 9.83 \\
OR5x100-0.50\_8.dat-0 & 50 & 4 & 20364 & 23945 & 29893 &  & 22283 & 8.61 \\
OR5x100-0.50\_9.dat-0 & 50 & 4 & 19037 & 24888 & 32863 &  & 23465 & 18.87 \\
OR5x100-0.50\_10.dat-0 & 50 & 4 & 19569 & 25360 & 27107 &  & 22596 & 13.40 \\
OR5x100-0.50\_1.dat-50 & 50 & 2 & 21207 & 30496 & 33475 &  & 29651 & 28.48 \\
OR5x100-0.50\_2.dat-50 & 50 & 2 & 20796 & 30461 & 32306 &  & 28254 & 26.40 \\
OR5x100-0.50\_3.dat-50 & 50 & 2 & 20033 & 30123 & 31051 &  & 28000 & 28.45 \\
OR5x100-0.50\_4.dat-50 & 50 & 2 & 21539 & 31945 & 33225 &  & 29334 & 26.57 \\
OR5x100-0.50\_5.dat-50 & 50 & 2 & 20100 & 29003 & 31008 &  & 27882 & 27.91 \\
OR5x100-0.50\_6.dat-50 & 50 & 2 & 19414 & 27930 & 30410 &  & 27331 & 28.97 \\
OR5x100-0.50\_7.dat-50 & 50 & 2 & 19888 & 28098 & 31826 &  & 26901 & 26.07 \\
OR5x100-0.50\_8.dat-50 & 50 & 2 & 21341 & 30971 & 34206 &  & 29952 & 28.75 \\
OR5x100-0.50\_9.dat-50 & 50 & 2 & 19847 & 31902 & 33819 &  & 29032 & 31.64 \\
OR5x100-0.50\_10.dat-50 & 50 & 2 & 19879 & 31436 & 31881 &  & 29152 & 31.81 \\
OR5x100-0.50\_1.dat-100 & 50 & 1 & 20708 & 30100 & 31172 & 28108 & 29437 & 4.51 \\
OR5x100-0.50\_2.dat-100 & 50 & 1 & 20894 & 29084 & 30477 & 27660 & 28530 & 3.05 \\
OR5x100-0.50\_3.dat-100 & 50 & 1 & 22308 & 29545 & 29887 & 27960 & 28767 & 2.81 \\
OR5x100-0.50\_4.dat-100 & 50 & 1 & 21678 & 31632 & 32359 & 29107 & 30499 & 4.56 \\
OR5x100-0.50\_5.dat-100 & 50 & 1 & 20966 & 28955 & 29142 & 27322 & 28269 & 3.35 \\
OR5x100-0.50\_6.dat-100 & 50 & 1 & 19588 & 28598 & 28746 & 26699 & 27621 & 3.34 \\
OR5x100-0.50\_7.dat-100 & 50 & 1 & 22087 & 28623 & 30947 & 26727 & 27915 & 4.26 \\
OR5x100-0.50\_8.dat-100 & 50 & 1 & 21521 & 31740 & 33174 & 29358 & 30424 & 3.50 \\
OR5x100-0.50\_9.dat-100 & 50 & 1 & 20773 & 30871 & 31618 & 28251 & 29163 & 3.13 \\
OR5x100-0.50\_10.dat-100 & 50 & 1 & 23071 & 30367 & 30448 & 28292 & 29232 & 3.22 \\
OR5x100-0.75\_1.dat-0 & 75 & 4 & 8855 & 8855 & 9350 &  & 8855 & 0.00 \\
OR5x100-0.75\_2.dat-0 & 75 & 4 & 11323 & 11323 & 12841 &  & 11323 & 0.00 \\
OR5x100-0.75\_3.dat-0 & 75 & 4 & 9344 & 9344 & 10231 &  & 9344 & 0.00 \\
OR5x100-0.75\_4.dat-0 & 75 & 4 & 9098 & 9098 & 10126 &  & 9098 & 0.00 \\
OR5x100-0.75\_5.dat-0 & 75 & 4 & 8492 & 8492 & 8895 &  & 8492 & 0.00 \\
OR5x100-0.75\_6.dat-0 & 75 & 4 & 8775 & 8775 & 10651 &  & 8775 & 0.00 \\
OR5x100-0.75\_7.dat-0 & 75 & 4 & 8596 & 8596 & 9311 &  & 8596 & 0.00 \\
OR5x100-0.75\_8.dat-0 & 75 & 4 & 9712 & 9712 & 10945 &  & 9712 & 0.00 \\
OR5x100-0.75\_9.dat-0 & 75 & 4 & 9212 & 9212 & 9964 &  & 9212 & 0.00 \\
OR5x100-0.75\_10.dat-0 & 75 & 4 & 8785 & 8785 & 9089 &  & 8785 & 0.00 \\
OR5x100-0.75\_1.dat-50 & 75 & 2 & 12477 & 12477 & 13783 &  & 12477 & 0.00 \\
OR5x100-0.75\_2.dat-50 & 75 & 2 & 13354 & 13354 & 14941 &  & 13354 & 0.00 \\
OR5x100-0.75\_3.dat-50 & 75 & 2 & 12774 & 12774 & 14365 &  & 12774 & 0.00 \\
OR5x100-0.75\_4.dat-50 & 75 & 2 & 13086 & 13086 & 15331 &  & 13086 & 0.00 \\
OR5x100-0.75\_5.dat-50 & 75 & 2 & 12907 & 12907 & 15056 &  & 12907 & 0.00 \\
OR5x100-0.75\_6.dat-50 & 75 & 2 & 12060 & 12060 & 13522 &  & 12060 & 0.00 \\
OR5x100-0.75\_7.dat-50 & 75 & 2 & 13344 & 13344 & 14571 &  & 13344 & 0.00 \\
OR5x100-0.75\_8.dat-50 & 75 & 2 & 13364 & 13364 & 15170 &  & 13364 & 0.00 \\
OR5x100-0.75\_9.dat-50 & 75 & 2 & 12077 & 12077 & 13567 &  & 12077 & 0.00 \\
OR5x100-0.75\_10.dat-50 & 75 & 2 & 12824 & 12824 & 14457 &  & 12824 & 0.00 \\
OR5x100-0.75\_1.dat-100 & 75 & 1 & 13104 & 13104 & 14018 & 11988 & 13104 & 0.00 \\
OR5x100-0.75\_2.dat-100 & 75 & 1 & 14603 & 14603 & 16122 & 13373 & 14603 & 0.00 \\
OR5x100-0.75\_3.dat-100 & 75 & 1 & 13214 & 13214 & 14408 & 12092 & 13214 & 0.00 \\
OR5x100-0.75\_4.dat-100 & 75 & 1 & 13555 & 13555 & 15053 & 12414 & 13555 & 0.00 \\
OR5x100-0.75\_5.dat-100 & 75 & 1 & 13604 & 13604 & 15116 & 12380 & 13604 & 0.00 \\
OR5x100-0.75\_6.dat-100 & 75 & 1 & 12929 & 12929 & 14960 & 11795 & 12929 & 0.00 \\
OR5x100-0.75\_7.dat-100 & 75 & 1 & 14103 & 14103 & 15902 & 12877 & 14103 & 0.00 \\
OR5x100-0.75\_8.dat-100 & 75 & 1 & 14004 & 14004 & 15169 & 12822 & 14004 & 0.00 \\
OR5x100-0.75\_9.dat-100 & 75 & 1 & 12783 & 12783 & 14027 & 11708 & 12783 & 0.00 \\
OR5x100-0.75\_10.dat-100 & 75 & 1 & 13322 & 13322 & 14702 & 12137 & 13322 & 0.00 \\
OR5x250-0.25\_1.dat-0 & 25 & 4 & 34494 & 51621 & 54798 &  & 50291 & 31.41 \\
OR5x250-0.25\_2.dat-0 & 25 & 4 & 36022 & 53453 & 56258 &  &  & 31.34 \\
OR5x250-0.25\_3.dat-0 & 25 & 4 & 35700 & 53506 & 56308 &  &  & 31.79 \\
OR5x250-0.25\_4.dat-0 & 25 & 4 & 35239 & 52573 & 55907 &  &  & 32.97 \\
OR5x250-0.25\_5.dat-0 & 25 & 4 & 33896 & 51603 & 54740 &  &  & 31.28 \\
OR5x250-0.25\_6.dat-0 & 25 & 4 & 34604 & 51849 & 55007 &  &  & 33.26 \\
OR5x250-0.25\_7.dat-0 & 25 & 4 & 35740 & 54014 & 55294 &  & 52020 & 31.30 \\
OR5x250-0.25\_8.dat-0 & 25 & 4 & 36632 & 55028 & 57526 &  & 53918 & 32.06 \\
OR5x250-0.25\_9.dat-0 & 25 & 4 & 35282 & 52910 & 55996 &  &  & 32.14 \\
OR5x250-0.25\_10.dat-0 & 25 & 4 & 35280 & 52031 & 55196 &  &  & 30.58 \\
OR5x250-0.25\_1.dat-50 & 25 & 2 & 38082 & 53836 & 55192 &  & 52384 & 27.30 \\
OR5x250-0.25\_2.dat-50 & 25 & 2 & 39285 & 55748 & 57851 &  & 54668 & 28.14 \\
OR5x250-0.25\_3.dat-50 & 25 & 2 & 38444 & 55392 & 57004 &  & 53485 & 28.12 \\
OR5x250-0.25\_4.dat-50 & 25 & 2 & 38716 & 56262 & 57428 &  & 54753 & 29.29 \\
OR5x250-0.25\_5.dat-50 & 25 & 2 & 38560 & 55651 & 55884 &  & 53890 & 28.45 \\
OR5x250-0.25\_6.dat-50 & 25 & 2 & 40088 & 56250 & 58016 &  & 54875 & 26.95 \\
OR5x250-0.25\_7.dat-50 & 25 & 2 & 40961 & 57130 & 59607 &  & 56511 & 27.52 \\
OR5x250-0.25\_8.dat-50 & 25 & 2 & 39968 & 58725 & 60600 &  & 57002 & 29.88 \\
OR5x250-0.25\_9.dat-50 & 25 & 2 & 38919 & 56112 & 57882 &  & 55402 & 29.75 \\
OR5x250-0.25\_10.dat-50 & 25 & 2 & 37815 & 53686 & 55473 &  & 52645 & 28.17 \\
OR5x250-0.25\_1.dat-100 & 25 & 1 & 46075 & 60613 & 61601 & 56913 & 58410 & 2.56 \\
OR5x250-0.25\_2.dat-100 & 25 & 1 & 45538 & 59501 & 60499 & 56239 & 57402 & 2.03 \\
OR5x250-0.25\_3.dat-100 & 25 & 1 & 43947 & 58421 & 59396 & 54982 & 56343 & 2.42 \\
OR5x250-0.25\_4.dat-100 & 25 & 1 & 44916 & 59212 & 59212 & 55270 & 56933 & 2.92 \\
OR5x250-0.25\_5.dat-100 & 25 & 1 & 45370 & 59118 & 59118 & 55404 & 56943 & 2.70 \\
OR5x250-0.25\_6.dat-100 & 25 & 1 & 46901 & 62072 & 63066 & 57987 & 59284 & 2.19 \\
OR5x250-0.25\_7.dat-100 & 25 & 1 & 48069 & 63569 & 63569 & 59937 & 61611 & 2.72 \\
OR5x250-0.25\_8.dat-100 & 25 & 1 & 45860 & 60675 & 60675 & 57505 & 58769 & 2.15 \\
OR5x250-0.25\_9.dat-100 & 25 & 1 & 46112 & 60663 & 60911 & 57097 & 58528 & 2.44 \\
OR5x250-0.25\_10.dat-100 & 25 & 1 & 47795 & 61775 & 61847 & 57576 & 59133 & 2.63 \\
OR5x250-0.50\_1.dat-0 & 50 & 4 & 40069 & 64651 & 84711 &  & 57509 & 30.33 \\
OR5x250-0.50\_2.dat-0 & 50 & 4 & 38908 & 62186 & 85349 &  & 54525 & 28.64 \\
OR5x250-0.50\_3.dat-0 & 50 & 4 & 37954 & 59247 & 76226 &  & 52500 & 27.71 \\
OR5x250-0.50\_4.dat-0 & 50 & 4 & 40005 & 65156 & 84713 &  & 59307 & 32.55 \\
OR5x250-0.50\_5.dat-0 & 50 & 4 & 40847 & 64714 & 88188 &  & 58188 & 29.80 \\
OR5x250-0.50\_6.dat-0 & 50 & 4 & 40393 & 63384 & 84965 &  & 55904 & 27.75 \\
OR5x250-0.50\_7.dat-0 & 50 & 4 & 39146 & 62593 & 87174 &  & 55731 & 29.76 \\
OR5x250-0.50\_8.dat-0 & 50 & 4 & 39927 & 62174 & 88969 &  & 56029 & 28.74 \\
OR5x250-0.50\_9.dat-0 & 50 & 4 & 39676 & 63075 & 87710 &  & 56625 & 29.93 \\
OR5x250-0.50\_10.dat-0 & 50 & 4 & 38309 & 60764 & 90119 &  & 54001 & 29.06 \\
OR5x250-0.50\_1.dat-50 & 50 & 2 & 47812 & 81492 & 83288 &  & 73732 & 35.15 \\
OR5x250-0.50\_2.dat-50 & 50 & 2 & 51097 & 79739 & 81981 &  & 70644 & 27.67 \\
OR5x250-0.50\_3.dat-50 & 50 & 2 & 46295 & 79097 & 83230 &  & 72948 & 36.54 \\
OR5x250-0.50\_4.dat-50 & 50 & 2 & 48764 & 80365 & 82044 &  & 72591 & 32.82 \\
OR5x250-0.50\_5.dat-50 & 50 & 2 & 48263 & 81400 & 87609 &  & 75882 & 36.40 \\
OR5x250-0.50\_6.dat-50 & 50 & 2 & 48037 & 80543 & 83800 &  & 74496 & 35.52 \\
OR5x250-0.50\_7.dat-50 & 50 & 2 & 45898 & 79433 & 84028 &  & 71471 & 35.78 \\
OR5x250-0.50\_8.dat-50 & 50 & 2 & 47169 & 81476 & 85274 &  & 73949 & 36.21 \\
OR5x250-0.50\_9.dat-50 & 50 & 2 & 47712 & 80901 & 83871 &  & 74801 & 36.21 \\
OR5x250-0.50\_10.dat-50 & 50 & 2 & 46116 & 77653 & 83708 &  & 71580 & 35.57 \\
OR5x250-0.50\_1.dat-100 & 50 & 1 & 51511 & 78004 & 78004 & 69671 & 73838 & 5.64 \\
OR5x250-0.50\_2.dat-100 & 50 & 1 & 51277 & 74726 & 74792 & 66618 & 70903 & 6.04 \\
OR5x250-0.50\_3.dat-100 & 50 & 1 & 49591 & 77700 & 78199 & 66098 & 73113 & 9.59 \\
OR5x250-0.50\_4.dat-100 & 50 & 1 & 52652 & 78077 & 78738 & 67268 & 73498 & 8.48 \\
OR5x250-0.50\_5.dat-100 & 50 & 1 & 55375 & 79729 & 80304 & 70969 & 75888 & 6.48 \\
OR5x250-0.50\_6.dat-100 & 50 & 1 & 52942 & 78894 & 78894 & 68667 & 75355 & 8.88 \\
OR5x250-0.50\_7.dat-100 & 50 & 1 & 53948 & 74755 & 76937 & 67519 & 72397 & 6.74 \\
OR5x250-0.50\_8.dat-100 & 50 & 1 & 51525 & 78290 & 78290 & 64771 & 73554 & 11.94 \\
OR5x250-0.50\_9.dat-100 & 50 & 1 & 53259 & 78999 & 78999 & 68884 & 74135 & 7.08 \\
OR5x250-0.50\_10.dat-100 & 50 & 1 & 52454 & 74272 & 75944 & 65988 & 72064 & 8.43 \\
OR5x250-0.75\_1.dat-0 & 75 & 4 & 24032 & 24032 & 26953 &  & 24032 & 0.00 \\
OR5x250-0.75\_2.dat-0 & 75 & 4 & 21799 & 21799 & 25425 &  & 21799 & 0.00 \\
OR5x250-0.75\_3.dat-0 & 75 & 4 & 25330 & 25330 & 31757 &  & 25330 & 0.00 \\
OR5x250-0.75\_4.dat-0 & 75 & 4 & 23637 & 23637 & 28816 &  & 23637 & 0.00 \\
OR5x250-0.75\_5.dat-0 & 75 & 4 & 24011 & 24011 & 32918 &  & 24011 & 0.00 \\
OR5x250-0.75\_6.dat-0 & 75 & 4 & 19766 & 19766 & 22525 &  & 19766 & 0.00 \\
OR5x250-0.75\_7.dat-0 & 75 & 4 & 21914 & 21914 & 25534 &  & 21914 & 0.00 \\
OR5x250-0.75\_8.dat-0 & 75 & 4 & 23874 & 23874 & 30417 &  & 23874 & 0.00 \\
OR5x250-0.75\_9.dat-0 & 75 & 4 & 25048 & 25048 & 30699 &  & 25048 & 0.00 \\
OR5x250-0.75\_10.dat-0 & 75 & 4 & 23965 & 23965 & 30455 &  & 23965 & 0.00 \\
OR5x250-0.75\_1.dat-50 & 75 & 2 & 31006 & 31006 & 37195 &  & 31006 & 0.00 \\
OR5x250-0.75\_2.dat-50 & 75 & 2 & 31726 & 31726 & 36581 &  & 31726 & 0.00 \\
OR5x250-0.75\_3.dat-50 & 75 & 2 & 32286 & 32286 & 39738 &  & 32286 & 0.00 \\
OR5x250-0.75\_4.dat-50 & 75 & 2 & 31989 & 31989 & 38905 &  & 31989 & 0.00 \\
OR5x250-0.75\_5.dat-50 & 75 & 2 & 31009 & 31009 & 38480 &  & 31009 & 0.00 \\
OR5x250-0.75\_6.dat-50 & 75 & 2 & 29841 & 29841 & 35169 &  & 29841 & 0.00 \\
OR5x250-0.75\_7.dat-50 & 75 & 2 & 31421 & 31421 & 37858 &  & 31421 & 0.00 \\
OR5x250-0.75\_8.dat-50 & 75 & 2 & 30763 & 30763 & 36473 &  & 30763 & 0.00 \\
OR5x250-0.75\_9.dat-50 & 75 & 2 & 32198 & 32198 & 37947 &  & 32198 & 0.00 \\
OR5x250-0.75\_10.dat-50 & 75 & 2 & 32569 & 32569 & 36519 &  & 32569 & 0.00 \\
OR5x250-0.75\_1.dat-100 & 75 & 1 & 33318 & 33318 & 40373 & 32031 & 33358 & 0.00 \\
OR5x250-0.75\_2.dat-100 & 75 & 1 & 33912 & 33912 & 38176 & 32705 & 33912 & 0.00 \\
OR5x250-0.75\_3.dat-100 & 75 & 1 & 34439 & 34439 & 41500 & 33198 & 34471 & 0.00 \\
OR5x250-0.75\_4.dat-100 & 75 & 1 & 34306 & 34306 & 39445 & 32946 & 34339 & 0.00 \\
OR5x250-0.75\_5.dat-100 & 75 & 1 & 33649 & 33649 & 38720 & 32400 & 33665 & 0.00 \\
OR5x250-0.75\_6.dat-100 & 75 & 1 & 31670 & 31670 & 36294 & 30306 & 31670 & 0.00 \\
OR5x250-0.75\_7.dat-100 & 75 & 1 & 33373 & 33373 & 38127 & 32122 & 33373 & 0.00 \\
OR5x250-0.75\_8.dat-100 & 75 & 1 & 32417 & 32417 & 37313 & 31260 & 32417 & 0.00 \\
OR5x250-0.75\_9.dat-100 & 75 & 1 & 34576 & 34576 & 38511 & 33186 & 34607 & 0.00 \\
OR5x250-0.75\_10.dat-100 & 75 & 1 & 33891 & 33891 & 39129 & 32657 & 33928 & 0.00 \\
OR5x500-0.25\_1.dat-0 & 25 & 4 & 69835 & 110214 & 117265 &  &  & 34.49 \\
OR5x500-0.25\_2.dat-0 & 25 & 4 & 67484 & 107643 & 111973 &  &  & 36.71 \\
OR5x500-0.25\_3.dat-0 & 25 & 4 & 70067 & 111192 & 118046 &  &  & 36.52 \\
OR5x500-0.25\_4.dat-0 & 25 & 4 & 69097 & 110303 & 120882 &  &  & 35.98 \\
OR5x500-0.25\_5.dat-0 & 25 & 4 & 68932 & 110666 & 117279 &  &  & 36.44 \\
OR5x500-0.25\_6.dat-0 & 25 & 4 & 67351 & 108756 &  &  &  & 36.14 \\
OR5x500-0.25\_7.dat-0 & 25 & 4 & 68077 & 108402 &  &  &  & 35.11 \\
OR5x500-0.25\_8.dat-0 & 25 & 4 & 67539 & 108202 & 112423 &  &  & 34.93 \\
OR5x500-0.25\_9.dat-0 & 25 & 4 & 69004 & 110234 & 113726 &  &  & 35.77 \\
OR5x500-0.25\_10.dat-0 & 25 & 4 & 68019 & 108009 &  &  &  & 35.73 \\
OR5x500-0.25\_1.dat-50 & 25 & 2 & 76052 & 113261 & 114719 &  & 111010 & 31.49 \\
OR5x500-0.25\_2.dat-50 & 25 & 2 & 75218 & 111681 & 114348 &  & 112174 & 32.65 \\
OR5x500-0.25\_3.dat-50 & 25 & 2 & 78945 & 118460 & 119657 &  & 114278 & 30.92 \\
OR5x500-0.25\_4.dat-50 & 25 & 2 & 79508 & 116473 & 116473 &  & 111381 & 28.62 \\
OR5x500-0.25\_5.dat-50 & 25 & 2 & 75546 & 114868 & 116012 &  & 110268 & 31.49 \\
OR5x500-0.25\_6.dat-50 & 25 & 2 & 77490 & 114252 & 116134 &  & 111698 & 30.63 \\
OR5x500-0.25\_7.dat-50 & 25 & 2 & 77511 & 112684 & 113888 &  & 109462 & 29.19 \\
OR5x500-0.25\_8.dat-50 & 25 & 2 & 76353 & 113134 & 113597 &  & 111252 & 31.37 \\
OR5x500-0.25\_9.dat-50 & 25 & 2 & 77845 & 115734 & 115734 &  & 113693 & 31.53 \\
OR5x500-0.25\_10.dat-50 & 25 & 2 & 76071 & 113094 & 114455 &  & 109429 & 30.48 \\
OR5x500-0.25\_1.dat-100 & 25 & 1 & 90594 & 123807 & 123813 & 112053 & 117212 & 4.40 \\
OR5x500-0.25\_2.dat-100 & 25 & 1 & 90110 & 120379 & 120379 & 110468 & 115579 & 4.42 \\
OR5x500-0.25\_3.dat-100 & 25 & 1 & 90361 & 122897 & 122897 & 111471 & 116685 & 4.47 \\
OR5x500-0.25\_4.dat-100 & 25 & 1 & 93316 & 126374 & 126374 & 113621 & 119915 & 5.25 \\
OR5x500-0.25\_5.dat-100 & 25 & 1 & 89472 & 120760 & 120760 & 110348 & 115732 & 4.65 \\
OR5x500-0.25\_6.dat-100 & 25 & 1 & 90508 & 117619 & 117619 & 108350 & 113962 & 4.92 \\
OR5x500-0.25\_7.dat-100 & 25 & 1 & 91222 & 118599 & 118599 & 109038 & 114264 & 4.57 \\
OR5x500-0.25\_8.dat-100 & 25 & 1 & 90045 & 117971 & 117971 & 106979 & 112554 & 4.95 \\
OR5x500-0.25\_9.dat-100 & 25 & 1 & 93289 & 123411 & 123411 & 113207 & 118639 & 4.58 \\
OR5x500-0.25\_10.dat-100 & 25 & 1 & 91103 & 120024 & 120024 & 109797 & 115683 & 5.09 \\
OR5x500-0.50\_1.dat-0 & 50 & 4 & 77109 & 123841 & 214120 &  & 108216 & 28.75 \\
OR5x500-0.50\_2.dat-0 & 50 & 4 & 74534 & 124608 &  &  & 109574 & 31.98 \\
OR5x500-0.50\_3.dat-0 & 50 & 4 & 71925 & 123468 & 211434 &  & 108233 & 33.55 \\
OR5x500-0.50\_4.dat-0 & 50 & 4 & 76000 & 129661 &  &  & 113945 & 33.30 \\
OR5x500-0.50\_5.dat-0 & 50 & 4 & 74417 & 126925 & 201516 &  & 110730 & 32.79 \\
OR5x500-0.50\_6.dat-0 & 50 & 4 & 76099 & 125246 & 197735 &  & 111703 & 31.87 \\
OR5x500-0.50\_7.dat-0 & 50 & 4 & 74173 & 127351 & 204545 &  & 110612 & 32.94 \\
OR5x500-0.50\_8.dat-0 & 50 & 4 & 74840 & 127420 & 218814 &  & 111627 & 32.96 \\
OR5x500-0.50\_9.dat-0 & 50 & 4 & 72145 & 121479 &  &  & 106432 & 32.21 \\
OR5x500-0.50\_10.dat-0 & 50 & 4 & 76215 & 129471 & 196699 &  & 115998 & 34.30 \\
OR5x500-0.50\_1.dat-50 & 50 & 2 & 93703 & 161392 & 164286 &  & 144300 & 35.06 \\
OR5x500-0.50\_2.dat-50 & 50 & 2 & 94588 & 163727 & 163727 &  & 147077 & 35.69 \\
OR5x500-0.50\_3.dat-50 & 50 & 2 & 91607 & 159366 & 163792 &  & 144492 & 36.60 \\
OR5x500-0.50\_4.dat-50 & 50 & 2 & 94797 & 167239 & 168310 &  & 150448 & 36.99 \\
OR5x500-0.50\_5.dat-50 & 50 & 2 & 92434 & 162645 & 165910 &  & 144796 & 36.16 \\
OR5x500-0.50\_6.dat-50 & 50 & 2 & 91819 & 160755 & 167066 &  & 147365 & 37.69 \\
OR5x500-0.50\_7.dat-50 & 50 & 2 & 91349 & 159842 & 168265 &  & 145558 & 37.24 \\
OR5x500-0.50\_8.dat-50 & 50 & 2 & 93238 & 161161 & 161684 &  & 146795 & 36.48 \\
OR5x500-0.50\_9.dat-50 & 50 & 2 & 91421 & 159578 & 166726 &  & 144903 & 36.91 \\
OR5x500-0.50\_10.dat-50 & 50 & 2 & 92710 & 164190 & 169033 &  & 147553 & 37.17 \\
OR5x500-0.50\_1.dat-100 & 50 & 1 & 95743 & 152050 & 152050 & 122445 & 143472 & 14.66 \\
OR5x500-0.50\_2.dat-100 & 50 & 1 & 96224 & 152674 & 152674 & 123075 & 145663 & 15.51 \\
OR5x500-0.50\_3.dat-100 & 50 & 1 & 95193 & 150648 & 150648 & 122060 & 142500 & 14.34 \\
OR5x500-0.50\_4.dat-100 & 50 & 1 & 98696 & 158065 & 158065 & 123483 & 148736 & 16.98 \\
OR5x500-0.50\_5.dat-100 & 50 & 1 & 95783 & 154977 & 154977 & 118933 & 146297 & 18.70 \\
OR5x500-0.50\_6.dat-100 & 50 & 1 & 96634 & 154727 & 154727 & 121111 & 146071 & 17.09 \\
OR5x500-0.50\_7.dat-100 & 50 & 1 & 95596 & 156375 & 156375 & 119143 & 145735 & 18.25 \\
OR5x500-0.50\_8.dat-100 & 50 & 1 & 96324 & 151145 & 151145 & 121654 & 144375 & 15.74 \\
OR5x500-0.50\_9.dat-100 & 50 & 1 & 94932 & 153060 & 153060 & 115963 & 144891 & 19.97 \\
OR5x500-0.50\_10.dat-100 & 50 & 1 & 95786 & 152964 & 155612 & 116185 & 147142 & 21.04 \\
OR5x500-0.75\_1.dat-0 & 75 & 4 & 33280 & 70213 & 275421 &  & 42843 & 22.32 \\
OR5x500-0.75\_2.dat-0 & 75 & 4 & 48611 & 66550 & 83975 &  & 48617 & 0.01 \\
OR5x500-0.75\_3.dat-0 & 75 & 4 & 35724 & 71632 & 97517 &  & 47256 & 24.40 \\
OR5x500-0.75\_4.dat-0 & 75 & 4 & 35106 & 71978 & 100542 &  & 48597 & 27.76 \\
OR5x500-0.75\_5.dat-0 & 75 & 4 & 35165 & 69985 & 96801 &  & 48644 & 27.71 \\
OR5x500-0.75\_6.dat-0 & 75 & 4 & 36713 & 65198 & 92928 &  & 46868 & 21.67 \\
OR5x500-0.75\_7.dat-0 & 75 & 4 & 45651 & 66476 & 94972 &  & 45668 & 0.04 \\
OR5x500-0.75\_8.dat-0 & 75 & 4 & 34991 & 66585 & 82364 &  & 46099 & 24.10 \\
OR5x500-0.75\_9.dat-0 & 75 & 4 & 37192 & 46817 & 100248 &  & 46817 & 20.56 \\
OR5x500-0.75\_10.dat-0 & 75 & 4 & 45895 & 65644 & 83529 &  & 45921 & 0.06 \\
OR5x500-0.75\_1.dat-50 & 75 & 2 & 44989 & 61770 & 97526 &  & 61697 & 27.08 \\
OR5x500-0.75\_2.dat-50 & 75 & 2 & 59252 & 62479 & 100543 &  & 62391 & 5.03 \\
OR5x500-0.75\_3.dat-50 & 75 & 2 & 59953 & 61213 & 91905 &  & 61034 & 1.77 \\
OR5x500-0.75\_4.dat-50 & 75 & 2 & 46017 & 63661 & 90979 &  & 63550 & 27.59 \\
OR5x500-0.75\_5.dat-50 & 75 & 2 & 46155 & 62096 & 96309 &  & 62048 & 25.61 \\
OR5x500-0.75\_6.dat-50 & 75 & 2 & 46315 & 63547 & 91224 &  & 63437 & 26.99 \\
OR5x500-0.75\_7.dat-50 & 75 & 2 & 61642 & 61745 & 99090 &  & 61693 & 0.08 \\
OR5x500-0.75\_8.dat-50 & 75 & 2 & 63201 & 63342 & 108200 &  & 63248 & 0.07 \\
OR5x500-0.75\_9.dat-50 & 75 & 2 & 45302 & 62819 & 101695 &  & 62717 & 27.77 \\
OR5x500-0.75\_10.dat-50 & 75 & 2 & 44977 & 61063 & 102403 &  & 60990 & 26.26 \\
OR5x500-0.75\_1.dat-100 & 75 & 1 & 53711 & 64823 & 91492 & 63493 & 64837 & 2.05 \\
OR5x500-0.75\_2.dat-100 & 75 & 1 & 64903 & 68349 & 96443 & 66939 & 68329 & 2.03 \\
OR5x500-0.75\_3.dat-100 & 75 & 1 & 62448 & 65458 & 90500 & 64125 & 65432 & 2.00 \\
OR5x500-0.75\_4.dat-100 & 75 & 1 & 67838 & 68073 & 95257 & 66760 & 68060 & 0.33 \\
OR5x500-0.75\_5.dat-100 & 75 & 1 & 62361 & 65516 & 93947 & 64190 & 65544 & 2.02 \\
OR5x500-0.75\_6.dat-100 & 75 & 1 & 63429 & 66482 & 93609 & 65119 & 66459 & 2.02 \\
OR5x500-0.75\_7.dat-100 & 75 & 1 & 64910 & 65076 & 97773 & 63786 & 65065 & 0.24 \\
OR5x500-0.75\_8.dat-100 & 75 & 1 & 68064 & 68254 & 96401 & 66870 & 68282 & 0.28 \\
OR5x500-0.75\_9.dat-100 & 75 & 1 & 63244 & 66562 & 91808 & 65278 & 66541 & 1.90 \\
OR5x500-0.75\_10.dat-100 & 75 & 1 & 62937 & 66051 & 92698 & 64709 & 66070 & 2.03 \\
OR10x100-0.25\_1.dat-0 & 25 & 9 & 14537 & 19029 & 21555 &  & 18761 & 22.51 \\
OR10x100-0.25\_2.dat-0 & 25 & 9 & 14481 & 18822 & 20668 &  &  & 23.06 \\
OR10x100-0.25\_3.dat-0 & 25 & 9 & 14432 & 18681 & 20698 &  & 18323 & 21.24 \\
OR10x100-0.25\_4.dat-0 & 25 & 9 & 14720 & 19808 & 21661 &  & 19703 & 25.29 \\
OR10x100-0.25\_5.dat-0 & 25 & 9 & 14857 & 19281 & 21437 &  &  & 22.94 \\
OR10x100-0.25\_6.dat-0 & 25 & 9 & 14875 & 19330 & 21616 &  & 19372 & 23.05 \\
OR10x100-0.25\_7.dat-0 & 25 & 9 & 13644 & 18104 & 20751 &  &  & 24.64 \\
OR10x100-0.25\_8.dat-0 & 25 & 9 & 14546 & 18738 & 21025 &  & 18854 & 22.37 \\
OR10x100-0.25\_9.dat-0 & 25 & 9 & 13989 & 18348 & 19917 &  & 18154 & 22.94 \\
OR10x100-0.25\_10.dat-0 & 25 & 9 & 14119 & 18673 & 20543 &  & 18565 & 23.95 \\
OR10x100-0.25\_1.dat-50 & 25 & 5 & 15584 & 20784 & 23483 &  & 20673 & 24.62 \\
OR10x100-0.25\_2.dat-50 & 25 & 5 & 15054 & 20134 & 21586 &  & 20000 & 24.73 \\
OR10x100-0.25\_3.dat-50 & 25 & 5 & 14995 & 20097 & 21085 &  & 19736 & 24.02 \\
OR10x100-0.25\_4.dat-50 & 25 & 5 & 15804 & 20754 & 22771 &  &  & 23.85 \\
OR10x100-0.25\_5.dat-50 & 25 & 5 & 15826 & 20765 & 22055 &  & 20480 & 22.72 \\
OR10x100-0.25\_6.dat-50 & 25 & 5 & 16183 & 21415 & 22920 &  & 21121 & 23.38 \\
OR10x100-0.25\_7.dat-50 & 25 & 5 & 14361 & 19438 & 20736 &  & 19009 & 24.45 \\
OR10x100-0.25\_8.dat-50 & 25 & 5 & 14996 & 20322 & 21458 &  & 19872 & 24.54 \\
OR10x100-0.25\_9.dat-50 & 25 & 5 & 14717 & 19527 & 21027 &  & 19106 & 22.97 \\
OR10x100-0.25\_10.dat-50 & 25 & 5 & 14825 & 20264 & 22328 &  & 19924 & 25.59 \\
OR10x100-0.25\_1.dat-100 & 25 & 1 & 21158 & 25534 & 26730 & 24632 & 24994 & 1.45 \\
OR10x100-0.25\_2.dat-100 & 25 & 1 & 21782 & 25551 & 26754 & 24132 & 24767 & 2.56 \\
OR10x100-0.25\_3.dat-100 & 25 & 1 & 20762 & 25175 & 27090 & 23863 & 24515 & 2.66 \\
OR10x100-0.25\_4.dat-100 & 25 & 1 & 20157 & 23669 & 24052 & 23011 & 23239 & 0.98 \\
OR10x100-0.25\_5.dat-100 & 25 & 1 & 20096 & 23151 & 24424 & 23016 & 23417 & 0.58 \\
OR10x100-0.25\_6.dat-100 & 25 & 1 & 21247 & 26003 & 27184 & 25036 & 25612 & 2.25 \\
OR10x100-0.25\_7.dat-100 & 25 & 1 & 20146 & 24079 & 25884 & 23051 & 23787 & 3.09 \\
OR10x100-0.25\_8.dat-100 & 25 & 1 & 21219 & 25907 & 26857 & 24635 & 25403 & 3.02 \\
OR10x100-0.25\_9.dat-100 & 25 & 1 & 19077 & 22758 & 23084 & 22091 & 22525 & 1.93 \\
OR10x100-0.25\_10.dat-100 & 25 & 1 & 20546 & 25246 & 26228 & 24062 & 24599 & 2.18 \\
OR10x100-0.50\_1.dat-0 & 50 & 9 & 15216 & 22209 & 26516 &  & 20476 & 25.69 \\
OR10x100-0.50\_2.dat-0 & 50 & 9 & 15426 & 21266 & 25742 &  & 20545 & 24.92 \\
OR10x100-0.50\_3.dat-0 & 50 & 9 & 19112 & 22787 & 34196 &  & 20916 & 8.62 \\
OR10x100-0.50\_4.dat-0 & 50 & 9 & 17568 & 26422 & 32127 &  & 24455 & 28.16 \\
OR10x100-0.50\_5.dat-0 & 50 & 9 & 15358 & 22399 & 29498 &  & 20730 & 25.91 \\
OR10x100-0.50\_6.dat-0 & 50 & 9 & 15440 & 23188 & 31185 &  & 21633 & 28.63 \\
OR10x100-0.50\_7.dat-0 & 50 & 9 & 18533 & 24346 & 30828 &  & 22981 & 19.36 \\
OR10x100-0.50\_8.dat-0 & 50 & 9 & 18216 & 23791 & 33462 &  & 21906 & 16.84 \\
OR10x100-0.50\_9.dat-0 & 50 & 9 & 16056 & 23364 & 28760 &  & 21609 & 25.70 \\
OR10x100-0.50\_10.dat-0 & 50 & 9 & 15718 & 21319 & 25989 &  & 19557 & 19.63 \\
OR10x100-0.50\_1.dat-50 & 50 & 5 & 18370 & 28595 & 31221 &  & 25614 & 28.28 \\
OR10x100-0.50\_2.dat-50 & 50 & 5 & 20276 & 29586 & 33445 &  & 27913 & 27.36 \\
OR10x100-0.50\_3.dat-50 & 50 & 5 & 19767 & 31145 & 33428 &  & 28340 & 30.25 \\
OR10x100-0.50\_4.dat-50 & 50 & 5 & 21026 & 32837 & 35267 &  & 29808 & 29.46 \\
OR10x100-0.50\_5.dat-50 & 50 & 5 & 18568 & 29650 & 32485 &  & 27197 & 31.73 \\
OR10x100-0.50\_6.dat-50 & 50 & 5 & 19411 & 30498 & 33486 &  & 27518 & 29.46 \\
OR10x100-0.50\_7.dat-50 & 50 & 5 & 19257 & 29791 & 33101 &  & 27821 & 30.78 \\
OR10x100-0.50\_8.dat-50 & 50 & 5 & 18556 & 30003 & 32252 &  & 26704 & 30.51 \\
OR10x100-0.50\_9.dat-50 & 50 & 5 & 18731 & 30808 & 31770 &  & 27713 & 32.41 \\
OR10x100-0.50\_10.dat-50 & 50 & 5 & 18813 & 29318 & 30758 &  & 26340 & 28.58 \\
OR10x100-0.50\_1.dat-100 & 50 & 1 & 19018 & 29329 & 29722 & 26565 & 27500 & 3.40 \\
OR10x100-0.50\_2.dat-100 & 50 & 1 & 20461 & 29499 & 30910 & 27630 & 28430 & 2.81 \\
OR10x100-0.50\_3.dat-100 & 50 & 1 & 21243 & 30415 & 32149 & 28023 & 29356 & 4.54 \\
OR10x100-0.50\_4.dat-100 & 50 & 1 & 23916 & 30985 & 31201 & 28913 & 29876 & 3.22 \\
OR10x100-0.50\_5.dat-100 & 50 & 1 & 20743 & 28407 & 29681 & 26550 & 27632 & 3.92 \\
OR10x100-0.50\_6.dat-100 & 50 & 1 & 20724 & 29761 & 30654 & 27684 & 28708 & 3.57 \\
OR10x100-0.50\_7.dat-100 & 50 & 1 & 21211 & 30595 & 30595 & 28508 & 29565 & 3.58 \\
OR10x100-0.50\_8.dat-100 & 50 & 1 & 20788 & 30178 & 30901 & 27463 & 28648 & 4.14 \\
OR10x100-0.50\_9.dat-100 & 50 & 1 & 19814 & 30015 & 30340 & 27070 & 28147 & 3.83 \\
OR10x100-0.50\_10.dat-100 & 50 & 1 & 18980 & 27905 & 27905 & 25774 & 26650 & 3.29 \\
OR10x100-0.75\_1.dat-0 & 75 & 9 & 9389 & 10212 & 10334 &  & 9426 & 0.39 \\
OR10x100-0.75\_2.dat-0 & 75 & 9 & 9057 & 9057 & 9436 &  & 9057 & 0.00 \\
OR10x100-0.75\_3.dat-0 & 75 & 9 & 8747 & 9310 & 9462 &  & 8749 & 0.02 \\
OR10x100-0.75\_4.dat-0$^{\dagger}$ & 75 & 9 & 8814 & 8813 & 9447 &  & 8813 & 0.00 \\
OR10x100-0.75\_5.dat-0 & 75 & 9 & 8616 & 9027 & 9850 &  & 8685 & 0.79 \\
OR10x100-0.75\_6.dat-0 & 75 & 9 & 9638 & 9638 & 10426 &  & 9638 & 0.00 \\
OR10x100-0.75\_7.dat-0 & 75 & 9 & 8374 & 8374 & 9661 &  & 8374 & 0.00 \\
OR10x100-0.75\_8.dat-0 & 75 & 9 & 8745 & 8745 & 9631 &  & 8745 & 0.00 \\
OR10x100-0.75\_9.dat-0 & 75 & 9 & 8364 & 8364 & 8853 &  & 8364 & 0.00 \\
OR10x100-0.75\_10.dat-0 & 75 & 9 & 7632 & 8183 & 8950 &  & 7671 & 0.51 \\
OR10x100-0.75\_1.dat-50 & 75 & 5 & 12090 & 13507 & 14330 &  & 12339 & 2.02 \\
OR10x100-0.75\_2.dat-50 & 75 & 5 & 12747 & 13752 & 14781 &  & 12970 & 1.72 \\
OR10x100-0.75\_3.dat-50 & 75 & 5 & 12421 & 13957 & 13957 &  & 12597 & 1.40 \\
OR10x100-0.75\_4.dat-50 & 75 & 5 & 13358 & 14742 & 15110 &  & 13571 & 1.57 \\
OR10x100-0.75\_5.dat-50 & 75 & 5 & 12750 & 14444 & 14866 &  & 12954 & 1.57 \\
OR10x100-0.75\_6.dat-50 & 75 & 5 & 14202 & 15594 & 16039 &  & 14401 & 1.38 \\
OR10x100-0.75\_7.dat-50 & 75 & 5 & 12214 & 13308 & 14590 &  & 12375 & 1.30 \\
OR10x100-0.75\_8.dat-50 & 75 & 5 & 12851 & 14951 & 14993 &  & 12996 & 1.12 \\
OR10x100-0.75\_9.dat-50 & 75 & 5 & 12896 & 14548 & 15142 &  & 13085 & 1.44 \\
OR10x100-0.75\_10.dat-50 & 75 & 5 & 11983 & 13232 & 13232 &  & 12240 & 2.10 \\
OR10x100-0.75\_1.dat-100 & 75 & 1 & 13599 & 13599 & 14962 & 12582 & 13599 & 0.00 \\
OR10x100-0.75\_2.dat-100 & 75 & 1 & 14155 & 14155 & 15768 & 12895 & 14317 & 0.00 \\
OR10x100-0.75\_3.dat-100 & 75 & 1 & 14216 & 14216 & 15618 & 12996 & 14216 & 0.00 \\
OR10x100-0.75\_4.dat-100 & 75 & 1 & 15290 & 15290 & 16627 & 13869 & 15332 & 0.00 \\
OR10x100-0.75\_5.dat-100 & 75 & 1 & 14439 & 14439 & 15413 & 13360 & 14439 & 0.00 \\
OR10x100-0.75\_6.dat-100 & 75 & 1 & 15347 & 15347 & 16867 & 14256 & 15347 & 0.00 \\
OR10x100-0.75\_7.dat-100 & 75 & 1 & 13934 & 13934 & 14996 & 12688 & 13984 & 0.00 \\
OR10x100-0.75\_8.dat-100 & 75 & 1 & 14137 & 14137 & 15508 & 13006 & 14137 & 0.00 \\
OR10x100-0.75\_9.dat-100 & 75 & 1 & 14185 & 14185 & 15078 & 12968 & 14252 & 0.00 \\
OR10x100-0.75\_10.dat-100 & 75 & 1 & 13814 & 13814 & 15169 & 12728 & 13814 & 0.00 \\
OR10x250-0.25\_1.dat-0 & 25 & 9 & 29491 & 52379 &  &  &  & 40.38 \\
OR10x250-0.25\_2.dat-0 & 25 & 9 & 28790 & 51927 &  &  &  & 41.88 \\
OR10x250-0.25\_3.dat-0 & 25 & 9 & 31377 & 51223 &  &  &  & 34.80 \\
OR10x250-0.25\_4.dat-0 & 25 & 9 & 33546 & 53025 &  &  &  & 35.61 \\
OR10x250-0.25\_5.dat-0 & 25 & 9 &  &  &  &  &  & 100.00 \\
OR10x250-0.25\_6.dat-0 & 25 & 9 & 32947 & 52205 &  &  &  & 34.00 \\
OR10x250-0.25\_7.dat-0 & 25 & 9 & 32666 & 53579 &  &  &  & 37.05 \\
OR10x250-0.25\_8.dat-0 & 25 & 9 & 31028 & 51707 &  &  &  & 38.65 \\
OR10x250-0.25\_9.dat-0 & 25 & 9 & 31302 & 51439 &  &  &  & 36.94 \\
OR10x250-0.25\_10.dat-0 & 25 & 9 & 33079 & 52856 &  &  &  & 35.43 \\
OR10x250-0.25\_1.dat-50 & 25 & 5 & 33844 & 53502 & 54589 &  &  & 33.89 \\
OR10x250-0.25\_2.dat-50 & 25 & 5 & 34108 & 53382 & 56458 &  &  & 35.76 \\
OR10x250-0.25\_3.dat-50 & 25 & 5 & 33876 & 52024 & 54475 &  &  & 33.43 \\
OR10x250-0.25\_4.dat-50 & 25 & 5 & 35545 & 54623 & 57372 &  &  & 33.41 \\
OR10x250-0.25\_5.dat-50 & 25 & 5 & 34172 & 53530 & 56657 &  &  & 33.51 \\
OR10x250-0.25\_6.dat-50 & 25 & 5 & 35036 & 53835 & 57033 &  &  & 32.52 \\
OR10x250-0.25\_7.dat-50 & 25 & 5 & 35749 & 55317 & 58507 &  &  & 32.44 \\
OR10x250-0.25\_8.dat-50 & 25 & 5 & 34152 & 53363 & 55998 &  &  & 33.62 \\
OR10x250-0.25\_9.dat-50 & 25 & 5 & 33469 & 53427 & 56737 &  &  & 35.30 \\
OR10x250-0.25\_10.dat-50 & 25 & 5 & 34721 & 54167 & 58614 &  &  & 33.22 \\
OR10x250-0.25\_1.dat-100 & 25 & 1 & 45733 & 60741 & 60741 & 57006 & 58509 & 2.57 \\
OR10x250-0.25\_2.dat-100 & 25 & 1 & 45395 & 61073 & 62057 & 56593 & 58389 & 3.08 \\
OR10x250-0.25\_3.dat-100 & 25 & 1 & 47383 & 62082 & 63543 & 57819 & 60158 & 3.89 \\
OR10x250-0.25\_4.dat-100 & 25 & 1 & 45811 & 61629 & 61629 & 57008 & 58720 & 2.92 \\
OR10x250-0.25\_5.dat-100 & 25 & 1 & 47078 & 63531 & 64835 & 58258 & 60139 & 3.13 \\
OR10x250-0.25\_6.dat-100 & 25 & 1 & 46276 & 61727 & 62410 & 56557 & 58847 & 3.89 \\
OR10x250-0.25\_7.dat-100 & 25 & 1 & 47040 & 61136 & 61691 & 56755 & 58419 & 2.85 \\
OR10x250-0.25\_8.dat-100 & 25 & 1 & 45924 & 60648 & 61592 & 56036 & 57691 & 2.87 \\
OR10x250-0.25\_9.dat-100 & 25 & 1 & 47463 & 61503 & 61503 & 57741 & 59446 & 2.87 \\
OR10x250-0.25\_10.dat-100 & 25 & 1 & 46128 & 62297 & 63736 & 57528 & 59535 & 3.37 \\
OR10x250-0.50\_1.dat-0 & 50 & 9 &  &  &  &  & 57818 & 100.00 \\
OR10x250-0.50\_2.dat-0 & 50 & 9 & 38126 & 62818 &  &  & 54748 & 30.36 \\
OR10x250-0.50\_3.dat-0 & 50 & 9 & 37856 & 61853 &  &  & 54326 & 30.32 \\
OR10x250-0.50\_4.dat-0 & 50 & 9 & 37582 & 59368 &  &  & 52717 & 28.71 \\
OR10x250-0.50\_5.dat-0 & 50 & 9 & 38115 & 62987 &  &  & 58115 & 34.41 \\
OR10x250-0.50\_6.dat-0 & 50 & 9 & 39300 & 62837 &  &  & 56345 & 30.25 \\
OR10x250-0.50\_7.dat-0 & 50 & 9 & 37201 & 61372 &  &  & 53302 & 30.21 \\
OR10x250-0.50\_8.dat-0 & 50 & 9 & 36697 & 60043 &  &  & 52577 & 30.20 \\
OR10x250-0.50\_9.dat-0 & 50 & 9 & 36004 & 57088 &  &  & 49519 & 27.29 \\
OR10x250-0.50\_10.dat-0 & 50 & 9 & 35602 & 57134 &  &  & 49964 & 28.74 \\
OR10x250-0.50\_1.dat-50 & 50 & 5 & 46783 & 83614 & 90802 &  & 73119 & 36.02 \\
OR10x250-0.50\_2.dat-50 & 50 & 5 & 45093 & 79698 & 87145 &  & 71142 & 36.62 \\
OR10x250-0.50\_3.dat-50 & 50 & 5 & 45192 & 46801 & 87484 &  & 71263 & 3.44 \\
OR10x250-0.50\_4.dat-50 & 50 & 5 & 46428 & 83324 & 86634 &  & 72536 & 35.99 \\
OR10x250-0.50\_5.dat-50 & 50 & 5 & 45521 & 81995 & 85921 &  & 69837 & 34.82 \\
OR10x250-0.50\_6.dat-50 & 50 & 5 & 46130 & 78957 & 88255 &  & 71000 & 35.03 \\
OR10x250-0.50\_7.dat-50 & 50 & 5 & 44281 & 78565 & 84387 &  & 68581 & 35.43 \\
OR10x250-0.50\_8.dat-50 & 50 & 5 & 44289 & 79996 & 86485 &  & 68859 & 35.68 \\
OR10x250-0.50\_9.dat-50 & 50 & 5 & 45700 & 81071 & 90949 &  & 70068 & 34.78 \\
OR10x250-0.50\_10.dat-50 & 50 & 5 & 43973 & 77887 & 81253 &  & 67139 & 34.50 \\
OR10x250-0.50\_1.dat-100 & 50 & 1 & 54353 & 80190 & 80190 & 70254 & 75854 & 7.38 \\
OR10x250-0.50\_2.dat-100 & 50 & 1 & 50538 & 74852 & 74852 & 66529 & 70788 & 6.02 \\
OR10x250-0.50\_3.dat-100 & 50 & 1 & 47200 & 75120 & 75120 & 66152 & 70219 & 5.79 \\
OR10x250-0.50\_4.dat-100 & 50 & 1 & 53571 & 76382 & 76382 & 67648 & 72297 & 6.43 \\
OR10x250-0.50\_5.dat-100 & 50 & 1 & 48841 & 76790 & 76790 & 68015 & 72938 & 6.75 \\
OR10x250-0.50\_6.dat-100 & 50 & 1 & 47428 & 76052 & 76052 & 67528 & 71923 & 6.11 \\
OR10x250-0.50\_7.dat-100 & 50 & 1 & 48261 & 73743 & 73743 & 65056 & 69311 & 6.14 \\
OR10x250-0.50\_8.dat-100 & 50 & 1 & 49499 & 73301 & 73403 & 64555 & 68763 & 6.12 \\
OR10x250-0.50\_9.dat-100 & 50 & 1 & 52821 & 79259 & 79836 & 69811 & 74526 & 6.33 \\
OR10x250-0.50\_10.dat-100 & 50 & 1 & 46410 & 71374 & 71374 & 63747 & 67988 & 6.24 \\
OR10x250-0.75\_1.dat-0 & 75 & 9 & 22839 & 25765 &  &  & 22888 & 0.21 \\
OR10x250-0.75\_2.dat-0 & 75 & 9 & 21009 & 23182 &  &  & 21041 & 0.15 \\
OR10x250-0.75\_3.dat-0 & 75 & 9 & 22398 & 25879 &  &  & 22460 & 0.28 \\
OR10x250-0.75\_4.dat-0 & 75 & 9 & 23056 & 25999 &  &  & 23083 & 0.12 \\
OR10x250-0.75\_5.dat-0 & 75 & 9 & 23022 & 26190 &  &  & 23031 & 0.04 \\
OR10x250-0.75\_6.dat-0 & 75 & 9 & 23154 & 25506 &  &  & 23173 & 0.08 \\
OR10x250-0.75\_7.dat-0 & 75 & 9 & 23248 & 25986 &  &  & 23276 & 0.12 \\
OR10x250-0.75\_8.dat-0 & 75 & 9 & 22123 & 25242 &  &  & 22156 & 0.15 \\
OR10x250-0.75\_9.dat-0 & 75 & 9 & 22775 & 25768 & 143153 &  & 22794 & 0.08 \\
OR10x250-0.75\_10.dat-0 & 75 & 9 & 22613 & 25872 &  &  & 22631 & 0.08 \\
OR10x250-0.75\_1.dat-50 & 75 & 5 & 32053 & 36970 & 38491 &  & 32193 & 0.43 \\
OR10x250-0.75\_2.dat-50 & 75 & 5 & 31550 & 37354 & 38526 &  & 31710 & 0.50 \\
OR10x250-0.75\_3.dat-50 & 75 & 5 & 32078 & 36456 & 38500 &  & 32237 & 0.49 \\
OR10x250-0.75\_4.dat-50 & 75 & 5 & 32823 & 38496 & 40143 &  & 32914 & 0.28 \\
OR10x250-0.75\_5.dat-50 & 75 & 5 & 32161 & 34949 & 39199 &  & 32305 & 0.45 \\
OR10x250-0.75\_6.dat-50 & 75 & 5 & 32547 & 36905 & 38195 &  & 32655 & 0.33 \\
OR10x250-0.75\_7.dat-50 & 75 & 5 & 31906 & 37389 & 37389 &  & 32012 & 0.33 \\
OR10x250-0.75\_8.dat-50 & 75 & 5 & 32573 & 38048 & 39745 &  & 32696 & 0.38 \\
OR10x250-0.75\_9.dat-50 & 75 & 5 & 32563 & 37021 & 40293 &  & 32743 & 0.55 \\
OR10x250-0.75\_10.dat-50 & 75 & 5 & 32629 & 37946 & 39600 &  & 32753 & 0.38 \\
OR10x250-0.75\_1.dat-100 & 75 & 1 & 34846 & 34846 & 38358 & 33561 & 34846 & 0.00 \\
OR10x250-0.75\_2.dat-100 & 75 & 1 & 34304 & 34351 & 40332 & 33042 & 34419 & 0.14 \\
OR10x250-0.75\_3.dat-100 & 75 & 1 & 35647 & 35699 & 39734 & 34336 & 35711 & 0.15 \\
OR10x250-0.75\_4.dat-100 & 75 & 1 & 35302 & 35355 & 41019 & 33997 & 35501 & 0.15 \\
OR10x250-0.75\_5.dat-100 & 75 & 1 & 35887 & 35920 & 40366 & 34662 & 36042 & 0.09 \\
OR10x250-0.75\_6.dat-100 & 75 & 1 & 35101 & 35101 & 39878 & 33724 & 35188 & 0.00 \\
OR10x250-0.75\_7.dat-100 & 75 & 1 & 35021 & 35021 & 40501 & 33697 & 35071 & 0.00 \\
OR10x250-0.75\_8.dat-100 & 75 & 1 & 36089 & 36139 & 41139 & 34875 & 36209 & 0.14 \\
OR10x250-0.75\_9.dat-100 & 75 & 1 & 34584 & 34625 & 38456 & 33260 & 34714 & 0.12 \\
OR10x250-0.75\_10.dat-100 & 75 & 1 & 34852 & 34889 & 40957 & 33600 & 34935 & 0.11 \\
OR10x500-0.25\_1.dat-0 & 25 & 9 &  &  &  &  &  & 100.00 \\
OR10x500-0.25\_2.dat-0 & 25 & 9 &  &  &  &  &  & 100.00 \\
OR10x500-0.25\_3.dat-0 & 25 & 9 &  &  &  &  &  & 100.00 \\
OR10x500-0.25\_4.dat-0 & 25 & 9 &  &  &  &  &  & 100.00 \\
OR10x500-0.25\_5.dat-0 & 25 & 9 &  &  &  &  &  & 100.00 \\
OR10x500-0.25\_6.dat-0 & 25 & 9 &  &  &  &  &  & 100.00 \\
OR10x500-0.25\_7.dat-0 & 25 & 9 &  &  &  &  &  & 100.00 \\
OR10x500-0.25\_8.dat-0 & 25 & 9 &  &  &  &  &  & 100.00 \\
OR10x500-0.25\_9.dat-0 & 25 & 9 &  &  &  &  &  & 100.00 \\
OR10x500-0.25\_10.dat-0 & 25 & 9 &  &  &  &  &  & 100.00 \\
OR10x500-0.25\_1.dat-50 & 25 & 5 & 67440 & 113396 &  &  &  & 36.39 \\
OR10x500-0.25\_2.dat-50 & 25 & 5 & 67265 & 112795 &  &  &  & 38.06 \\
OR10x500-0.25\_3.dat-50 & 25 & 5 & 67129 & 113842 &  &  &  & 38.00 \\
OR10x500-0.25\_4.dat-50 & 25 & 5 & 67267 & 113364 &  &  &  & 36.00 \\
OR10x500-0.25\_5.dat-50$^{\dagger}$ & 25 & 5 & 67175 & 20909 &  &  &  & 36.89 \\
OR10x500-0.25\_6.dat-50 & 25 & 5 & 66729 & 112484 &  &  &  & 37.63 \\
OR10x500-0.25\_7.dat-50 & 25 & 5 & 66534 & 113266 &  &  &  & 38.07 \\
OR10x500-0.25\_8.dat-50 & 25 & 5 & 67532 & 113484 &  &  &  & 37.37 \\
OR10x500-0.25\_9.dat-50 & 25 & 5 & 66584 & 112885 &  &  &  & 37.55 \\
OR10x500-0.25\_10.dat-50 & 25 & 5 & 66920 & 113424 &  &  &  & 38.21 \\
OR10x500-0.25\_1.dat-100 & 25 & 1 & 94743 & 127585 & 127585 & 115162 & 121246 & 5.02 \\
OR10x500-0.25\_2.dat-100 & 25 & 1 & 95933 & 127453 & 127453 & 115492 & 122627 & 5.82 \\
OR10x500-0.25\_3.dat-100 & 25 & 1 & 94994 & 125471 & 125471 & 113509 & 119730 & 5.20 \\
OR10x500-0.25\_4.dat-100 & 25 & 1 & 89911 & 123364 & 123364 & 111864 & 117808 & 5.05 \\
OR10x500-0.25\_5.dat-100 & 25 & 1 & 91545 & 120705 & 120705 & 108352 & 114902 & 5.70 \\
OR10x500-0.25\_6.dat-100 & 25 & 1 & 87880 & 117664 & 117664 & 107695 & 112855 & 4.57 \\
OR10x500-0.25\_7.dat-100 & 25 & 1 & 89663 & 117580 & 117580 & 108463 & 115125 & 5.79 \\
OR10x500-0.25\_8.dat-100 & 25 & 1 & 88272 & 118222 & 118222 & 108534 & 113884 & 4.70 \\
OR10x500-0.25\_9.dat-100 & 25 & 1 & 94528 & 123911 & 123911 & 112103 & 118877 & 5.70 \\
OR10x500-0.25\_10.dat-100 & 25 & 1 & 91633 & 125687 & 125687 & 113444 & 119347 & 4.95 \\
OR10x500-0.50\_1.dat-0 & 50 & 9 &  &  &  &  & 110228 & 100.00 \\
OR10x500-0.50\_2.dat-0 & 50 & 9 &  &  &  &  & 108994 & 100.00 \\
OR10x500-0.50\_3.dat-0 & 50 & 9 &  &  &  &  & 109258 & 100.00 \\
OR10x500-0.50\_4.dat-0 & 50 & 9 &  &  &  &  & 109277 & 100.00 \\
OR10x500-0.50\_5.dat-0 & 50 & 9 &  &  &  &  & 109764 & 100.00 \\
OR10x500-0.50\_6.dat-0 & 50 & 9 &  &  &  &  & 109294 & 100.00 \\
OR10x500-0.50\_7.dat-0 & 50 & 9 &  &  &  &  & 108537 & 100.00 \\
OR10x500-0.50\_8.dat-0 & 50 & 9 &  &  &  &  & 111621 & 100.00 \\
OR10x500-0.50\_9.dat-0 & 50 & 9 &  &  &  &  & 107753 & 100.00 \\
OR10x500-0.50\_10.dat-0 & 50 & 9 &  &  &  &  & 111864 & 100.00 \\
OR10x500-0.50\_1.dat-50 & 50 & 5 & 90553 & 168280 &  &  & 143800 & 37.03 \\
OR10x500-0.50\_2.dat-50 & 50 & 5 & 89644 & 168134 &  &  & 144652 & 38.03 \\
OR10x500-0.50\_3.dat-50 & 50 & 5 & 88559 & 167065 &  &  & 142099 & 37.68 \\
OR10x500-0.50\_4.dat-50 & 50 & 5 & 89058 & 165223 &  &  & 138428 & 35.66 \\
OR10x500-0.50\_5.dat-50 & 50 & 5 & 88025 & 165182 &  &  & 141199 & 37.66 \\
OR10x500-0.50\_6.dat-50 & 50 & 5 & 88666 & 163977 &  &  & 140895 & 37.07 \\
OR10x500-0.50\_7.dat-50 & 50 & 5 & 89541 & 165363 &  &  & 142685 & 37.25 \\
OR10x500-0.50\_8.dat-50 & 50 & 5 & 90548 & 171282 &  &  & 139449 & 35.07 \\
OR10x500-0.50\_9.dat-50 & 50 & 5 & 88168 & 163683 &  &  & 136827 & 35.56 \\
OR10x500-0.50\_10.dat-50 & 50 & 5 & 91352 & 171052 &  &  & 144345 & 36.71 \\
OR10x500-0.50\_1.dat-100 & 50 & 1 & 94989 & 151252 & 151252 & 126056 & 143544 & 12.18 \\
OR10x500-0.50\_2.dat-100 & 50 & 1 & 95347 & 152805 & 152805 & 126165 & 144228 & 12.52 \\
OR10x500-0.50\_3.dat-100 & 50 & 1 & 93960 & 150071 & 150081 & 124877 & 141697 & 11.87 \\
OR10x500-0.50\_4.dat-100 & 50 & 1 & 92771 & 144641 & 144641 & 124705 & 139035 & 10.31 \\
OR10x500-0.50\_5.dat-100 & 50 & 1 & 93142 & 149765 & 149765 & 124880 & 140684 & 11.23 \\
OR10x500-0.50\_6.dat-100 & 50 & 1 & 92805 & 147967 & 148618 & 123818 & 139239 & 11.08 \\
OR10x500-0.50\_7.dat-100 & 50 & 1 & 93245 & 149764 & 149764 & 124042 & 141572 & 12.38 \\
OR10x500-0.50\_8.dat-100 & 50 & 1 & 94131 & 147014 & 147020 & 124989 & 140147 & 10.82 \\
OR10x500-0.50\_9.dat-100 & 50 & 1 & 91805 & 149977 & 149977 & 123180 & 140613 & 12.40 \\
OR10x500-0.50\_10.dat-100 & 50 & 1 & 95673 & 155420 & 155420 & 128690 & 146691 & 12.27 \\
OR10x500-0.75\_1.dat-0 & 75 & 9 & 33983 & 52026 &  &  & 45235 & 24.87 \\
OR10x500-0.75\_2.dat-0 & 75 & 9 &  &  &  &  & 45759 & 100.00 \\
OR10x500-0.75\_3.dat-0 & 75 & 9 &  &  &  &  & 46872 & 100.00 \\
OR10x500-0.75\_4.dat-0 & 75 & 9 &  &  &  &  & 43498 & 100.00 \\
OR10x500-0.75\_5.dat-0 & 75 & 9 &  &  &  &  & 46681 & 100.00 \\
OR10x500-0.75\_6.dat-0 & 75 & 9 &  &  &  &  & 47558 & 100.00 \\
OR10x500-0.75\_7.dat-0 & 75 & 9 &  &  &  &  & 47506 & 100.00 \\
OR10x500-0.75\_8.dat-0 & 75 & 9 &  &  &  &  & 42652 & 100.00 \\
OR10x500-0.75\_9.dat-0 & 75 & 9 &  &  &  &  & 44192 & 100.00 \\
OR10x500-0.75\_10.dat-0 & 75 & 9 &  &  &  &  & 45448 & 100.00 \\
OR10x500-0.75\_1.dat-50 & 75 & 5 & 63789 & 74925 &  &  & 63888 & 0.15 \\
OR10x500-0.75\_2.dat-50 & 75 & 5 & 63960 & 74542 &  &  & 64075 & 0.18 \\
OR10x500-0.75\_3.dat-50 & 75 & 5 & 46714 & 77304 & 290174 &  & 65792 & 29.00 \\
OR10x500-0.75\_4.dat-50 & 75 & 5 & 45938 & 74701 &  &  & 63697 & 27.88 \\
OR10x500-0.75\_5.dat-50 & 75 & 5 & 45616 & 75575 &  &  & 65907 & 30.79 \\
OR10x500-0.75\_6.dat-50 & 75 & 5 & 64635 & 74937 &  &  & 64738 & 0.16 \\
OR10x500-0.75\_7.dat-50 & 75 & 5 & 63946 & 74679 & 101020 &  & 64068 & 0.19 \\
OR10x500-0.75\_8.dat-50 & 75 & 5 & 63071 & 74193 & 299817 &  & 63159 & 0.14 \\
OR10x500-0.75\_9.dat-50 & 75 & 5 & 45646 & 74809 &  &  & 63140 & 27.71 \\
OR10x500-0.75\_10.dat-50 & 75 & 5 & 46380 & 77141 & 96962 &  & 65197 & 28.86 \\
OR10x500-0.75\_1.dat-100 & 75 & 1 & 69335 & 69713 & 96504 & 68272 & 69743 & 0.54 \\
OR10x500-0.75\_2.dat-100 & 75 & 1 & 68643 & 69027 & 98003 & 67621 & 69069 & 0.56 \\
OR10x500-0.75\_3.dat-100 & 75 & 1 & 54281 & 70802 & 94624 & 69492 & 70802 & 1.85 \\
OR10x500-0.75\_4.dat-100 & 75 & 1 & 68445 & 68843 & 95396 & 67373 & 68881 & 0.58 \\
OR10x500-0.75\_5.dat-100 & 75 & 1 & 69744 & 70096 & 94674 & 68543 & 69998 & 0.36 \\
OR10x500-0.75\_6.dat-100 & 75 & 1 & 70268 & 70641 & 95464 & 69207 & 70720 & 0.53 \\
OR10x500-0.75\_7.dat-100 & 75 & 1 & 69661 & 70123 & 97355 & 68671 & 70134 & 0.66 \\
OR10x500-0.75\_8.dat-100 & 75 & 1 & 68849 & 68975 & 93518 & 67667 & 69081 & 0.18 \\
OR10x500-0.75\_9.dat-100 & 75 & 1 & 67835 & 68188 & 95135 & 66816 & 68228 & 0.52 \\
OR10x500-0.75\_10.dat-100 & 75 & 1 & 69841 & 70117 & 98288 & 68727 & 70190 & 0.39 \\
OR30x100-0.25\_1.dat-0 & 25 & 29 & 12896 & 18176 & 20967 &  &  & 28.58 \\
OR30x100-0.25\_2.dat-0 & 25 & 29 & 11887 & 19013 & 20883 &  &  & 34.52 \\
OR30x100-0.25\_3.dat-0 & 25 & 29 & 12408 & 17619 & 20050 &  &  & 29.58 \\
OR30x100-0.25\_4.dat-0 & 25 & 29 &  & 19822 &  &  &  & 100.00 \\
OR30x100-0.25\_5.dat-0 & 25 & 29 &  & 20626 &  &  &  & 100.00 \\
OR30x100-0.25\_6.dat-0 & 25 & 29 & 13153 & 18387 &  &  &  & 28.47 \\
OR30x100-0.25\_7.dat-0 & 25 & 29 & 12125 & 18688 &  &  &  & 32.94 \\
OR30x100-0.25\_8.dat-0 & 25 & 29 &  & 19967 &  &  &  & 100.00 \\
OR30x100-0.25\_9.dat-0 & 25 & 29 &  & 21618 &  &  &  & 100.00 \\
OR30x100-0.25\_10.dat-0 & 25 & 29 & 974 & 18151 & 19939 &  &  & 94.58 \\
OR30x100-0.25\_1.dat-50 & 25 & 15 & 14358 & 19539 & 21678 &  &  & 26.52 \\
OR30x100-0.25\_2.dat-50 & 25 & 15 & 13955 & 19468 & 20351 &  &  & 28.32 \\
OR30x100-0.25\_3.dat-50 & 25 & 15 & 13826 & 18956 & 20892 &  &  & 27.06 \\
OR30x100-0.25\_4.dat-50 & 25 & 15 & 13997 & 19452 & 20591 &  &  & 27.69 \\
OR30x100-0.25\_5.dat-50 & 25 & 15 & 14772 & 20202 & 22158 &  &  & 26.88 \\
OR30x100-0.25\_6.dat-50 & 25 & 15 & 14488 & 19788 & 21075 &  &  & 26.78 \\
OR30x100-0.25\_7.dat-50 & 25 & 15 & 14106 & 19498 & 21055 &  & 19443 & 27.45 \\
OR30x100-0.25\_8.dat-50 & 25 & 15 & 14328 & 19699 & 21345 &  &  & 27.27 \\
OR30x100-0.25\_9.dat-50 & 25 & 15 & 15173 & 20935 & 21651 &  &  & 27.52 \\
OR30x100-0.25\_10.dat-50 & 25 & 15 & 13710 & 18470 & 20780 &  &  & 25.77 \\
OR30x100-0.25\_1.dat-100 & 25 & 1 & 19857 & 24023 & 24689 & 22442 & 23431 & 4.22 \\
OR30x100-0.25\_2.dat-100 & 25 & 1 & 21423 & 25500 & 26280 & 23845 & 24863 & 4.09 \\
OR30x100-0.25\_3.dat-100 & 25 & 1 & 18927 & 23033 & 23988 & 21604 & 22607 & 4.44 \\
OR30x100-0.25\_4.dat-100 & 25 & 1 & 19203 & 23311 & 24900 & 22150 & 23073 & 4.00 \\
OR30x100-0.25\_5.dat-100 & 25 & 1 & 20638 & 25265 & 27414 & 24017 & 24955 & 3.76 \\
OR30x100-0.25\_6.dat-100 & 25 & 1 & 20272 & 24540 & 25644 & 23099 & 24116 & 4.22 \\
OR30x100-0.25\_7.dat-100 & 25 & 1 & 22418 & 27396 & 28746 & 25726 & 26993 & 4.69 \\
OR30x100-0.25\_8.dat-100 & 25 & 1 & 19527 & 23856 & 25096 & 22336 & 23319 & 4.22 \\
OR30x100-0.25\_9.dat-100 & 25 & 1 & 23302 & 27345 & 28174 & 25753 & 26760 & 3.76 \\
OR30x100-0.25\_10.dat-100 & 25 & 1 & 19848 & 24084 & 24908 & 22825 & 23762 & 3.94 \\
OR30x100-0.50\_1.dat-0 & 50 & 29 & 14360 & 22276 & 28821 &  & 20246 & 29.07 \\
OR30x100-0.50\_2.dat-0 & 50 & 29 & 15060 & 23891 & 37927 &  & 21279 & 29.23 \\
OR30x100-0.50\_3.dat-0 & 50 & 29 & 15798 & 25271 &  &  & 23024 & 31.38 \\
OR30x100-0.50\_4.dat-0 & 50 & 29 & 15637 & 24998 & 29193 &  & 22792 & 31.39 \\
OR30x100-0.50\_5.dat-0 & 50 & 29 & 14858 & 22171 &  &  & 20805 & 28.58 \\
OR30x100-0.50\_6.dat-0 & 50 & 29 & 15306 & 23182 & 35501 &  & 21594 & 29.12 \\
OR30x100-0.50\_7.dat-0 & 50 & 29 & 15351 & 24028 & 31149 &  & 21739 & 29.38 \\
OR30x100-0.50\_8.dat-0 & 50 & 29 & 16082 & 24989 & 39030 &  & 23398 & 31.27 \\
OR30x100-0.50\_9.dat-0 & 50 & 29 & 15390 & 24269 & 31571 &  & 22260 & 30.86 \\
OR30x100-0.50\_10.dat-0 & 50 & 29 & 14678 & 22218 & 29171 &  & 20776 & 29.35 \\
OR30x100-0.50\_1.dat-50 & 50 & 15 & 18108 & 30596 & 32058 &  & 27226 & 33.49 \\
OR30x100-0.50\_2.dat-50 & 50 & 15 & 18572 & 31150 & 34004 &  & 27576 & 32.65 \\
OR30x100-0.50\_3.dat-50 & 50 & 15 & 18267 & 30402 & 32721 &  & 27358 & 33.23 \\
OR30x100-0.50\_4.dat-50 & 50 & 15 & 18029 & 30630 & 31866 &  & 26504 & 31.98 \\
OR30x100-0.50\_5.dat-50 & 50 & 15 & 18421 & 31891 & 33177 &  & 27982 & 34.17 \\
OR30x100-0.50\_6.dat-50 & 50 & 15 & 17915 & 29987 & 30968 &  & 27098 & 33.89 \\
OR30x100-0.50\_7.dat-50 & 50 & 15 & 18097 & 30614 & 32053 &  & 27330 & 33.78 \\
OR30x100-0.50\_8.dat-50 & 50 & 15 & 18584 & 31674 & 32923 &  & 28026 & 33.69 \\
OR30x100-0.50\_9.dat-50 & 50 & 15 & 18432 & 31308 & 32481 &  & 27966 & 34.09 \\
OR30x100-0.50\_10.dat-50 & 50 & 15 & 18020 & 29750 & 30905 &  & 27420 & 34.28 \\
OR30x100-0.50\_1.dat-100 & 50 & 1 & 19936 & 29647 & 29647 & 25597 & 27647 & 7.41 \\
OR30x100-0.50\_2.dat-100 & 50 & 1 & 20863 & 29888 & 31795 & 27075 & 28969 & 6.54 \\
OR30x100-0.50\_3.dat-100 & 50 & 1 & 20153 & 29784 & 30364 & 26400 & 28282 & 6.65 \\
OR30x100-0.50\_4.dat-100 & 50 & 1 & 19939 & 28972 & 28972 & 25836 & 27529 & 6.15 \\
OR30x100-0.50\_5.dat-100 & 50 & 1 & 21653 & 30673 & 30673 & 27834 & 29571 & 5.87 \\
OR30x100-0.50\_6.dat-100 & 50 & 1 & 21042 & 30180 & 31040 & 27033 & 28618 & 5.54 \\
OR30x100-0.50\_7.dat-100 & 50 & 1 & 20283 & 30451 & 30452 & 26681 & 28877 & 7.60 \\
OR30x100-0.50\_8.dat-100 & 50 & 1 & 20947 & 29461 & 30154 & 27156 & 28624 & 5.13 \\
OR30x100-0.50\_9.dat-100 & 50 & 1 & 21885 & 29984 & 30731 & 27090 & 28534 & 5.06 \\
OR30x100-0.50\_10.dat-100 & 50 & 1 & 20037 & 30229 & 30229 & 27100 & 28530 & 5.01 \\
OR30x100-0.75\_1.dat-0 & 75 & 29 & 9652 & 10302 & 10772 &  & 9657 & 0.05 \\
OR30x100-0.75\_2.dat-0 & 75 & 29 & 9301 & 9756 & 10657 &  & 9397 & 1.02 \\
OR30x100-0.75\_3.dat-0 & 75 & 29 & 9041 & 9545 & 10543 &  & 9157 & 1.27 \\
OR30x100-0.75\_4.dat-0 & 75 & 29 & 10023 & 10623 & 10847 &  & 10119 & 0.95 \\
OR30x100-0.75\_5.dat-0 & 75 & 29 & 7334 & 7334 & 8411 &  & 7334 & 0.00 \\
OR30x100-0.75\_6.dat-0 & 75 & 29 & 8783 & 9699 & 9906 &  & 8785 & 0.02 \\
OR30x100-0.75\_7.dat-0 & 75 & 29 & 9261 & 9757 & 9757 &  & 9307 & 0.49 \\
OR30x100-0.75\_8.dat-0 & 75 & 29 & 9225 & 9562 & 10097 &  & 9274 & 0.53 \\
OR30x100-0.75\_9.dat-0 & 75 & 29 & 9660 & 10089 & 10937 &  & 9733 & 0.75 \\
OR30x100-0.75\_10.dat-0 & 75 & 29 & 8612 & 9138 & 9138 &  & 8688 & 0.87 \\
OR30x100-0.75\_1.dat-50 & 75 & 15 & 13858 & 15421 & 15421 &  & 14319 & 3.22 \\
OR30x100-0.75\_2.dat-50 & 75 & 15 & 14431 & 16556 & 16626 &  & 14835 & 2.72 \\
OR30x100-0.75\_3.dat-50 & 75 & 15 & 14312 & 16042 & 16731 &  & 14622 & 2.12 \\
OR30x100-0.75\_4.dat-50 & 75 & 15 & 14154 & 16050 & 16716 &  & 14599 & 3.05 \\
OR30x100-0.75\_5.dat-50 & 75 & 15 & 13870 & 15596 & 16815 &  & 14321 & 3.15 \\
OR30x100-0.75\_6.dat-50 & 75 & 15 & 14160 & 15919 & 15919 &  & 14557 & 2.73 \\
OR30x100-0.75\_7.dat-50 & 75 & 15 & 10328 & 15766 & 16527 &  & 14263 & 27.59 \\
OR30x100-0.75\_8.dat-50 & 75 & 15 & 13826 & 15902 & 16438 &  & 14368 & 3.77 \\
OR30x100-0.75\_9.dat-50 & 75 & 15 & 14221 & 15871 & 15871 &  & 14668 & 3.05 \\
OR30x100-0.75\_10.dat-50 & 75 & 15 & 14605 & 16851 & 17182 &  & 15050 & 2.96 \\
OR30x100-0.75\_1.dat-100 & 75 & 1 & 15145 & 15145 & 16706 & 13665 & 15235 & 0.00 \\
OR30x100-0.75\_2.dat-100 & 75 & 1 & 15901 & 15901 & 17154 & 14421 & 15931 & 0.00 \\
OR30x100-0.75\_3.dat-100 & 75 & 1 & 15501 & 15501 & 16795 & 14020 & 15561 & 0.00 \\
OR30x100-0.75\_4.dat-100 & 75 & 1 & 15693 & 15693 & 16687 & 14349 & 15704 & 0.00 \\
OR30x100-0.75\_5.dat-100 & 75 & 1 & 15066 & 15066 & 16563 & 13726 & 15066 & 0.00 \\
OR30x100-0.75\_6.dat-100 & 75 & 1 & 15608 & 15608 & 17235 & 14204 & 15608 & 0.00 \\
OR30x100-0.75\_7.dat-100 & 75 & 1 & 14810 & 14810 & 16254 & 13330 & 14815 & 0.00 \\
OR30x100-0.75\_8.dat-100 & 75 & 1 & 15300 & 15300 & 16672 & 13793 & 15339 & 0.00 \\
OR30x100-0.75\_9.dat-100 & 75 & 1 & 12626 & 15543 & 16977 & 14059 & 15588 & 9.55 \\
OR30x100-0.75\_10.dat-100 & 75 & 1 & 12507 & 16354 & 17842 & 14845 & 16383 & 9.23 \\
OR30x250-0.25\_1.dat-0 & 25 & 29 &  &  &  &  &  & 100.00 \\
OR30x250-0.25\_2.dat-0 & 25 & 29 &  &  &  &  &  & 100.00 \\
OR30x250-0.25\_3.dat-0 & 25 & 29 &  &  &  &  &  & 100.00 \\
OR30x250-0.25\_4.dat-0 & 25 & 29 &  &  &  &  &  & 100.00 \\
OR30x250-0.25\_5.dat-0 & 25 & 29 &  &  &  &  &  & 100.00 \\
OR30x250-0.25\_6.dat-0 & 25 & 29 &  &  &  &  &  & 100.00 \\
OR30x250-0.25\_7.dat-0 & 25 & 29 &  &  &  &  &  & 100.00 \\
OR30x250-0.25\_8.dat-0 & 25 & 29 &  &  &  &  &  & 100.00 \\
OR30x250-0.25\_9.dat-0 & 25 & 29 &  &  &  &  &  & 100.00 \\
OR30x250-0.25\_10.dat-0 & 25 & 29 &  &  &  &  &  & 100.00 \\
OR30x250-0.25\_1.dat-50 & 25 & 15 &  &  &  &  &  & 100.00 \\
OR30x250-0.25\_2.dat-50 & 25 & 15 &  &  &  &  &  & 100.00 \\
OR30x250-0.25\_3.dat-50 & 25 & 15 &  &  &  &  &  & 100.00 \\
OR30x250-0.25\_4.dat-50 & 25 & 15 &  &  &  &  &  & 100.00 \\
OR30x250-0.25\_5.dat-50 & 25 & 15 &  &  &  &  &  & 100.00 \\
OR30x250-0.25\_6.dat-50 & 25 & 15 &  &  &  &  &  & 100.00 \\
OR30x250-0.25\_7.dat-50 & 25 & 15 &  &  &  &  &  & 100.00 \\
OR30x250-0.25\_8.dat-50 & 25 & 15 &  &  &  &  &  & 100.00 \\
OR30x250-0.25\_9.dat-50 & 25 & 15 &  &  &  &  &  & 100.00 \\
OR30x250-0.25\_10.dat-50 & 25 & 15 &  &  &  &  &  & 100.00 \\
OR30x250-0.25\_1.dat-100 & 25 & 1 & 47685 & 62454 & 62454 & 57151 & 59568 & 4.06 \\
OR30x250-0.25\_2.dat-100 & 25 & 1 & 47281 & 63145 & 63145 & 57589 & 60467 & 4.76 \\
OR30x250-0.25\_3.dat-100 & 25 & 1 & 44762 & 61389 & 61552 & 55280 & 58423 & 5.38 \\
OR30x250-0.25\_4.dat-100 & 25 & 1 & 44104 & 59622 & 59622 & 54631 & 57031 & 4.21 \\
OR30x250-0.25\_5.dat-100 & 25 & 1 & 49566 & 64027 & 64027 & 59421 & 62358 & 4.71 \\
OR30x250-0.25\_6.dat-100 & 25 & 1 & 46755 & 62882 & 64184 & 56954 & 59705 & 4.61 \\
OR30x250-0.25\_7.dat-100 & 25 & 1 & 48882 & 63143 & 63143 & 57442 & 60532 & 5.10 \\
OR30x250-0.25\_8.dat-100 & 25 & 1 & 45023 & 60008 & 60008 & 55218 & 57832 & 4.52 \\
OR30x250-0.25\_9.dat-100 & 25 & 1 & 49081 & 65036 & 65036 & 58335 & 61704 & 5.46 \\
OR30x250-0.25\_10.dat-100 & 25 & 1 & 44752 & 60515 & 60515 & 55118 & 57781 & 4.61 \\
OR30x250-0.50\_1.dat-0 & 50 & 29 &  &  &  &  & 57878 & 100.00 \\
OR30x250-0.50\_2.dat-0 & 50 & 29 &  &  &  &  & 55862 & 100.00 \\
OR30x250-0.50\_3.dat-0 & 50 & 29 &  &  &  &  & 53747 & 100.00 \\
OR30x250-0.50\_4.dat-0 & 50 & 29 &  &  &  &  & 53794 & 100.00 \\
OR30x250-0.50\_5.dat-0 & 50 & 29 &  &  &  &  & 51556 & 100.00 \\
OR30x250-0.50\_6.dat-0 & 50 & 29 &  &  &  &  & 54634 & 100.00 \\
OR30x250-0.50\_7.dat-0 & 50 & 29 &  &  &  &  & 52579 & 100.00 \\
OR30x250-0.50\_8.dat-0 & 50 & 29 &  &  &  &  & 55431 & 100.00 \\
OR30x250-0.50\_9.dat-0 & 50 & 29 &  &  &  &  & 54788 & 100.00 \\
OR30x250-0.50\_10.dat-0 & 50 & 29 &  &  &  &  & 53757 & 100.00 \\
OR30x250-0.50\_1.dat-50 & 50 & 15 &  &  &  &  & 69476 & 100.00 \\
OR30x250-0.50\_2.dat-50 & 50 & 15 &  &  &  &  & 71164 & 100.00 \\
OR30x250-0.50\_3.dat-50 & 50 & 15 &  &  &  &  & 69324 & 100.00 \\
OR30x250-0.50\_4.dat-50 & 50 & 15 &  &  &  &  & 70897 & 100.00 \\
OR30x250-0.50\_5.dat-50 & 50 & 15 &  &  &  &  & 68846 & 100.00 \\
OR30x250-0.50\_6.dat-50 & 50 & 15 &  &  &  &  & 71186 & 100.00 \\
OR30x250-0.50\_7.dat-50 & 50 & 15 &  &  &  &  & 67987 & 100.00 \\
OR30x250-0.50\_8.dat-50 & 50 & 15 &  &  &  &  & 67765 & 100.00 \\
OR30x250-0.50\_9.dat-50 & 50 & 15 &  &  &  &  & 69846 & 100.00 \\
OR30x250-0.50\_10.dat-50 & 50 & 15 &  &  &  &  & 68031 & 100.00 \\
OR30x250-0.50\_1.dat-100 & 50 & 1 & 47439 & 66391 & 76643 & 66200 & 71552 & 0.29 \\
OR30x250-0.50\_2.dat-100 & 50 & 1 & 51684 & 78437 & 78796 & 66953 & 72889 & 8.14 \\
OR30x250-0.50\_3.dat-100 & 50 & 1 & 47608 & 77047 & 77047 & 66136 & 72226 & 8.43 \\
OR30x250-0.50\_4.dat-100 & 50 & 1 & 46976 & 74221 & 74221 & 65658 & 70800 & 7.26 \\
OR30x250-0.50\_5.dat-100 & 50 & 1 & 46896 & 75191 & 75191 & 65974 & 71547 & 7.79 \\
OR30x250-0.50\_6.dat-100 & 50 & 1 & 52961 & 75033 & 75033 & 66853 & 71874 & 6.99 \\
OR30x250-0.50\_7.dat-100 & 50 & 1 & 51425 & 74651 & 74651 & 65133 & 70073 & 7.05 \\
OR30x250-0.50\_8.dat-100 & 50 & 1 & 45524 & 72609 & 72609 & 63845 & 68983 & 7.45 \\
OR30x250-0.50\_9.dat-100 & 50 & 1 & 51372 & 73939 & 73939 & 65269 & 70341 & 7.21 \\
OR30x250-0.50\_10.dat-100 & 50 & 1 & 51170 & 72859 & 72859 & 64963 & 69569 & 6.62 \\
OR30x250-0.75\_1.dat-0 & 75 & 29 &  &  &  &  & 20531 & 100.00 \\
OR30x250-0.75\_2.dat-0 & 75 & 29 &  &  &  &  & 23922 & 100.00 \\
OR30x250-0.75\_3.dat-0 & 75 & 29 &  &  &  &  & 22432 & 100.00 \\
OR30x250-0.75\_4.dat-0 & 75 & 29 &  &  &  &  & 24409 & 100.00 \\
OR30x250-0.75\_5.dat-0 & 75 & 29 & 17158 & 24747 &  &  & 21944 & 21.81 \\
OR30x250-0.75\_6.dat-0 & 75 & 29 &  &  &  &  & 22434 & 100.00 \\
OR30x250-0.75\_7.dat-0 & 75 & 29 &  &  &  &  & 22141 & 100.00 \\
OR30x250-0.75\_8.dat-0 & 75 & 29 &  &  &  &  & 22793 & 100.00 \\
OR30x250-0.75\_9.dat-0 & 75 & 29 &  &  &  &  & 22874 & 100.00 \\
OR30x250-0.75\_10.dat-0 & 75 & 29 & 17213 & 27191 &  &  & 23629 & 27.15 \\
OR30x250-0.75\_1.dat-50 & 75 & 15 & 23195 & 40774 &  &  & 34678 & 33.11 \\
OR30x250-0.75\_2.dat-50 & 75 & 15 & 23642 & 41815 &  &  & 34967 & 32.39 \\
OR30x250-0.75\_3.dat-50 & 75 & 15 & 23943 & 40383 &  &  & 34184 & 29.96 \\
OR30x250-0.75\_4.dat-50 & 75 & 15 & 23596 & 41271 &  &  & 35785 & 34.06 \\
OR30x250-0.75\_5.dat-50 & 75 & 15 & 24119 & 41200 &  &  & 35026 & 31.14 \\
OR30x250-0.75\_6.dat-50 & 75 & 15 & 23627 & 40846 &  &  & 35722 & 33.86 \\
OR30x250-0.75\_7.dat-50 & 75 & 15 & 23374 & 40010 &  &  & 34015 & 31.28 \\
OR30x250-0.75\_8.dat-50 & 75 & 15 & 24081 & 41760 &  &  & 35598 & 32.35 \\
OR30x250-0.75\_9.dat-50 & 75 & 15 & 23681 & 41135 &  &  & 35634 & 33.54 \\
OR30x250-0.75\_10.dat-50 & 75 & 15 & 23516 & 41566 &  &  & 34659 & 32.15 \\
OR30x250-0.75\_1.dat-100 & 75 & 1 & 35895 & 36063 & 40966 & 34538 & 36237 & 0.47 \\
OR30x250-0.75\_2.dat-100 & 75 & 1 & 36220 & 36430 & 40815 & 34945 & 36504 & 0.58 \\
OR30x250-0.75\_3.dat-100 & 75 & 1 & 36391 & 36609 & 41476 & 35083 & 36684 & 0.60 \\
OR30x250-0.75\_4.dat-100 & 75 & 1 & 36846 & 37032 & 41307 & 35193 & 37261 & 0.50 \\
OR30x250-0.75\_5.dat-100 & 75 & 1 & 36130 & 36271 & 41647 & 34861 & 36418 & 0.39 \\
OR30x250-0.75\_6.dat-100 & 75 & 1 & 36439 & 36648 & 41492 & 35080 & 36771 & 0.57 \\
OR30x250-0.75\_7.dat-100 & 75 & 1 & 35243 & 35461 & 40829 & 33960 & 35508 & 0.61 \\
OR30x250-0.75\_8.dat-100 & 75 & 1 & 37215 & 37353 & 42512 & 35944 & 37562 & 0.37 \\
OR30x250-0.75\_9.dat-100 & 75 & 1 & 36702 & 36862 & 41515 & 34968 & 37109 & 0.43 \\
OR30x250-0.75\_10.dat-100 & 75 & 1 & 36601 & 36808 & 41351 & 35348 & 36915 & 0.56 \\
OR30x500-0.25\_1.dat-0 & 25 & 29 &  &  &  &  &  & 100.00 \\
OR30x500-0.25\_2.dat-0 & 25 & 29 &  &  &  &  &  & 100.00 \\
OR30x500-0.25\_3.dat-0 & 25 & 29 &  &  &  &  &  & 100.00 \\
OR30x500-0.25\_4.dat-0 & 25 & 29 &  &  &  &  &  & 100.00 \\
OR30x500-0.25\_5.dat-0 & 25 & 29 &  &  &  &  &  & 100.00 \\
OR30x500-0.25\_6.dat-0 & 25 & 29 &  &  &  &  &  & 100.00 \\
OR30x500-0.25\_7.dat-0 & 25 & 29 &  &  &  &  &  & 100.00 \\
OR30x500-0.25\_8.dat-0 & 25 & 29 &  &  &  &  &  & 100.00 \\
OR30x500-0.25\_9.dat-0 & 25 & 29 &  &  &  &  &  & 100.00 \\
OR30x500-0.25\_10.dat-0 & 25 & 29 &  &  &  &  &  & 100.00 \\
OR30x500-0.25\_1.dat-50 & 25 & 15 &  &  &  &  &  & 100.00 \\
OR30x500-0.25\_2.dat-50 & 25 & 15 &  &  &  &  &  & 100.00 \\
OR30x500-0.25\_3.dat-50 & 25 & 15 &  &  &  &  &  & 100.00 \\
OR30x500-0.25\_4.dat-50 & 25 & 15 &  &  &  &  &  & 100.00 \\
OR30x500-0.25\_5.dat-50 & 25 & 15 &  &  &  &  &  & 100.00 \\
OR30x500-0.25\_6.dat-50 & 25 & 15 &  &  &  &  &  & 100.00 \\
OR30x500-0.25\_7.dat-50 & 25 & 15 &  &  &  &  &  & 100.00 \\
OR30x500-0.25\_8.dat-50 & 25 & 15 &  &  &  &  &  & 100.00 \\
OR30x500-0.25\_9.dat-50 & 25 & 15 &  &  &  &  &  & 100.00 \\
OR30x500-0.25\_10.dat-50 & 25 & 15 &  &  &  &  &  & 100.00 \\
OR30x500-0.25\_1.dat-100 & 25 & 1 & 89570 & 122025 & 122025 & 110666 & 117538 & 5.85 \\
OR30x500-0.25\_2.dat-100 & 25 & 1 & 91797 & 121733 & 121733 & 109040 & 116234 & 6.19 \\
OR30x500-0.25\_3.dat-100 & 25 & 1 & 95947 & 127974 & 127974 & 115167 & 123627 & 6.84 \\
OR30x500-0.25\_4.dat-100 & 25 & 1 & 92444 & 123519 & 123519 & 110414 & 117614 & 6.12 \\
OR30x500-0.25\_5.dat-100 & 25 & 1 & 91811 & 121916 & 121916 & 110217 & 116721 & 5.57 \\
OR30x500-0.25\_6.dat-100 & 25 & 1 & 92537 & 127402 & 127402 & 112511 & 120018 & 6.25 \\
OR30x500-0.25\_7.dat-100 & 25 & 1 & 86807 & 118761 & 118761 & 106181 & 113072 & 6.09 \\
OR30x500-0.25\_8.dat-100 & 25 & 1 & 88265 & 123281 & 123281 & 107987 & 114185 & 5.43 \\
OR30x500-0.25\_9.dat-100 & 25 & 1 & 88718 & 121264 & 121264 & 108776 & 115608 & 5.91 \\
OR30x500-0.25\_10.dat-100 & 25 & 1 & 93353 & 124000 & 124000 & 111669 & 119793 & 6.78 \\
OR30x500-0.50\_1.dat-0 & 50 & 29 &  &  &  &  & 108021 & 100.00 \\
OR30x500-0.50\_2.dat-0 & 50 & 29 &  &  &  &  & 107585 & 100.00 \\
OR30x500-0.50\_3.dat-0 & 50 & 29 &  &  &  &  & 103067 & 100.00 \\
OR30x500-0.50\_4.dat-0 & 50 & 29 &  &  &  &  & 108935 & 100.00 \\
OR30x500-0.50\_5.dat-0 & 50 & 29 &  &  &  &  & 108370 & 100.00 \\
OR30x500-0.50\_6.dat-0 & 50 & 29 &  &  &  &  & 109852 & 100.00 \\
OR30x500-0.50\_7.dat-0 & 50 & 29 &  &  &  &  & 101196 & 100.00 \\
OR30x500-0.50\_8.dat-0 & 50 & 29 &  &  &  &  & 104091 & 100.00 \\
OR30x500-0.50\_9.dat-0 & 50 & 29 &  &  &  &  & 109002 & 100.00 \\
OR30x500-0.50\_10.dat-0 & 50 & 29 &  &  &  &  & 109755 & 100.00 \\
OR30x500-0.50\_1.dat-50 & 50 & 15 &  &  &  &  & 142560 & 100.00 \\
OR30x500-0.50\_2.dat-50 & 50 & 15 &  &  &  &  & 140696 & 100.00 \\
OR30x500-0.50\_3.dat-50 & 50 & 15 &  &  &  &  & 140130 & 100.00 \\
OR30x500-0.50\_4.dat-50 & 50 & 15 &  &  &  &  & 140690 & 100.00 \\
OR30x500-0.50\_5.dat-50 & 50 & 15 &  &  &  &  & 139185 & 100.00 \\
OR30x500-0.50\_6.dat-50 & 50 & 15 &  &  &  &  & 140397 & 100.00 \\
OR30x500-0.50\_7.dat-50 & 50 & 15 &  &  &  &  & 137570 & 100.00 \\
OR30x500-0.50\_8.dat-50 & 50 & 15 &  &  &  &  & 138542 & 100.00 \\
OR30x500-0.50\_9.dat-50 & 50 & 15 &  &  &  &  & 143174 & 100.00 \\
OR30x500-0.50\_10.dat-50 & 50 & 15 &  &  &  &  & 138253 & 100.00 \\
OR30x500-0.50\_1.dat-100 & 50 & 1 & 92736 & 148902 & 148902 & 121643 & 140962 & 13.71 \\
OR30x500-0.50\_2.dat-100 & 50 & 1 & 92207 & 147816 & 147816 & 123890 & 138905 & 10.81 \\
OR30x500-0.50\_3.dat-100 & 50 & 1 & 92248 & 148928 & 148928 & 123199 & 139553 & 11.72 \\
OR30x500-0.50\_4.dat-100 & 50 & 1 & 94098 & 147747 & 147747 & 125054 & 141383 & 11.55 \\
OR30x500-0.50\_5.dat-100 & 50 & 1 & 92187 & 145957 & 145957 & 123385 & 138983 & 11.22 \\
OR30x500-0.50\_6.dat-100 & 50 & 1 & 93745 & 147913 & 147913 & 125045 & 142158 & 12.04 \\
OR30x500-0.50\_7.dat-100 & 50 & 1 & 92239 & 147833 & 147833 & 125268 & 140414 & 10.79 \\
OR30x500-0.50\_8.dat-100 & 50 & 1 & 92092 & 146197 & 146197 & 120078 & 137682 & 12.79 \\
OR30x500-0.50\_9.dat-100 & 50 & 1 & 92908 & 147501 & 147501 & 123625 & 139773 & 11.55 \\
OR30x500-0.50\_10.dat-100 & 50 & 1 & 91157 & 146327 & 146435 & 121197 & 137127 & 11.62 \\
OR30x500-0.75\_1.dat-0 & 75 & 29 &  &  &  &  & 43478 & 100.00 \\
OR30x500-0.75\_2.dat-0 & 75 & 29 &  &  &  &  & 44990 & 100.00 \\
OR30x500-0.75\_3.dat-0 & 75 & 29 &  &  &  &  & 42837 & 100.00 \\
OR30x500-0.75\_4.dat-0 & 75 & 29 &  &  &  &  & 44294 & 100.00 \\
OR30x500-0.75\_5.dat-0 & 75 & 29 &  &  &  &  & 46644 & 100.00 \\
OR30x500-0.75\_6.dat-0 & 75 & 29 &  &  &  &  & 45373 & 100.00 \\
OR30x500-0.75\_7.dat-0 & 75 & 29 &  &  &  &  & 45390 & 100.00 \\
OR30x500-0.75\_8.dat-0 & 75 & 29 &  &  &  &  & 48711 & 100.00 \\
OR30x500-0.75\_9.dat-0 & 75 & 29 &  &  &  &  & 43827 & 100.00 \\
OR30x500-0.75\_10.dat-0 & 75 & 29 &  &  &  &  & 45775 & 100.00 \\
OR30x500-0.75\_1.dat-50 & 75 & 15 &  &  &  &  & 68720 & 100.00 \\
OR30x500-0.75\_2.dat-50 & 75 & 15 &  &  &  &  & 69640 & 100.00 \\
OR30x500-0.75\_3.dat-50 & 75 & 15 &  &  &  &  & 67961 & 100.00 \\
OR30x500-0.75\_4.dat-50 & 75 & 15 &  &  &  &  & 68523 & 100.00 \\
OR30x500-0.75\_5.dat-50 & 75 & 15 &  &  &  &  & 70609 & 100.00 \\
OR30x500-0.75\_6.dat-50 & 75 & 15 &  &  &  &  & 68309 & 100.00 \\
OR30x500-0.75\_7.dat-50 & 75 & 15 &  &  &  &  & 70946 & 100.00 \\
OR30x500-0.75\_8.dat-50 & 75 & 15 &  &  &  &  & 69553 & 100.00 \\
OR30x500-0.75\_9.dat-50 & 75 & 15 &  &  &  &  & 69181 & 100.00 \\
OR30x500-0.75\_10.dat-50 & 75 & 15 &  &  &  &  & 69553 & 100.00 \\
OR30x500-0.75\_1.dat-100 & 75 & 1 & 70260 & 71066 & 102997 & 69453 & 71091 & 1.13 \\
OR30x500-0.75\_2.dat-100 & 75 & 1 & 54615 & 72321 & 102633 & 70617 & 72265 & 2.28 \\
OR30x500-0.75\_3.dat-100 & 75 & 1 & 69924 & 70774 & 97761 & 69109 & 70777 & 1.20 \\
OR30x500-0.75\_4.dat-100 & 75 & 1 & 70645 & 71580 & 95980 & 68950 & 71456 & 1.13 \\
OR30x500-0.75\_5.dat-100 & 75 & 1 & 72234 & 73108 & 103856 & 70425 & 73245 & 1.20 \\
OR30x500-0.75\_6.dat-100 & 75 & 1 & 70312 & 71195 & 95892 & 69397 & 71025 & 1.00 \\
OR30x500-0.75\_7.dat-100 & 75 & 1 & 72604 & 74658 & 92638 & 70376 & 73314 & 0.97 \\
OR30x500-0.75\_8.dat-100 & 75 & 1 & 70708 & 71683 & 103231 & 69876 & 71592 & 1.23 \\
OR30x500-0.75\_9.dat-100 & 75 & 1 & 71479 & 72484 & 96781 & 70708 & 72384 & 1.25 \\
OR30x500-0.75\_10.dat-100 & 75 & 1 & 71384 & 72146 & 91017 & 70444 & 72101 & 0.99 \\
\end{longtable}
}


\subsection{Clique interdiction}
% Generated by paper/scripts/gen_appendix_tables.py from bobilib/results/; do not edit by hand.
% Needs booktabs, longtable, siunitx (v3) and multicol in the preamble.
% ROUND2: update when scripts/check_clique_certificates.py has checked all 170 case certificates; until then the density-0.9 lower bounds are class (ii) and the caption must say so.

{\footnotesize\sisetup{group-separator={,},group-minimum-digits=4,group-digits=integer}\setlength{\tabcolsep}{4pt}\setlength{\LTcapwidth}{\linewidth}
\begin{longtable}{@{}l S[table-format=2.0] S[table-format=1.1] S[table-format=4.0] S[table-format=3.0] S[table-format=2.0] S[table-format=1.0] l S[table-format=1.0] l l@{}}
\caption{Edge clique interdiction, BOBILib set \texttt{clique}: $n$ vertices, density $d$, $m$ edges, budget $k$ (the leader deletes at most $k$ edges). Lower bound: Mantel's theorem ($m-k>\lfloor n^2/4\rfloor$, so a triangle survives) or Brouwer's 1981 Theorem~2 ($m-k\ge t_3(n)-\lfloor n/3\rfloor+2$, so a $K_4$-free remainder is 3-partite) together with the proof that every 3-partition has more than $k$ edges inside its parts ($w_3(G)>k$): \NumCliquePartCases{} part-size cases, each decided by solvers and by exhaustive integer searches, with case certificates checked by the integer checker \texttt{scripts/check\_clique\_certificates.py}. Upper bound: the solution file deletes the edges inside an $r$-partition (certificate: $r$ and the number of inside edges, at most $k$), leaving an $r$-partite graph. Both bounds are exact certificates (class (i) of Section~\ref{app:classes}). Open: the best 3-partitions found for normalClique60-70-2, normalClique60-70-3, normalClique60-70-5 have 311, 310, 310 inside edges, more than $k=309$, so only $3\le{}$value${}\le 4$ is known.}\label{tab:clique}\\
\toprule
Instance & {$n$} & {$d$} & {$m$} & {$k$} & {BOBILib} & {LB} & Argument & {UB} & Certificate & Status \\
\midrule
\endfirsthead
\multicolumn{11}{l}{\emph{(continued)}}\\
\toprule
Instance & {$n$} & {$d$} & {$m$} & {$k$} & {BOBILib} & {LB} & Argument & {UB} & Certificate & Status \\
\midrule
\endhead
\midrule
\multicolumn{11}{r}{\emph{continued on the next page}}\\
\endfoot
\bottomrule
\endlastfoot
normalClique40-70-1 & 40 & 0.7 & 546 & 136 & 6 & 3 & Mantel & 3 & 3-partition, 125 & optimal \\
normalClique40-70-2 & 40 & 0.7 & 546 & 136 & 6 & 3 & Mantel & 3 & 3-partition, 125 & optimal \\
normalClique40-70-3 & 40 & 0.7 & 546 & 136 & 6 & 3 & Mantel & 3 & 3-partition, 126 & optimal \\
normalClique40-70-4 & 40 & 0.7 & 546 & 136 & 5 & 3 & Mantel & 3 & 3-partition, 125 & optimal \\
normalClique40-70-5 & 40 & 0.7 & 546 & 136 & 7 & 3 & Mantel & 3 & 3-partition, 128 & optimal \\
normalClique40-70-6 & 40 & 0.7 & 546 & 136 & 6 & 3 & Mantel & 3 & 3-partition, 126 & optimal \\
normalClique40-70-7 & 40 & 0.7 & 546 & 136 & 5 & 3 & Mantel & 3 & 3-partition, 123 & optimal \\
normalClique40-70-8 & 40 & 0.7 & 546 & 136 & 6 & 3 & Mantel & 3 & 3-partition, 126 & optimal \\
normalClique40-70-9 & 40 & 0.7 & 546 & 136 & 6 & 3 & Mantel & 3 & 3-partition, 125 & optimal \\
normalClique40-70-10 & 40 & 0.7 & 546 & 136 & 6 & 3 & Mantel & 3 & 3-partition, 125 & optimal \\
normalClique40-90-1 & 40 & 0.9 & 702 & 175 & 10 & 4 & Brouwer${}+w_3(G)>k$ & 4 & 4-partition, 131 & optimal \\
normalClique40-90-2 & 40 & 0.9 & 702 & 175 & 10 & 4 & Brouwer${}+w_3(G)>k$ & 4 & 4-partition, 129 & optimal \\
normalClique40-90-3 & 40 & 0.9 & 702 & 175 & 9 & 4 & Brouwer${}+w_3(G)>k$ & 4 & 4-partition, 128 & optimal \\
normalClique40-90-4 & 40 & 0.9 & 702 & 175 & 9 & 4 & Brouwer${}+w_3(G)>k$ & 4 & 4-partition, 127 & optimal \\
normalClique40-90-5 & 40 & 0.9 & 702 & 175 & 10 & 4 & Brouwer${}+w_3(G)>k$ & 4 & 4-partition, 128 & optimal \\
normalClique40-90-6 & 40 & 0.9 & 702 & 175 & 10 & 4 & Brouwer${}+w_3(G)>k$ & 4 & 4-partition, 127 & optimal \\
normalClique40-90-7 & 40 & 0.9 & 702 & 175 & 9 & 4 & Brouwer${}+w_3(G)>k$ & 4 & 4-partition, 127 & optimal \\
normalClique40-90-8 & 40 & 0.9 & 702 & 175 & 11 & 4 & Brouwer${}+w_3(G)>k$ & 4 & 4-partition, 129 & optimal \\
normalClique40-90-9 & 40 & 0.9 & 702 & 175 & 10 & 4 & Brouwer${}+w_3(G)>k$ & 4 & 4-partition, 130 & optimal \\
normalClique40-90-10 & 40 & 0.9 & 702 & 175 & 10 & 4 & Brouwer${}+w_3(G)>k$ & 4 & 4-partition, 130 & optimal \\
normalClique50-70-1 & 50 & 0.7 & 857 & 214 & 7 & 3 & Mantel & 3 & 3-partition, 205 & optimal \\
normalClique50-70-2 & 50 & 0.7 & 857 & 214 & 7 & 3 & Mantel & 3 & 3-partition, 205 & optimal \\
normalClique50-70-3 & 50 & 0.7 & 857 & 214 & 7 & 3 & Mantel & 3 & 3-partition, 208 & optimal \\
normalClique50-70-4 & 50 & 0.7 & 857 & 214 & 7 & 3 & Mantel & 3 & 3-partition, 203 & optimal \\
normalClique50-70-5 & 50 & 0.7 & 857 & 214 & 7 & 3 & Mantel & 3 & 3-partition, 205 & optimal \\
normalClique50-70-6 & 50 & 0.7 & 857 & 214 & 8 & 3 & Mantel & 3 & 3-partition, 209 & optimal \\
normalClique50-70-7 & 50 & 0.7 & 857 & 214 & 8 & 3 & Mantel & 3 & 3-partition, 205 & optimal \\
normalClique50-70-8 & 50 & 0.7 & 857 & 214 & 7 & 3 & Mantel & 3 & 3-partition, 204 & optimal \\
normalClique50-70-9 & 50 & 0.7 & 857 & 214 & 7 & 3 & Mantel & 3 & 3-partition, 210 & optimal \\
normalClique50-70-10 & 50 & 0.7 & 857 & 214 & 8 & 3 & Mantel & 3 & 3-partition, 208 & optimal \\
normalClique50-90-1 & 50 & 0.9 & 1102 & 275 & 11 & 4 & Brouwer${}+w_3(G)>k$ & 4 & 4-partition, 210 & optimal \\
normalClique50-90-2 & 50 & 0.9 & 1102 & 275 & 11 & 4 & Brouwer${}+w_3(G)>k$ & 4 & 4-partition, 212 & optimal \\
normalClique50-90-3 & 50 & 0.9 & 1102 & 275 & 12 & 4 & Brouwer${}+w_3(G)>k$ & 4 & 4-partition, 210 & optimal \\
normalClique50-90-4 & 50 & 0.9 & 1102 & 275 & 11 & 4 & Brouwer${}+w_3(G)>k$ & 4 & 4-partition, 210 & optimal \\
normalClique50-90-5 & 50 & 0.9 & 1102 & 275 & 12 & 4 & Brouwer${}+w_3(G)>k$ & 4 & 4-partition, 212 & optimal \\
normalClique50-90-6 & 50 & 0.9 & 1102 & 275 & 12 & 4 & Brouwer${}+w_3(G)>k$ & 4 & 4-partition, 212 & optimal \\
normalClique50-90-7 & 50 & 0.9 & 1102 & 275 & 11 & 4 & Brouwer${}+w_3(G)>k$ & 4 & 4-partition, 210 & optimal \\
normalClique50-90-8 & 50 & 0.9 & 1102 & 275 & 13 & 4 & Brouwer${}+w_3(G)>k$ & 4 & 4-partition, 211 & optimal \\
normalClique50-90-9 & 50 & 0.9 & 1102 & 275 & 12 & 4 & Brouwer${}+w_3(G)>k$ & 4 & 4-partition, 209 & optimal \\
normalClique50-90-10 & 50 & 0.9 & 1102 & 275 & 12 & 4 & Brouwer${}+w_3(G)>k$ & 4 & 4-partition, 209 & optimal \\
normalClique60-70-1 & 60 & 0.7 & 1239 & 309 & 8 & 3 & Mantel & 3 & 3-partition, 306 & optimal \\
normalClique60-70-2 & 60 & 0.7 & 1239 & 309 & 9 & 3 & Mantel & 4 & 4-partition, 202 & \textbf{open} \\
normalClique60-70-3 & 60 & 0.7 & 1239 & 309 & 9 & 3 & Mantel & 4 & 4-partition, 200 & \textbf{open} \\
normalClique60-70-4 & 60 & 0.7 & 1239 & 309 & 9 & 3 & Mantel & 3 & 3-partition, 309 & optimal \\
normalClique60-70-5 & 60 & 0.7 & 1239 & 309 & 8 & 3 & Mantel & 4 & 4-partition, 202 & \textbf{open} \\
normalClique60-70-6 & 60 & 0.7 & 1239 & 309 & 9 & 3 & Mantel & 3 & 3-partition, 306 & optimal \\
normalClique60-70-7 & 60 & 0.7 & 1239 & 309 & 8 & 3 & Mantel & 3 & 3-partition, 307 & optimal \\
normalClique60-70-8 & 60 & 0.7 & 1239 & 309 & 8 & 3 & Mantel & 3 & 3-partition, 308 & optimal \\
normalClique60-70-9 & 60 & 0.7 & 1239 & 309 & 8 & 3 & Mantel & 3 & 3-partition, 307 & optimal \\
normalClique60-70-10 & 60 & 0.7 & 1239 & 309 & 8 & 3 & Mantel & 3 & 3-partition, 308 & optimal \\
normalClique60-90-1 & 60 & 0.9 & 1593 & 398 & 13 & 4 & Brouwer${}+w_3(G)>k$ & 4 & 4-partition, 314 & optimal \\
normalClique60-90-2 & 60 & 0.9 & 1593 & 398 & 13 & 4 & Brouwer${}+w_3(G)>k$ & 4 & 4-partition, 312 & optimal \\
normalClique60-90-3 & 60 & 0.9 & 1593 & 398 & 13 & 4 & Brouwer${}+w_3(G)>k$ & 4 & 4-partition, 311 & optimal \\
normalClique60-90-4 & 60 & 0.9 & 1593 & 398 & 12 & 4 & Brouwer${}+w_3(G)>k$ & 4 & 4-partition, 313 & optimal \\
normalClique60-90-5 & 60 & 0.9 & 1593 & 398 & 12 & 4 & Brouwer${}+w_3(G)>k$ & 4 & 4-partition, 311 & optimal \\
normalClique60-90-6 & 60 & 0.9 & 1593 & 398 & 13 & 4 & Brouwer${}+w_3(G)>k$ & 4 & 4-partition, 314 & optimal \\
normalClique60-90-7 & 60 & 0.9 & 1593 & 398 & 13 & 4 & Brouwer${}+w_3(G)>k$ & 4 & 4-partition, 313 & optimal \\
normalClique60-90-8 & 60 & 0.9 & 1593 & 398 & 12 & 4 & Brouwer${}+w_3(G)>k$ & 4 & 4-partition, 310 & optimal \\
normalClique60-90-9 & 60 & 0.9 & 1593 & 398 & 12 & 4 & Brouwer${}+w_3(G)>k$ & 4 & 4-partition, 312 & optimal \\
normalClique60-90-10 & 60 & 0.9 & 1593 & 398 & 12 & 4 & Brouwer${}+w_3(G)>k$ & 4 & 4-partition, 317 & optimal \\
\end{longtable}
}


\subsection{Network interdiction}
% Generated by paper/scripts/gen_appendix_tables.py from bobilib/results/; do not edit by hand.
% Needs booktabs, longtable, siunitx (v3) and multicol in the preamble.

{\footnotesize\sisetup{group-separator={,},group-minimum-digits=4,group-digits=integer}\setlength{\tabcolsep}{4pt}\setlength{\LTcapwidth}{\linewidth}
\begin{longtable}{@{}l S[table-format=2.0] S[table-format=2.0] S[table-format=1.0] S[table-format=1.0] S[table-format=7.0] S[table-format=2.0] S[table-format=2.0] S[table-format=2.0] S[table-format=2.0]@{}}
\caption{The two attack-protect games (\texttt{inter-fire}) BOBILib lists as feasible, proven optimal by enumerating every attack of at most $\Phi$ nodes, with the follower (protect at most $\Lambda$ nodes) solved exactly by a tree dynamic program and the minimiser re-solved by HiGHS. Time: wall seconds of the enumeration on one thread of a loaded machine.}\label{tab:network}\\
\toprule
Instance & {Nodes} & {Edges} & {$\Phi$} & {$\Lambda$} & {Attacks} & {Optimum} & {\shortstack{Optimal\\attacks}} & {BOBILib} & {Time (s)} \\
\midrule
\endfirsthead
\multicolumn{10}{l}{\emph{(continued)}}\\
\toprule
Instance & {Nodes} & {Edges} & {$\Phi$} & {$\Lambda$} & {Attacks} & {Optimum} & {\shortstack{Optimal\\attacks}} & {BOBILib} & {Time (s)} \\
\midrule
\endhead
\midrule
\multicolumn{10}{r}{\emph{continued on the next page}}\\
\endfoot
\bottomrule
\endlastfoot
tree-50\_1-5-5\_006 & 49 & 42 & 5 & 5 & 2138410 & 26 & 4 & 26 & 80 \\
tree-50\_1-5-5\_010 & 49 & 44 & 5 & 5 & 2138410 & 26 & 28 & 26 & 70 \\
\end{longtable}
}


{\footnotesize\sisetup{group-separator={,},group-minimum-digits=4,group-digits=integer}\setlength{\tabcolsep}{4pt}\setlength{\LTcapwidth}{\linewidth}
\begin{longtable}{@{}l S[table-format=1.0] S[table-format=2.0] S[table-format=1.0] S[table-format=1.0] S[table-format=7.0] l l@{}}
\caption{The 34 tree instances BOBILib lists as optimal, re-derived by the same enumeration: instances per group, nodes (after removing the vaccinated ones), budgets, attacks per instance, the optimal values found, and how many equal BOBILib's value with a certified solution file.}\label{tab:network-trees}\\
\toprule
Group & {Inst.} & {Nodes} & {$\Phi$} & {$\Lambda$} & {Attacks} & Optima & Agree \\
\midrule
\endfirsthead
\multicolumn{8}{l}{\emph{(continued)}}\\
\toprule
Group & {Inst.} & {Nodes} & {$\Phi$} & {$\Lambda$} & {Attacks} & Optima & Agree \\
\midrule
\endhead
\midrule
\multicolumn{8}{r}{\emph{continued on the next page}}\\
\endfoot
\bottomrule
\endlastfoot
tree-25\_1-3-3 & 4 & 24 & 3 & 3 & 2325 & 12, 13, 13, 13 & 4/4 \\
tree-25\_2-2-2 & 4 & 23 & 2 & 2 & 277 & 17, 17, 17, 18 & 4/4 \\
tree-25\_3-3-3 & 4 & 22 & 3 & 3 & 1794 & 15, 16, 16, 17 & 4/4 \\
tree-50\_1-3-3 & 4 & 49 & 3 & 3 & 19650 & 28, 28, 30, 32 & 4/4 \\
tree-50\_1-5-5 & 2 & 49 & 5 & 5 & 2138410 & 24, 25 & 2/2 \\
tree-50\_3-3-1 & 4 & 47 & 3 & 1 & 17344 & 26, 27, 30, 31 & 4/4 \\
tree-80\_2-2-2 & 4 & 78 & 2 & 2 & 3082 & 56, 57, 61, 63 & 4/4 \\
tree-80\_3-3-1 & 4 & 77 & 3 & 1 & 76154 & 45, 45, 48, 48 & 4/4 \\
tree-80\_3-3-3 & 4 & 77 & 3 & 3 & 76154 & 55, 56, 56, 58 & 4/4 \\
\end{longtable}
}


\subsection{General instances}
% Generated by paper/scripts/gen_appendix_tables.py from bobilib/results/; do not edit by hand.
% Needs booktabs, longtable, siunitx (v3) and multicol in the preamble.

{\footnotesize\sisetup{group-separator={,},group-minimum-digits=4,group-digits=integer}\setlength{\tabcolsep}{4pt}\setlength{\LTcapwidth}{\linewidth}
\begin{longtable}{@{}l l S[table-format=5.0] S[table-format=-9.4] l@{}}
\caption{The 40 open MIPLIB-derived instances with no bilevel feasible point: the follower's own rows admit a recession direction $d$ (follower variables only) with $f^\top d<0$, for every leader choice (Meyer 1974). Support: nonzero entries of $d$; $f^\top d$ rounded to four decimals; exact: checked in exact rational arithmetic by \texttt{scripts/verify\_ray.py} (the two supportcase42 directions leave residuals of about $10^{-12}$ relative and rest on floating point).}\label{tab:rays}\\
\toprule
Instance & Collection & {Support} & {$f^\top d$} & Exact \\
\midrule
\endfirsthead
\multicolumn{5}{l}{\emph{(continued)}}\\
\toprule
Instance & Collection & {Support} & {$f^\top d$} & Exact \\
\midrule
\endhead
\midrule
\multicolumn{5}{r}{\emph{continued on the next page}}\\
\endfoot
\bottomrule
\endlastfoot
biella1\_50\_50 & miplib2010 & 621 & -444178600 & yes \\
binkar10\_1\_50\_50 & miplib2010 & 208 & -90000000 & yes \\
brazil3\_0\_100 & miplib2017 & 3 & -1 & yes \\
brazil3\_50\_50 & miplib2017 & 11 & -1 & yes \\
buildingenergy\_0\_100 & miplib2017 & 17 & -10766.212 & yes \\
buildingenergy\_50\_50 & miplib2017 & 8777 & -11236.372 & yes \\
core2536-691\_50\_50 & miplib2010 & 9 & -200.9 & yes \\
cost266-UUE\_50\_50 & miplib2017 & 2024 & -13041.23 & yes \\
gen-ip054\_0\_100 & miplib2017 & 9 & -1475.2477 & yes \\
gen-ip054\_50\_50 & miplib2017 & 13 & -2718.065 & yes \\
highschool1-aigio\_0\_100 & miplib2017 & 3 & -1 & yes \\
irish-electricity\_0\_100 & miplib2017 & 5206 & -5206 & yes \\
irish-electricity\_50\_50 & miplib2017 & 5206 & -5206 & yes \\
lotsize\_50\_50 & miplib2017 & 297 & -871 & yes \\
mad\_0\_100 & miplib2017 & 20 & -20 & yes \\
mad\_50\_50 & miplib2017 & 20 & -20 & yes \\
neos-476283\_0\_100 & miplib2010 & 369 & -369 & yes \\
neos-476283\_50\_50 & miplib2010 & 369 & -369 & yes \\
neos-1396125\_0\_100 & miplib2010 & 3 & -240 & yes \\
neos-1396125\_50\_50 & miplib2010 & 3 & -240 & yes \\
neos-3627168-kasai\_50\_50 & miplib2017 & 3 & -3 & yes \\
neos-5093327-huahum\_0\_100 & miplib2017 & 3968 & -3968 & yes \\
neos-5093327-huahum\_50\_50 & miplib2017 & 3968 & -3968 & yes \\
nu25-pr12\_0\_100 & miplib2017 & 36 & -9000 & yes \\
rail507\_0\_100 & miplib2010 & 10 & -201 & yes \\
rail507\_50\_50 & miplib2010 & 10 & -201 & yes \\
rmatr100-p5\_0\_100 & miplib2010 & 4392 & -117792 & yes \\
rmatr100-p5\_50\_50 & miplib2010 & 4392 & -117792 & yes \\
rmatr100-p10\_0\_100 & miplib2010 & 3679 & -61179 & yes \\
rmatr100-p10\_50\_50 & miplib2010 & 3679 & -61179 & yes \\
rmatr200-p5\_0\_100 & miplib2017 & 18908 & -988708 & yes \\
rmatr200-p5\_50\_50 & miplib2017 & 18908 & -988708 & yes \\
snp-02-004-104\_0\_100 & miplib2017 & 104 & -10400000 & yes \\
snp-02-004-104\_50\_50 & miplib2017 & 104 & -10400000 & yes \\
supportcase40\_0\_100 & miplib2017 & 3360 & -338410.0838 & yes \\
supportcase40\_50\_50 & miplib2017 & 3360 & -338410.0838 & yes \\
supportcase42\_0\_100 & miplib2017 & 9733 & -8706 & no (floating point) \\
supportcase42\_50\_50 & miplib2017 & 9733 & -8706 & no (floating point) \\
tr12-30\_50\_50 & miplib2017 & 180 & -405 & yes \\
trento1\_50\_50 & miplib2017 & 655 & -33097000 & yes \\
\end{longtable}
}

% Generated by paper/scripts/gen_appendix_tables.py from bobilib/results/; do not edit by hand.
% Needs booktabs, longtable, siunitx (v3) and multicol in the preamble.

{\footnotesize\sisetup{group-separator={,},group-minimum-digits=4,group-digits=integer}\setlength{\tabcolsep}{4pt}\setlength{\LTcapwidth}{\linewidth}
\begin{longtable}{@{}l l l r r l@{}}
\caption{Five MIPLIB-derived instances whose high point relaxation is exactly the MIPLIB model they were built from (checked in exact arithmetic by \texttt{scripts/compare\_hpr.py}), so the published MIPLIB optimum (MIPLIB 2017's solution file) is a lower bound; a verified bilevel feasible point at that value is optimal by accepted benchmark conventions. The MIPLIB optima are floating-point results of MILP solvers (class (ii) of Section~\ref{app:classes}), not exact-arithmetic proofs: the objective of noswot takes integer values only, so its published bound rounds up to $-41$ and the result is safe if that bound is valid; pigeon-10 inherits MIPLIB's optimal status; for timtab1 the agreement of point and bound is numerical. Status: BOBILib's before this work. (neos-950242\_0\_100 meets its bound too, but was already listed as optimal in an earlier BOBILib release and is counted as published.)}\label{tab:hpr}\\
\toprule
Instance & BOBILib & MIPLIB instance & MIPLIB optimum & Bilevel optimum & Point from \\
\midrule
\endfirsthead
\multicolumn{6}{l}{\emph{(continued)}}\\
\toprule
Instance & BOBILib & MIPLIB instance & MIPLIB optimum & Bilevel optimum & Point from \\
\midrule
\endhead
\midrule
\multicolumn{6}{r}{\emph{continued on the next page}}\\
\endfoot
\bottomrule
\endlastfoot
noswot\_0\_100 & feasible & noswot & \num{-41.00000885} & \num{-41} & BOBILib's solution file \\
noswot\_50\_50 & feasible & noswot & \num{-41.00000885} & \num{-41} & BOBILib's solution file \\
pigeon-10\_0\_100 & open & pigeon-10 & \num{-9000} & \num{-9000} & \texttt{bobi general} \\
timtab1\_0\_100 & open & timtab1 & \num{764771.99999978} & \num{764771.99999978} & \texttt{bobi general} \\
pigeon-10\_50\_50 & open & pigeon-10 & \num{-9000} & \num{-9000} & \texttt{bobi coupled} \\
\end{longtable}
}

% Generated by paper/scripts/gen_appendix_tables.py from bobilib/results/; do not edit by hand.
% Needs booktabs, longtable, siunitx (v3) and multicol in the preamble.

{\footnotesize\sisetup{group-separator={,},group-minimum-digits=4,group-digits=integer}\setlength{\tabcolsep}{4pt}\setlength{\LTcapwidth}{\linewidth}
\begin{longtable}{@{}l l c l S[table-format=-10.4] S[table-format=-10.4]@{}}
\caption{The 88 verified bilevel feasible points for open MIPLIB-derived instances, for which BOBILib lists no point. $^{\ddagger}$: a point had been published elsewhere (its value and source are in the notes of \texttt{instances.csv}); for the other 79 this is the first point known. Coupling: whether leader rows contain follower variables. Method: \texttt{hpr-f-t}, the high point relaxation, then the follower solved with the leader fixed, then an optimistic tie-break (\texttt{bobi general}); \texttt{hpr}, \texttt{penalty}, \texttt{follow}: the part of the \texttt{bobi coupled} portfolio that found the point (the HPR's own leader values, a penalised HPR, best-response dynamics). Values rounded to four decimals (full precision in \texttt{instances.csv}); the leader minimises, so does the follower. Every file passes \texttt{bobi verify -{}-tight} (the follower's value within $\max(10^{-6},10^{-9}|\varphi|)$ of its proven optimum), so every point is solver-certified (class (ii) of Section~\ref{app:classes}). $^{\S}$: one of the 26 points re-derived under that strict tolerance by \texttt{bobi coupled-tighten}, at the leader values of a superseded file that had used BOBILib's tolerance of $10^{-4}$ on the follower's value (kept in \texttt{tolerance/}, Section~\ref{app:general}). Five further points are optimal (Table~\ref{tab:hpr}).}\label{tab:general-points}\\
\toprule
Instance & Collection & Coupling & Method & {Leader value} & {Follower value} \\
\midrule
\endfirsthead
\multicolumn{6}{l}{\emph{(continued)}}\\
\toprule
Instance & Collection & Coupling & Method & {Leader value} & {Follower value} \\
\midrule
\endhead
\midrule
\multicolumn{6}{r}{\emph{continued on the next page}}\\
\endfoot
\bottomrule
\endlastfoot
air05-0.900000$^{\ddagger}$ & miplib3 & no & hpr-f-t & 64955 & -63283 \\
beasleyC3\_0\_100 & miplib2010 & no & hpr-f-t & 6844 & -6844 \\
beasleyC3\_50\_50 & miplib2010 & yes & hpr & 6844 & -6844 \\
binkar10\_1\_0\_100$^{\S}$ & miplib2010 & no & hpr-f-t & 5006811.08 & -5000741.82 \\
cap6000-0.500000$^{\ddagger}$ & miplib3 & no & hpr-f-t & -1169217 & 0 \\
cap6000-0.900000$^{\ddagger}$ & miplib3 & no & hpr-f-t & -100780 & 45829 \\
cmflsp50-24-8-8\_0\_100$^{\S}$ & miplib2017 & no & hpr-f-t & 59480214.6542 & -12534898.0379 \\
core2536-691\_0\_100$^{\S}$ & miplib2010 & no & hpr-f-t & 10000000000 & -9999999544 \\
cost266-UUE\_0\_100$^{\S}$ & miplib2017 & no & hpr-f-t & 35801082 & -10431022.02 \\
exp-1-500-5-5\_0\_100$^{\S}$ & miplib2017 & no & hpr-f-t & 912511 & -879352 \\
h80x6320d\_0\_100 & miplib2017 & no & hpr-f-t & 12461.7042 & -10960.0714 \\
ic97\_potential\_0\_100$^{\S}$ & miplib2017 & no & hpr-f-t & 4003 & -353 \\
ic97\_potential\_50\_50$^{\S}$ & miplib2017 & yes & penalty & 4197 & -592 \\
lotsize\_0\_100 & miplib2017 & no & hpr-f-t & 3691725 & -3059545 \\
mc11\_0\_100 & miplib2017 & no & hpr-f-t & 128960 & -128960 \\
mc11\_50\_50 & miplib2017 & yes & hpr & 128960 & -128960 \\
mik-250-1-100-1\_0\_100 & miplib2010 & no & hpr-f-t & 890671 & -1000000 \\
mik-250-1-100-1\_50\_50 & miplib2010 & yes & hpr & 890671 & -1000000 \\
mik-250-20-75-4\_0\_100 & miplib2017 & no & hpr-f-t & 663799 & -750000 \\
mik-250-20-75-4\_50\_50 & miplib2017 & yes & hpr & 663799 & -750000 \\
mitre-0.500000$^{\ddagger}$ & miplib3 & no & hpr-f-t & 147075 & -96600 \\
mitre-0.900000$^{\ddagger}$ & miplib3 & no & hpr-f-t & 169500 & -169500 \\
mushroom-best\_0\_100 & miplib2017 & no & hpr-f-t & 4230.0553 & -4230 \\
mushroom-best\_50\_50 & miplib2017 & yes & hpr & 4230.0553 & -4230 \\
n4-3\_0\_100$^{\S}$ & miplib2010 & no & hpr-f-t & 34802896 & -34801026 \\
n4-3\_50\_50$^{\S}$ & miplib2010 & yes & penalty & 50621410 & -42710515 \\
n5-3\_0\_100$^{\S}$ & miplib2017 & no & hpr-f-t & 30001635 & -30000665 \\
n5-3\_50\_50$^{\S}$ & miplib2017 & yes & hpr & 30001775 & -30000805 \\
n9-3\_0\_100$^{\S}$ & miplib2017 & no & hpr-f-t & 50404633 & -50401629 \\
n9-3\_50\_50$^{\S}$ & miplib2017 & yes & penalty & 74423855 & -62411805 \\
neos17\_0\_100 & miplib2017 & no & hpr-f-t & 0.6872 & -0.564 \\
neos17\_50\_50$^{\S}$ & miplib2017 & yes & hpr & 0.6872 & -0.564 \\
neos-631710\_0\_100$^{\ddagger}$ & miplib2017 & no & hpr-f-t & 556 & -556 \\
neos-827175\_0\_100 & miplib2017 & no & hpr-f-t & 122.0015 & -122 \\
neos-827175\_50\_50 & miplib2017 & yes & hpr & 122.0015 & -122 \\
neos-860300\_0\_100 & miplib2017 & no & hpr-f-t & 5920 & -3702 \\
neos-873061\_0\_100$^{\S}$ & miplib2017 & no & hpr-f-t & 258644.4199 & -258644.1933 \\
neos-916792\_0\_100 & miplib2010 & no & hpr-f-t & 32.6951 & -5.6656 \\
neos-933966\_0\_100 & miplib2017 & no & hpr-f-t & 326 & -189 \\
neos-933966\_50\_50 & miplib2017 & yes & hpr & 47506 & -224 \\
neos-934278\_0\_100 & miplib2010 & no & hpr-f-t & 6339 & -267 \\
neos-934278\_50\_50 & miplib2010 & yes & hpr & 32490 & -299 \\
neos-950242\_50\_50 & miplib2017 & yes & follow & 24 & -24 \\
neos-1122047\_0\_100 & miplib2017 & no & hpr-f-t & 472 & -472 \\
neos-1337307\_50\_50 & miplib2010 & yes & hpr & -201272 & 27257 \\
neos-1445765\_50\_50 & miplib2017 & yes & penalty & 0 & 0 \\
neos-1456979\_0\_100 & miplib2017 & no & hpr-f-t & 250 & -157 \\
neos-2978193-inde\_0\_100$^{\S}$ & miplib2017 & no & hpr-f-t & 61.454 & -62.3292 \\
neos-2978193-inde\_50\_50$^{\S}$ & miplib2017 & yes & hpr & 69.2751 & -69.5863 \\
neos-3216931-puriri\_0\_100 & miplib2017 & no & hpr-f-t & 81430 & 900 \\
neos-3627168-kasai\_0\_100 & miplib2017 & no & hpr-f-t & 1591085.62 & -1591085.62 \\
neos-3754480-nidda\_0\_100 & miplib2017 & no & hpr-f-t & 13973.2905 & -2052644.3591 \\
neos-4300652-rahue\_0\_100 & miplib2017 & no & hpr-f-t & 6.4075 & -2.3316 \\
neos-4722843-widden\_0\_100 & miplib2017 & no & hpr-f-t & 25109.6634 & -7400 \\
neos-4763324-toguru\_0\_100 & miplib2017 & no & hpr-f-t & 17296.7271 & -13725.617 \\
neos-4954672-berkel\_0\_100 & miplib2017 & no & hpr-f-t & 8891491 & -7334836 \\
neos-4954672-berkel\_50\_50 & miplib2017 & yes & hpr & 8861094 & -7334836 \\
neos-5188808-nattai\_0\_100 & miplib2017 & no & hpr-f-t & 6.1216 & -6.1216 \\
neos-5188808-nattai\_50\_50 & miplib2017 & yes & hpr & 8.1556 & -8.1556 \\
netdiversion\_0\_100$^{\S}$ & miplib2010 & no & hpr-f-t & 19850250 & -19850250 \\
ns1830653\_0\_100 & miplib2010 & no & hpr-f-t & 23622 & -1910 \\
opm2-z10-s4\_0\_100 & miplib2017 & no & hpr-f-t & -11630 & 4812 \\
p200x1188c\_0\_100 & miplib2017 & no & hpr-f-t & 3584430 & -3584430 \\
p200x1188c\_50\_50 & miplib2017 & yes & hpr & 3584430 & -3584430 \\
p0548-0.500000$^{\ddagger}$ & miplib3 & no & hpr-f-t & 22197 & -17695 \\
pg5\_34\_0\_100$^{\S}$ & miplib2010 & no & hpr-f-t & -11516.2869 & 1824.44 \\
pg5\_34\_50\_50 & miplib2010 & yes & penalty & -11591.1699 & 0 \\
pg\_0\_100$^{\S}$ & miplib2017 & no & hpr-f-t & -5364.2592 & 2018.342 \\
pg\_50\_50$^{\S}$ & miplib2017 & yes & follow & -6061.0681 & 575.8461 \\
ran16x16\_0\_100 & miplib2010 & no & hpr-f-t & 43000 & -42126 \\
ran16x16\_50\_50 & miplib2010 & yes & hpr & 43000 & -42126 \\
rmine6\_0\_100$^{\ddagger}$ & miplib2010 & no & hpr-f-t & -440.6503 & 201.3995 \\
rmine6\_50\_50$^{\ddagger}$ & miplib2010 & yes & follow & -153.7498 & 119.897 \\
s100\_0\_100 & miplib2017 & no & hpr-f-t & -0.0758 & 0 \\
s250r10\_0\_100 & miplib2017 & no & hpr-f-t & -0.0608 & 0.0035 \\
satellites1-25\_0\_100 & miplib2010 & no & hpr-f-t & 32 & -32 \\
satellites2-40\_0\_100 & miplib2017 & no & hpr-f-t & 0 & 0 \\
seymour1\_0\_100$^{\S}$ & miplib2017 & no & hpr-f-t & 959.536 & -686 \\
seymour1\_50\_50$^{\S}$ & miplib2017 & yes & hpr & 960.6332 & -686 \\
sing44\_0\_100 & miplib2017 & no & hpr-f-t & 10059257.3791 & -112109.371 \\
sing326\_0\_100 & miplib2017 & no & hpr-f-t & 20652331.0689 & -181252.322 \\
sp150x300d\_0\_100 & miplib2017 & no & hpr-f-t & 300 & -300 \\
sp150x300d\_50\_50 & miplib2017 & yes & hpr & 300 & -300 \\
supportcase12\_0\_100 & miplib2017 & no & hpr-f-t & -1811.9515 & 0 \\
timtab1\_50\_50 & miplib2010 & yes & hpr & 764774 & -20655.996 \\
tr12-30\_0\_100$^{\S}$ & miplib2017 & no & hpr-f-t & 201019 & -92641 \\
uct-subprob\_0\_100$^{\S}$ & miplib2017 & no & hpr-f-t & 828 & -682 \\
uct-subprob\_50\_50$^{\S}$ & miplib2017 & yes & hpr & 848 & -682 \\
\end{longtable}
}


\subsection{Unchanged instances}
% Generated by paper/scripts/gen_appendix_tables.py from bobilib/results/; do not edit by hand.
% Needs booktabs, longtable, siunitx (v3) and multicol in the preamble.

{\footnotesize\sisetup{group-separator={,},group-minimum-digits=4,group-digits=integer}\setlength{\tabcolsep}{4pt}\setlength{\LTcapwidth}{\linewidth}
\begin{longtable}{@{}l p{0.34\linewidth} S[table-format=3.0] S[table-format=3.0] S[table-format=3.0] S[table-format=3.0]@{}}
\caption{Instances whose BOBILib status is unchanged, by family: BOBILib status (feasible or open), and how many now have a verified point better than BOBILib's (marked $^{*}$ in the lists below).}\label{tab:unchanged}\\
\toprule
Family & Collections & {Feasible} & {Open} & {Better point} & {Total} \\
\midrule
\endfirsthead
\multicolumn{6}{l}{\emph{(continued)}}\\
\toprule
Family & Collections & {Feasible} & {Open} & {Better point} & {Total} \\
\midrule
\endhead
\midrule
\multicolumn{6}{r}{\emph{continued on the next page}}\\
\endfoot
\bottomrule
\endlastfoot
multidimensional knapsack interdiction & or 582 & 503 & 79 & 462 & 582 \\
general, pure integer & miplib2017 42, miplib3 20, miplib2010 13, misc 1 & 60 & 16 & 0 & 76 \\
general, mixed integer & miplib2017 41, miplib2010 6 & 11 & 36 & 0 & 47 \\
generalized interdiction & generalized 8 & 8 & 0 & 0 & 8 \\
clique interdiction & clique 3 & 3 & 0 & 3 & 3 \\
\midrule
Total &  & 585 & 131 & 465 & 716 \\
\end{longtable}
}

\par\medskip\noindent\textbf{Multidimensional knapsack interdiction (582)}\par\nopagebreak\smallskip
{\setlength{\multicolsep}{3pt plus 1pt}\begin{multicols}{4}
\footnotesize\raggedright\setlength{\parindent}{0pt}\setlength{\parskip}{0pt}
OR5x100-0.25\_1.dat-0$^{*}$\par
OR5x100-0.25\_1.dat-50$^{*}$\par
OR5x100-0.25\_1.dat-100$^{*}$\par
OR5x100-0.25\_2.dat-0$^{*}$\par
OR5x100-0.25\_2.dat-50$^{*}$\par
OR5x100-0.25\_2.dat-100$^{*}$\par
OR5x100-0.25\_3.dat-0$^{*}$\par
OR5x100-0.25\_3.dat-50$^{*}$\par
OR5x100-0.25\_3.dat-100$^{*}$\par
OR5x100-0.25\_4.dat-0$^{*}$\par
OR5x100-0.25\_4.dat-50$^{*}$\par
OR5x100-0.25\_4.dat-100$^{*}$\par
OR5x100-0.25\_5.dat-0$^{*}$\par
OR5x100-0.25\_5.dat-50$^{*}$\par
OR5x100-0.25\_5.dat-100$^{*}$\par
OR5x100-0.25\_6.dat-0$^{*}$\par
OR5x100-0.25\_6.dat-50$^{*}$\par
OR5x100-0.25\_6.dat-100$^{*}$\par
OR5x100-0.25\_7.dat-0$^{*}$\par
OR5x100-0.25\_7.dat-50$^{*}$\par
OR5x100-0.25\_7.dat-100$^{*}$\par
OR5x100-0.25\_8.dat-0$^{*}$\par
OR5x100-0.25\_8.dat-50$^{*}$\par
OR5x100-0.25\_8.dat-100$^{*}$\par
OR5x100-0.25\_9.dat-0$^{*}$\par
OR5x100-0.25\_9.dat-50$^{*}$\par
OR5x100-0.25\_9.dat-100$^{*}$\par
OR5x100-0.25\_10.dat-0$^{*}$\par
OR5x100-0.25\_10.dat-50$^{*}$\par
OR5x100-0.25\_10.dat-100$^{*}$\par
OR5x100-0.50\_1.dat-0$^{*}$\par
OR5x100-0.50\_1.dat-50$^{*}$\par
OR5x100-0.50\_1.dat-100$^{*}$\par
OR5x100-0.50\_2.dat-0$^{*}$\par
OR5x100-0.50\_2.dat-50$^{*}$\par
OR5x100-0.50\_2.dat-100$^{*}$\par
OR5x100-0.50\_3.dat-0$^{*}$\par
OR5x100-0.50\_3.dat-50$^{*}$\par
OR5x100-0.50\_3.dat-100$^{*}$\par
OR5x100-0.50\_4.dat-0$^{*}$\par
OR5x100-0.50\_4.dat-50$^{*}$\par
OR5x100-0.50\_4.dat-100$^{*}$\par
OR5x100-0.50\_5.dat-0$^{*}$\par
OR5x100-0.50\_5.dat-50$^{*}$\par
OR5x100-0.50\_5.dat-100$^{*}$\par
OR5x100-0.50\_6.dat-0$^{*}$\par
OR5x100-0.50\_6.dat-50$^{*}$\par
OR5x100-0.50\_6.dat-100$^{*}$\par
OR5x100-0.50\_7.dat-0$^{*}$\par
OR5x100-0.50\_7.dat-50$^{*}$\par
OR5x100-0.50\_7.dat-100$^{*}$\par
OR5x100-0.50\_8.dat-0$^{*}$\par
OR5x100-0.50\_8.dat-50$^{*}$\par
OR5x100-0.50\_8.dat-100$^{*}$\par
OR5x100-0.50\_9.dat-0$^{*}$\par
OR5x100-0.50\_9.dat-50$^{*}$\par
OR5x100-0.50\_9.dat-100$^{*}$\par
OR5x100-0.50\_10.dat-0$^{*}$\par
OR5x100-0.50\_10.dat-50$^{*}$\par
OR5x100-0.50\_10.dat-100$^{*}$\par
OR5x250-0.25\_1.dat-0$^{*}$\par
OR5x250-0.25\_1.dat-50$^{*}$\par
OR5x250-0.25\_1.dat-100$^{*}$\par
OR5x250-0.25\_2.dat-0\par
OR5x250-0.25\_2.dat-50$^{*}$\par
OR5x250-0.25\_2.dat-100$^{*}$\par
OR5x250-0.25\_3.dat-0\par
OR5x250-0.25\_3.dat-50$^{*}$\par
OR5x250-0.25\_3.dat-100$^{*}$\par
OR5x250-0.25\_4.dat-0\par
OR5x250-0.25\_4.dat-50$^{*}$\par
OR5x250-0.25\_4.dat-100$^{*}$\par
OR5x250-0.25\_5.dat-0\par
OR5x250-0.25\_5.dat-50$^{*}$\par
OR5x250-0.25\_5.dat-100$^{*}$\par
OR5x250-0.25\_6.dat-0\par
OR5x250-0.25\_6.dat-50$^{*}$\par
OR5x250-0.25\_6.dat-100$^{*}$\par
OR5x250-0.25\_7.dat-0$^{*}$\par
OR5x250-0.25\_7.dat-50$^{*}$\par
OR5x250-0.25\_7.dat-100$^{*}$\par
OR5x250-0.25\_8.dat-0$^{*}$\par
OR5x250-0.25\_8.dat-50$^{*}$\par
OR5x250-0.25\_8.dat-100$^{*}$\par
OR5x250-0.25\_9.dat-0\par
OR5x250-0.25\_9.dat-50$^{*}$\par
OR5x250-0.25\_9.dat-100$^{*}$\par
OR5x250-0.25\_10.dat-0\par
OR5x250-0.25\_10.dat-50$^{*}$\par
OR5x250-0.25\_10.dat-100$^{*}$\par
OR5x250-0.50\_1.dat-0$^{*}$\par
OR5x250-0.50\_1.dat-50$^{*}$\par
OR5x250-0.50\_1.dat-100$^{*}$\par
OR5x250-0.50\_2.dat-0$^{*}$\par
OR5x250-0.50\_2.dat-50$^{*}$\par
OR5x250-0.50\_2.dat-100$^{*}$\par
OR5x250-0.50\_3.dat-0$^{*}$\par
OR5x250-0.50\_3.dat-50$^{*}$\par
OR5x250-0.50\_3.dat-100$^{*}$\par
OR5x250-0.50\_4.dat-0$^{*}$\par
OR5x250-0.50\_4.dat-50$^{*}$\par
OR5x250-0.50\_4.dat-100$^{*}$\par
OR5x250-0.50\_5.dat-0$^{*}$\par
OR5x250-0.50\_5.dat-50$^{*}$\par
OR5x250-0.50\_5.dat-100$^{*}$\par
OR5x250-0.50\_6.dat-0$^{*}$\par
OR5x250-0.50\_6.dat-50$^{*}$\par
OR5x250-0.50\_6.dat-100$^{*}$\par
OR5x250-0.50\_7.dat-0$^{*}$\par
OR5x250-0.50\_7.dat-50$^{*}$\par
OR5x250-0.50\_7.dat-100$^{*}$\par
OR5x250-0.50\_8.dat-0$^{*}$\par
OR5x250-0.50\_8.dat-50$^{*}$\par
OR5x250-0.50\_8.dat-100$^{*}$\par
OR5x250-0.50\_9.dat-0$^{*}$\par
OR5x250-0.50\_9.dat-50$^{*}$\par
OR5x250-0.50\_9.dat-100$^{*}$\par
OR5x250-0.50\_10.dat-0$^{*}$\par
OR5x250-0.50\_10.dat-50$^{*}$\par
OR5x250-0.50\_10.dat-100$^{*}$\par
OR5x500-0.25\_1.dat-0\par
OR5x500-0.25\_1.dat-50$^{*}$\par
OR5x500-0.25\_1.dat-100$^{*}$\par
OR5x500-0.25\_2.dat-0\par
OR5x500-0.25\_2.dat-50$^{*}$\par
OR5x500-0.25\_2.dat-100$^{*}$\par
OR5x500-0.25\_3.dat-0\par
OR5x500-0.25\_3.dat-50$^{*}$\par
OR5x500-0.25\_3.dat-100$^{*}$\par
OR5x500-0.25\_4.dat-0\par
OR5x500-0.25\_4.dat-50$^{*}$\par
OR5x500-0.25\_4.dat-100$^{*}$\par
OR5x500-0.25\_5.dat-0\par
OR5x500-0.25\_5.dat-50$^{*}$\par
OR5x500-0.25\_5.dat-100$^{*}$\par
OR5x500-0.25\_6.dat-0\par
OR5x500-0.25\_6.dat-50$^{*}$\par
OR5x500-0.25\_6.dat-100$^{*}$\par
OR5x500-0.25\_7.dat-0\par
OR5x500-0.25\_7.dat-50$^{*}$\par
OR5x500-0.25\_7.dat-100$^{*}$\par
OR5x500-0.25\_8.dat-0\par
OR5x500-0.25\_8.dat-50$^{*}$\par
OR5x500-0.25\_8.dat-100$^{*}$\par
OR5x500-0.25\_9.dat-0\par
OR5x500-0.25\_9.dat-50$^{*}$\par
OR5x500-0.25\_9.dat-100$^{*}$\par
OR5x500-0.25\_10.dat-0\par
OR5x500-0.25\_10.dat-50$^{*}$\par
OR5x500-0.25\_10.dat-100$^{*}$\par
OR5x500-0.50\_1.dat-0$^{*}$\par
OR5x500-0.50\_1.dat-50$^{*}$\par
OR5x500-0.50\_1.dat-100$^{*}$\par
OR5x500-0.50\_2.dat-50$^{*}$\par
OR5x500-0.50\_2.dat-100$^{*}$\par
OR5x500-0.50\_3.dat-0$^{*}$\par
OR5x500-0.50\_3.dat-50$^{*}$\par
OR5x500-0.50\_3.dat-100$^{*}$\par
OR5x500-0.50\_4.dat-50$^{*}$\par
OR5x500-0.50\_4.dat-100$^{*}$\par
OR5x500-0.50\_5.dat-0$^{*}$\par
OR5x500-0.50\_5.dat-50$^{*}$\par
OR5x500-0.50\_5.dat-100$^{*}$\par
OR5x500-0.50\_6.dat-0$^{*}$\par
OR5x500-0.50\_6.dat-50$^{*}$\par
OR5x500-0.50\_6.dat-100$^{*}$\par
OR5x500-0.50\_7.dat-0$^{*}$\par
OR5x500-0.50\_7.dat-50$^{*}$\par
OR5x500-0.50\_7.dat-100$^{*}$\par
OR5x500-0.50\_8.dat-0$^{*}$\par
OR5x500-0.50\_8.dat-50$^{*}$\par
OR5x500-0.50\_8.dat-100$^{*}$\par
OR5x500-0.50\_9.dat-50$^{*}$\par
OR5x500-0.50\_9.dat-100$^{*}$\par
OR5x500-0.50\_10.dat-0$^{*}$\par
OR5x500-0.50\_10.dat-50$^{*}$\par
OR5x500-0.50\_10.dat-100$^{*}$\par
OR5x500-0.75\_1.dat-0$^{*}$\par
OR5x500-0.75\_1.dat-50$^{*}$\par
OR5x500-0.75\_1.dat-100$^{*}$\par
OR5x500-0.75\_2.dat-0$^{*}$\par
OR5x500-0.75\_2.dat-50$^{*}$\par
OR5x500-0.75\_2.dat-100$^{*}$\par
OR5x500-0.75\_3.dat-0$^{*}$\par
OR5x500-0.75\_3.dat-50$^{*}$\par
OR5x500-0.75\_3.dat-100$^{*}$\par
OR5x500-0.75\_4.dat-0$^{*}$\par
OR5x500-0.75\_4.dat-50$^{*}$\par
OR5x500-0.75\_4.dat-100$^{*}$\par
OR5x500-0.75\_5.dat-0$^{*}$\par
OR5x500-0.75\_5.dat-50$^{*}$\par
OR5x500-0.75\_5.dat-100$^{*}$\par
OR5x500-0.75\_6.dat-0$^{*}$\par
OR5x500-0.75\_6.dat-50$^{*}$\par
OR5x500-0.75\_6.dat-100$^{*}$\par
OR5x500-0.75\_7.dat-0$^{*}$\par
OR5x500-0.75\_7.dat-50$^{*}$\par
OR5x500-0.75\_7.dat-100$^{*}$\par
OR5x500-0.75\_8.dat-0$^{*}$\par
OR5x500-0.75\_8.dat-50$^{*}$\par
OR5x500-0.75\_8.dat-100$^{*}$\par
OR5x500-0.75\_9.dat-0$^{*}$\par
OR5x500-0.75\_9.dat-50$^{*}$\par
OR5x500-0.75\_9.dat-100$^{*}$\par
OR5x500-0.75\_10.dat-0$^{*}$\par
OR5x500-0.75\_10.dat-50$^{*}$\par
OR5x500-0.75\_10.dat-100$^{*}$\par
OR10x100-0.25\_1.dat-0$^{*}$\par
OR10x100-0.25\_1.dat-50$^{*}$\par
OR10x100-0.25\_1.dat-100$^{*}$\par
OR10x100-0.25\_2.dat-0\par
OR10x100-0.25\_2.dat-50$^{*}$\par
OR10x100-0.25\_2.dat-100$^{*}$\par
OR10x100-0.25\_3.dat-0$^{*}$\par
OR10x100-0.25\_3.dat-50$^{*}$\par
OR10x100-0.25\_3.dat-100$^{*}$\par
OR10x100-0.25\_4.dat-0$^{*}$\par
OR10x100-0.25\_4.dat-50\par
OR10x100-0.25\_4.dat-100$^{*}$\par
OR10x100-0.25\_5.dat-0\par
OR10x100-0.25\_5.dat-50$^{*}$\par
OR10x100-0.25\_5.dat-100$^{*}$\par
OR10x100-0.25\_6.dat-0$^{*}$\par
OR10x100-0.25\_6.dat-50$^{*}$\par
OR10x100-0.25\_6.dat-100$^{*}$\par
OR10x100-0.25\_7.dat-0\par
OR10x100-0.25\_7.dat-50$^{*}$\par
OR10x100-0.25\_7.dat-100$^{*}$\par
OR10x100-0.25\_8.dat-0$^{*}$\par
OR10x100-0.25\_8.dat-50$^{*}$\par
OR10x100-0.25\_8.dat-100$^{*}$\par
OR10x100-0.25\_9.dat-0$^{*}$\par
OR10x100-0.25\_9.dat-50$^{*}$\par
OR10x100-0.25\_9.dat-100$^{*}$\par
OR10x100-0.25\_10.dat-0$^{*}$\par
OR10x100-0.25\_10.dat-50$^{*}$\par
OR10x100-0.25\_10.dat-100$^{*}$\par
OR10x100-0.50\_1.dat-0$^{*}$\par
OR10x100-0.50\_1.dat-50$^{*}$\par
OR10x100-0.50\_1.dat-100$^{*}$\par
OR10x100-0.50\_2.dat-0$^{*}$\par
OR10x100-0.50\_2.dat-50$^{*}$\par
OR10x100-0.50\_2.dat-100$^{*}$\par
OR10x100-0.50\_3.dat-0$^{*}$\par
OR10x100-0.50\_3.dat-50$^{*}$\par
OR10x100-0.50\_3.dat-100$^{*}$\par
OR10x100-0.50\_4.dat-0$^{*}$\par
OR10x100-0.50\_4.dat-50$^{*}$\par
OR10x100-0.50\_4.dat-100$^{*}$\par
OR10x100-0.50\_5.dat-0$^{*}$\par
OR10x100-0.50\_5.dat-50$^{*}$\par
OR10x100-0.50\_5.dat-100$^{*}$\par
OR10x100-0.50\_6.dat-0$^{*}$\par
OR10x100-0.50\_6.dat-50$^{*}$\par
OR10x100-0.50\_6.dat-100$^{*}$\par
OR10x100-0.50\_7.dat-0$^{*}$\par
OR10x100-0.50\_7.dat-50$^{*}$\par
OR10x100-0.50\_7.dat-100$^{*}$\par
OR10x100-0.50\_8.dat-0$^{*}$\par
OR10x100-0.50\_8.dat-50$^{*}$\par
OR10x100-0.50\_8.dat-100$^{*}$\par
OR10x100-0.50\_9.dat-0$^{*}$\par
OR10x100-0.50\_9.dat-50$^{*}$\par
OR10x100-0.50\_9.dat-100$^{*}$\par
OR10x100-0.50\_10.dat-0$^{*}$\par
OR10x100-0.50\_10.dat-50$^{*}$\par
OR10x100-0.50\_10.dat-100$^{*}$\par
OR10x100-0.75\_1.dat-0$^{*}$\par
OR10x100-0.75\_1.dat-50$^{*}$\par
OR10x100-0.75\_2.dat-50$^{*}$\par
OR10x100-0.75\_3.dat-0$^{*}$\par
OR10x100-0.75\_3.dat-50$^{*}$\par
OR10x100-0.75\_4.dat-50$^{*}$\par
OR10x100-0.75\_5.dat-0$^{*}$\par
OR10x100-0.75\_5.dat-50$^{*}$\par
OR10x100-0.75\_6.dat-50$^{*}$\par
OR10x100-0.75\_7.dat-50$^{*}$\par
OR10x100-0.75\_8.dat-50$^{*}$\par
OR10x100-0.75\_9.dat-50$^{*}$\par
OR10x100-0.75\_10.dat-0$^{*}$\par
OR10x100-0.75\_10.dat-50$^{*}$\par
OR10x250-0.25\_1.dat-0\par
OR10x250-0.25\_1.dat-50\par
OR10x250-0.25\_1.dat-100$^{*}$\par
OR10x250-0.25\_2.dat-0\par
OR10x250-0.25\_2.dat-50\par
OR10x250-0.25\_2.dat-100$^{*}$\par
OR10x250-0.25\_3.dat-0\par
OR10x250-0.25\_3.dat-50\par
OR10x250-0.25\_3.dat-100$^{*}$\par
OR10x250-0.25\_4.dat-0\par
OR10x250-0.25\_4.dat-50\par
OR10x250-0.25\_4.dat-100$^{*}$\par
OR10x250-0.25\_5.dat-0\par
OR10x250-0.25\_5.dat-50\par
OR10x250-0.25\_5.dat-100$^{*}$\par
OR10x250-0.25\_6.dat-0\par
OR10x250-0.25\_6.dat-50\par
OR10x250-0.25\_6.dat-100$^{*}$\par
OR10x250-0.25\_7.dat-0\par
OR10x250-0.25\_7.dat-50\par
OR10x250-0.25\_7.dat-100$^{*}$\par
OR10x250-0.25\_8.dat-0\par
OR10x250-0.25\_8.dat-50\par
OR10x250-0.25\_8.dat-100$^{*}$\par
OR10x250-0.25\_9.dat-0\par
OR10x250-0.25\_9.dat-50\par
OR10x250-0.25\_9.dat-100$^{*}$\par
OR10x250-0.25\_10.dat-0\par
OR10x250-0.25\_10.dat-50\par
OR10x250-0.25\_10.dat-100$^{*}$\par
OR10x250-0.50\_1.dat-50$^{*}$\par
OR10x250-0.50\_1.dat-100$^{*}$\par
OR10x250-0.50\_2.dat-50$^{*}$\par
OR10x250-0.50\_2.dat-100$^{*}$\par
OR10x250-0.50\_3.dat-50$^{*}$\par
OR10x250-0.50\_3.dat-100$^{*}$\par
OR10x250-0.50\_4.dat-50$^{*}$\par
OR10x250-0.50\_4.dat-100$^{*}$\par
OR10x250-0.50\_5.dat-50$^{*}$\par
OR10x250-0.50\_5.dat-100$^{*}$\par
OR10x250-0.50\_6.dat-50$^{*}$\par
OR10x250-0.50\_6.dat-100$^{*}$\par
OR10x250-0.50\_7.dat-50$^{*}$\par
OR10x250-0.50\_7.dat-100$^{*}$\par
OR10x250-0.50\_8.dat-50$^{*}$\par
OR10x250-0.50\_8.dat-100$^{*}$\par
OR10x250-0.50\_9.dat-50$^{*}$\par
OR10x250-0.50\_9.dat-100$^{*}$\par
OR10x250-0.50\_10.dat-50$^{*}$\par
OR10x250-0.50\_10.dat-100$^{*}$\par
OR10x250-0.75\_1.dat-50$^{*}$\par
OR10x250-0.75\_2.dat-50$^{*}$\par
OR10x250-0.75\_2.dat-100$^{*}$\par
OR10x250-0.75\_3.dat-50$^{*}$\par
OR10x250-0.75\_3.dat-100$^{*}$\par
OR10x250-0.75\_4.dat-50$^{*}$\par
OR10x250-0.75\_4.dat-100$^{*}$\par
OR10x250-0.75\_5.dat-50$^{*}$\par
OR10x250-0.75\_5.dat-100$^{*}$\par
OR10x250-0.75\_6.dat-50$^{*}$\par
OR10x250-0.75\_7.dat-50$^{*}$\par
OR10x250-0.75\_8.dat-50$^{*}$\par
OR10x250-0.75\_8.dat-100$^{*}$\par
OR10x250-0.75\_9.dat-0$^{*}$\par
OR10x250-0.75\_9.dat-50$^{*}$\par
OR10x250-0.75\_9.dat-100$^{*}$\par
OR10x250-0.75\_10.dat-50$^{*}$\par
OR10x250-0.75\_10.dat-100$^{*}$\par
OR10x500-0.25\_1.dat-0\par
OR10x500-0.25\_1.dat-50\par
OR10x500-0.25\_1.dat-100$^{*}$\par
OR10x500-0.25\_2.dat-0\par
OR10x500-0.25\_2.dat-50\par
OR10x500-0.25\_2.dat-100$^{*}$\par
OR10x500-0.25\_3.dat-0\par
OR10x500-0.25\_3.dat-50\par
OR10x500-0.25\_3.dat-100$^{*}$\par
OR10x500-0.25\_4.dat-0\par
OR10x500-0.25\_4.dat-50\par
OR10x500-0.25\_4.dat-100$^{*}$\par
OR10x500-0.25\_5.dat-0\par
OR10x500-0.25\_5.dat-50\par
OR10x500-0.25\_5.dat-100$^{*}$\par
OR10x500-0.25\_6.dat-0\par
OR10x500-0.25\_6.dat-50\par
OR10x500-0.25\_6.dat-100$^{*}$\par
OR10x500-0.25\_7.dat-0\par
OR10x500-0.25\_7.dat-50\par
OR10x500-0.25\_7.dat-100$^{*}$\par
OR10x500-0.25\_8.dat-0\par
OR10x500-0.25\_8.dat-50\par
OR10x500-0.25\_8.dat-100$^{*}$\par
OR10x500-0.25\_9.dat-0\par
OR10x500-0.25\_9.dat-50\par
OR10x500-0.25\_9.dat-100$^{*}$\par
OR10x500-0.25\_10.dat-0\par
OR10x500-0.25\_10.dat-50\par
OR10x500-0.25\_10.dat-100$^{*}$\par
OR10x500-0.50\_1.dat-100$^{*}$\par
OR10x500-0.50\_2.dat-100$^{*}$\par
OR10x500-0.50\_3.dat-100$^{*}$\par
OR10x500-0.50\_4.dat-100$^{*}$\par
OR10x500-0.50\_5.dat-100$^{*}$\par
OR10x500-0.50\_6.dat-100$^{*}$\par
OR10x500-0.50\_7.dat-100$^{*}$\par
OR10x500-0.50\_8.dat-100$^{*}$\par
OR10x500-0.50\_9.dat-100$^{*}$\par
OR10x500-0.50\_10.dat-100$^{*}$\par
OR10x500-0.75\_1.dat-100$^{*}$\par
OR10x500-0.75\_2.dat-100$^{*}$\par
OR10x500-0.75\_3.dat-50$^{*}$\par
OR10x500-0.75\_3.dat-100$^{*}$\par
OR10x500-0.75\_4.dat-100$^{*}$\par
OR10x500-0.75\_5.dat-100$^{*}$\par
OR10x500-0.75\_6.dat-100$^{*}$\par
OR10x500-0.75\_7.dat-50$^{*}$\par
OR10x500-0.75\_7.dat-100$^{*}$\par
OR10x500-0.75\_8.dat-50$^{*}$\par
OR10x500-0.75\_8.dat-100$^{*}$\par
OR10x500-0.75\_9.dat-100$^{*}$\par
OR10x500-0.75\_10.dat-50$^{*}$\par
OR10x500-0.75\_10.dat-100$^{*}$\par
OR30x100-0.25\_1.dat-0\par
OR30x100-0.25\_1.dat-50\par
OR30x100-0.25\_1.dat-100$^{*}$\par
OR30x100-0.25\_2.dat-0\par
OR30x100-0.25\_2.dat-50\par
OR30x100-0.25\_2.dat-100$^{*}$\par
OR30x100-0.25\_3.dat-0\par
OR30x100-0.25\_3.dat-50\par
OR30x100-0.25\_3.dat-100$^{*}$\par
OR30x100-0.25\_4.dat-0\par
OR30x100-0.25\_4.dat-50\par
OR30x100-0.25\_4.dat-100$^{*}$\par
OR30x100-0.25\_5.dat-0\par
OR30x100-0.25\_5.dat-50\par
OR30x100-0.25\_5.dat-100$^{*}$\par
OR30x100-0.25\_6.dat-0\par
OR30x100-0.25\_6.dat-50\par
OR30x100-0.25\_6.dat-100$^{*}$\par
OR30x100-0.25\_7.dat-0\par
OR30x100-0.25\_7.dat-50$^{*}$\par
OR30x100-0.25\_7.dat-100$^{*}$\par
OR30x100-0.25\_8.dat-0\par
OR30x100-0.25\_8.dat-50\par
OR30x100-0.25\_8.dat-100$^{*}$\par
OR30x100-0.25\_9.dat-0\par
OR30x100-0.25\_9.dat-50\par
OR30x100-0.25\_9.dat-100$^{*}$\par
OR30x100-0.25\_10.dat-0\par
OR30x100-0.25\_10.dat-50\par
OR30x100-0.25\_10.dat-100$^{*}$\par
OR30x100-0.50\_1.dat-0$^{*}$\par
OR30x100-0.50\_1.dat-50$^{*}$\par
OR30x100-0.50\_1.dat-100$^{*}$\par
OR30x100-0.50\_2.dat-0$^{*}$\par
OR30x100-0.50\_2.dat-50$^{*}$\par
OR30x100-0.50\_2.dat-100$^{*}$\par
OR30x100-0.50\_3.dat-50$^{*}$\par
OR30x100-0.50\_3.dat-100$^{*}$\par
OR30x100-0.50\_4.dat-0$^{*}$\par
OR30x100-0.50\_4.dat-50$^{*}$\par
OR30x100-0.50\_4.dat-100$^{*}$\par
OR30x100-0.50\_5.dat-50$^{*}$\par
OR30x100-0.50\_5.dat-100$^{*}$\par
OR30x100-0.50\_6.dat-0$^{*}$\par
OR30x100-0.50\_6.dat-50$^{*}$\par
OR30x100-0.50\_6.dat-100$^{*}$\par
OR30x100-0.50\_7.dat-0$^{*}$\par
OR30x100-0.50\_7.dat-50$^{*}$\par
OR30x100-0.50\_7.dat-100$^{*}$\par
OR30x100-0.50\_8.dat-0$^{*}$\par
OR30x100-0.50\_8.dat-50$^{*}$\par
OR30x100-0.50\_8.dat-100$^{*}$\par
OR30x100-0.50\_9.dat-0$^{*}$\par
OR30x100-0.50\_9.dat-50$^{*}$\par
OR30x100-0.50\_9.dat-100$^{*}$\par
OR30x100-0.50\_10.dat-0$^{*}$\par
OR30x100-0.50\_10.dat-50$^{*}$\par
OR30x100-0.50\_10.dat-100$^{*}$\par
OR30x100-0.75\_1.dat-0$^{*}$\par
OR30x100-0.75\_1.dat-50$^{*}$\par
OR30x100-0.75\_2.dat-0$^{*}$\par
OR30x100-0.75\_2.dat-50$^{*}$\par
OR30x100-0.75\_3.dat-0$^{*}$\par
OR30x100-0.75\_3.dat-50$^{*}$\par
OR30x100-0.75\_4.dat-0$^{*}$\par
OR30x100-0.75\_4.dat-50$^{*}$\par
OR30x100-0.75\_5.dat-50$^{*}$\par
OR30x100-0.75\_6.dat-0$^{*}$\par
OR30x100-0.75\_6.dat-50$^{*}$\par
OR30x100-0.75\_7.dat-0$^{*}$\par
OR30x100-0.75\_7.dat-50$^{*}$\par
OR30x100-0.75\_8.dat-0$^{*}$\par
OR30x100-0.75\_8.dat-50$^{*}$\par
OR30x100-0.75\_9.dat-0$^{*}$\par
OR30x100-0.75\_9.dat-50$^{*}$\par
OR30x100-0.75\_9.dat-100$^{*}$\par
OR30x100-0.75\_10.dat-0$^{*}$\par
OR30x100-0.75\_10.dat-50$^{*}$\par
OR30x100-0.75\_10.dat-100$^{*}$\par
OR30x250-0.25\_1.dat-0\par
OR30x250-0.25\_1.dat-50\par
OR30x250-0.25\_1.dat-100$^{*}$\par
OR30x250-0.25\_2.dat-0\par
OR30x250-0.25\_2.dat-50\par
OR30x250-0.25\_2.dat-100$^{*}$\par
OR30x250-0.25\_3.dat-0\par
OR30x250-0.25\_3.dat-50\par
OR30x250-0.25\_3.dat-100$^{*}$\par
OR30x250-0.25\_4.dat-0\par
OR30x250-0.25\_4.dat-50\par
OR30x250-0.25\_4.dat-100$^{*}$\par
OR30x250-0.25\_5.dat-0\par
OR30x250-0.25\_5.dat-50\par
OR30x250-0.25\_5.dat-100$^{*}$\par
OR30x250-0.25\_6.dat-0\par
OR30x250-0.25\_6.dat-50\par
OR30x250-0.25\_6.dat-100$^{*}$\par
OR30x250-0.25\_7.dat-0\par
OR30x250-0.25\_7.dat-50\par
OR30x250-0.25\_7.dat-100$^{*}$\par
OR30x250-0.25\_8.dat-0\par
OR30x250-0.25\_8.dat-50\par
OR30x250-0.25\_8.dat-100$^{*}$\par
OR30x250-0.25\_9.dat-0\par
OR30x250-0.25\_9.dat-50\par
OR30x250-0.25\_9.dat-100$^{*}$\par
OR30x250-0.25\_10.dat-0\par
OR30x250-0.25\_10.dat-50\par
OR30x250-0.25\_10.dat-100$^{*}$\par
OR30x250-0.50\_1.dat-100$^{*}$\par
OR30x250-0.50\_2.dat-100$^{*}$\par
OR30x250-0.50\_3.dat-100$^{*}$\par
OR30x250-0.50\_4.dat-100$^{*}$\par
OR30x250-0.50\_5.dat-100$^{*}$\par
OR30x250-0.50\_6.dat-100$^{*}$\par
OR30x250-0.50\_7.dat-100$^{*}$\par
OR30x250-0.50\_8.dat-100$^{*}$\par
OR30x250-0.50\_9.dat-100$^{*}$\par
OR30x250-0.50\_10.dat-100$^{*}$\par
OR30x250-0.75\_1.dat-100$^{*}$\par
OR30x250-0.75\_2.dat-100$^{*}$\par
OR30x250-0.75\_3.dat-100$^{*}$\par
OR30x250-0.75\_4.dat-100$^{*}$\par
OR30x250-0.75\_5.dat-100$^{*}$\par
OR30x250-0.75\_6.dat-100$^{*}$\par
OR30x250-0.75\_7.dat-100$^{*}$\par
OR30x250-0.75\_8.dat-100$^{*}$\par
OR30x250-0.75\_9.dat-100$^{*}$\par
OR30x250-0.75\_10.dat-100$^{*}$\par
OR30x500-0.25\_1.dat-0\par
OR30x500-0.25\_1.dat-50\par
OR30x500-0.25\_1.dat-100$^{*}$\par
OR30x500-0.25\_2.dat-0\par
OR30x500-0.25\_2.dat-50\par
OR30x500-0.25\_2.dat-100$^{*}$\par
OR30x500-0.25\_3.dat-0\par
OR30x500-0.25\_3.dat-50\par
OR30x500-0.25\_3.dat-100$^{*}$\par
OR30x500-0.25\_4.dat-0\par
OR30x500-0.25\_4.dat-50\par
OR30x500-0.25\_4.dat-100$^{*}$\par
OR30x500-0.25\_5.dat-0\par
OR30x500-0.25\_5.dat-50\par
OR30x500-0.25\_5.dat-100$^{*}$\par
OR30x500-0.25\_6.dat-0\par
OR30x500-0.25\_6.dat-50\par
OR30x500-0.25\_6.dat-100$^{*}$\par
OR30x500-0.25\_7.dat-0\par
OR30x500-0.25\_7.dat-50\par
OR30x500-0.25\_7.dat-100$^{*}$\par
OR30x500-0.25\_8.dat-0\par
OR30x500-0.25\_8.dat-50\par
OR30x500-0.25\_8.dat-100$^{*}$\par
OR30x500-0.25\_9.dat-0\par
OR30x500-0.25\_9.dat-50\par
OR30x500-0.25\_9.dat-100$^{*}$\par
OR30x500-0.25\_10.dat-0\par
OR30x500-0.25\_10.dat-50\par
OR30x500-0.25\_10.dat-100$^{*}$\par
OR30x500-0.50\_1.dat-100$^{*}$\par
OR30x500-0.50\_2.dat-100$^{*}$\par
OR30x500-0.50\_3.dat-100$^{*}$\par
OR30x500-0.50\_4.dat-100$^{*}$\par
OR30x500-0.50\_5.dat-100$^{*}$\par
OR30x500-0.50\_6.dat-100$^{*}$\par
OR30x500-0.50\_7.dat-100$^{*}$\par
OR30x500-0.50\_8.dat-100$^{*}$\par
OR30x500-0.50\_9.dat-100$^{*}$\par
OR30x500-0.50\_10.dat-100$^{*}$\par
OR30x500-0.75\_1.dat-100$^{*}$\par
OR30x500-0.75\_2.dat-100$^{*}$\par
OR30x500-0.75\_3.dat-100$^{*}$\par
OR30x500-0.75\_4.dat-100$^{*}$\par
OR30x500-0.75\_5.dat-100$^{*}$\par
OR30x500-0.75\_6.dat-100$^{*}$\par
OR30x500-0.75\_7.dat-100$^{*}$\par
OR30x500-0.75\_8.dat-100$^{*}$\par
OR30x500-0.75\_9.dat-100$^{*}$\par
OR30x500-0.75\_10.dat-100$^{*}$\par
\end{multicols}}

\par\medskip\noindent\textbf{General, pure integer (76)}\par\nopagebreak\smallskip
{\setlength{\multicolsep}{3pt plus 1pt}\begin{multicols}{3}
\footnotesize\raggedright\setlength{\parindent}{0pt}\setlength{\parskip}{0pt}
air03-0.500000\par
air03-0.900000\par
air04-0.100000\par
air04-0.500000\par
air04-0.900000\par
air05-0.500000\par
bab2\_0\_100\par
bab5\_0\_100\par
bab6\_0\_100\par
co-100\_0\_100\par
co-100\_50\_50\par
comp07-2idx\_0\_100\par
comp21-2idx\_0\_100\par
cvs16r128-89\_0\_100\par
cvs16r128-89\_50\_50\par
eilA101-2\_0\_100\par
ex10\_50\_50\par
gen-ip002\_0\_100\par
gen-ip002\_50\_50\par
harp2-0.100000\par
harp2-0.500000\par
harp2-0.900000\par
irp\_0\_100\par
l152lav-0.500000\par
l152lav-0.900000\par
lectsched-5-obj\_50\_50\par
milp\_10\_20\_50\_2310\par
mine-90-10\_0\_100\par
mine-166-5\_0\_100\par
mine-166-5\_50\_50\par
mitre-0.100000\par
mod010-0.500000\par
mod010-0.900000\par
neos-631710\_50\_50\par
neos-662469\_0\_100\par
neos-957323\_0\_100\par
neos-1109824\_0\_100\par
neos-1337307\_0\_100\par
neos-1582420\_0\_100\par
neos-1582420\_50\_50\par
neos-2657525-crna\_0\_100\par
neos-3024952-loue\_0\_100\par
neos-3083819-nubu\_0\_100\par
neos-5052403-cygnet\_0\_100\par
neos-5052403-cygnet\_50\_50\par
nursesched-medium-hint03\_0\_100\par
nw04-0.500000\par
nw04-0.900000\par
opm2-z7-s2\_0\_100\par
p2756-0.100000\par
p2756-0.500000\par
p2756-0.900000\par
reblock67\_0\_100\par
reblock115\_0\_100\par
s100\_50\_50\par
seymour-0.100000\par
sorrell3\_0\_100\par
sorrell3\_50\_50\par
sp97ar\_0\_100\par
sp97ar\_50\_50\par
sp98ar\_0\_100\par
sp98ar\_50\_50\par
sp98ir\_0\_100\par
square41\_0\_100\par
square41\_50\_50\par
square47\_0\_100\par
square47\_50\_50\par
supportcase6\_0\_100\par
supportcase6\_50\_50\par
supportcase19\_0\_100\par
tanglegram1\_0\_100\par
tanglegram1\_50\_50\par
tanglegram2\_0\_100\par
tanglegram2\_50\_50\par
toll-like\_0\_100\par
toll-like\_50\_50\par
\end{multicols}}

\par\medskip\noindent\textbf{General, mixed integer (47)}\par\nopagebreak\smallskip
{\setlength{\multicolsep}{3pt plus 1pt}\begin{multicols}{3}
\footnotesize\raggedright\setlength{\parindent}{0pt}\setlength{\parskip}{0pt}
atlanta-ip\_0\_100\par
biella1\_0\_100\par
cmflsp50-24-8-8\_50\_50\par
drayage-25-23\_0\_100\par
drayage-100-23\_0\_100\par
exp-1-500-5-5\_50\_50\par
fiball\_0\_100\par
gmu-35-40\_0\_100\par
gmu-35-40\_50\_50\par
gmu-35-50\_0\_100\par
gmu-35-50\_50\_50\par
h80x6320d\_50\_50\par
momentum1\_0\_100\par
momentum1\_50\_50\par
neos-848589\_0\_100\par
neos-848589\_50\_50\par
neos-873061\_50\_50\par
neos-916792\_50\_50\par
neos-1122047\_50\_50\par
neos-1456979\_50\_50\par
neos-2746589-doon\_0\_100\par
neos-3216931-puriri\_50\_50\par
neos-3754224-navua\_0\_100\par
neos-3754224-navua\_50\_50\par
neos-3754480-nidda\_50\_50\par
neos-4300652-rahue\_50\_50\par
neos-4387871-tavua\_0\_100\par
neos-4387871-tavua\_50\_50\par
neos-4722843-widden\_50\_50\par
neos-4763324-toguru\_50\_50\par
neos-5049753-cuanza\_0\_100\par
neos-5049753-cuanza\_50\_50\par
neos-5104907-jarama\_0\_100\par
neos-5104907-jarama\_50\_50\par
neos-5114902-kasavu\_0\_100\par
neos-5114902-kasavu\_50\_50\par
nexp-150-20-8-5\_0\_100\par
ns1830653\_50\_50\par
ran14x18-disj-8\_0\_100\par
satellites1-25\_50\_50\par
satellites2-40\_50\_50\par
satellites2-60-fs\_0\_100\par
satellites2-60-fs\_50\_50\par
sing44\_50\_50\par
sing326\_50\_50\par
supportcase12\_50\_50\par
trento1\_0\_100\par
\end{multicols}}

\par\medskip\noindent\textbf{Generalized interdiction (8)}\par\nopagebreak\smallskip
{\setlength{\multicolsep}{3pt plus 1pt}\begin{multicols}{4}
\footnotesize\raggedright\setlength{\parindent}{0pt}\setlength{\parskip}{0pt}
general20-30-20-20-20-2\par
general20-30-20-20-20-9\par
general20-30-20-20-20-10\par
general30-30-20-20-20-1\par
general30-30-20-20-20-2\par
general30-30-30-20-20-1\par
general30-30-30-20-20-2\par
general30-30-30-20-20-6\par
\end{multicols}}

\par\medskip\noindent\textbf{Clique interdiction (3)}\par\nopagebreak\smallskip
{\setlength{\multicolsep}{3pt plus 1pt}\begin{multicols}{4}
\footnotesize\raggedright\setlength{\parindent}{0pt}\setlength{\parskip}{0pt}
normalClique60-70-2$^{*}$\par
normalClique60-70-3$^{*}$\par
normalClique60-70-5$^{*}$\par
\end{multicols}}



\bibliographystyle{unsrt}
\bibliography{../refs}

\end{document}
